\documentclass[11pt,letterpaper]{book}
\usepackage[margin=1in]{geometry}
\usepackage[T1]{fontenc}
\usepackage[utf8]{inputenc}
\usepackage{lmodern}
\usepackage{microtype}
\usepackage{xcolor}
\usepackage{listings}
\usepackage{booktabs}
\usepackage{array}
\usepackage{emptypage}
\usepackage[hidelinks]{hyperref}

\definecolor{kw}{RGB}{20,60,130}
\definecolor{cm}{RGB}{80,120,80}
\definecolor{st}{RGB}{150,50,50}
\definecolor{rulecol}{RGB}{170,175,190}
\definecolor{numcol}{RGB}{150,150,150}

\lstdefinestyle{code}{language=Python,basicstyle=\ttfamily\small,
  keywordstyle=\color{kw}\bfseries,commentstyle=\color{cm}\itshape,
  stringstyle=\color{st},showstringspaces=false,keepspaces=true,
  columns=fullflexible,breaklines=true,breakatwhitespace=true,
  postbreak=\mbox{\textcolor{numcol}{$\hookrightarrow$}\space},
  upquote=true,aboveskip=8pt,belowskip=8pt,xleftmargin=1.1em,
  frame=leftline,framerule=0.8pt,rulecolor=\color{rulecol}}
\lstdefinestyle{spec}{basicstyle=\ttfamily\small,keepspaces=true,
  columns=fullflexible,breaklines=true,breakatwhitespace=true,
  postbreak=\mbox{\textcolor{numcol}{$\hookrightarrow$}\space},
  upquote=true,aboveskip=6pt,belowskip=6pt,xleftmargin=1.1em,
  frame=leftline,framerule=0.8pt,rulecolor=\color{rulecol},language={}}
\lstdefinestyle{file}{language=Python,basicstyle=\ttfamily\footnotesize,
  keywordstyle=\color{kw}\bfseries,commentstyle=\color{cm}\itshape,
  stringstyle=\color{st},showstringspaces=false,keepspaces=true,
  columns=fullflexible,breaklines=true,breakatwhitespace=true,
  postbreak=\mbox{\textcolor{numcol}{$\hookrightarrow$}\space},
  upquote=true,numbers=left,numberstyle=\tiny\color{numcol},numbersep=8pt,
  aboveskip=6pt,belowskip=6pt}

\setcounter{tocdepth}{1}
\setlength{\parskip}{2pt}

\title{\Huge\bfseries PythonFP\\[6pt]
  \Large Functional Programming in Python, the Prelude Way\\[16pt]
  \normalsize A course in one file and seventy-four problems}
\author{J.~L.~Clements~III}
\date{\today}

\begin{document}
\frontmatter
\maketitle

\chapter{Preface}
This book teaches functional programming in Python through a single artifact:
\texttt{prelude.py}, a Haskell-style prelude rebuilt in pure Python with no
imports, read top to bottom in strict define-before-use order. Around it sit
seventy-four problems, each solved by \emph{composing} prelude tools rather
than narrating loops.

The course has three movements. Part~I teaches how to \emph{think} in this
style: what functional Python buys you, and the recognition drill that maps the
words of a problem statement to the right tool. Part~II walks the prelude
section by section -- every definition, why it exists, and what it displaces.
Part~III applies the toolkit to nine classes of problems; every problem appears
with its statement, contract, hint, doctests and a worked solution.

The book is generated from the course files themselves, so every listing is the
real code. The intended use is active: keep \texttt{prelude.py} open beside the
book, attempt each problem in its practice file before reading the solution,
and let the doctests judge you. Appendix~D describes that workflow.

\tableofcontents
\mainmatter
\part{Thinking in PythonFP}
\chapter{Why Functional Python}
\section{One file, no imports}
Everything in this course runs on bare \texttt{python3}. That constraint is not
asceticism; it is the point. When \texttt{collections.Counter} is an import,
it is a black box; when it is six lines you have retyped from memory, it is
knowledge. The prelude rebuilds the working parts of Haskell's standard
vocabulary -- and the useful corners of \texttt{Data.List}, \texttt{Data.Map},
\texttt{Data.Maybe} and friends -- as Python one-liners and small definitions
you can carry into any interview editor, any locked-down environment, any
whiteboard.

\section{What ``functional'' buys you}
Three habits, none of which require giving up Python:

\begin{itemize}
\item \textbf{Functions are values.} Sort keys, predicates, combining
functions, closures over window state: passing behaviour as data is what lets
one skeleton (\texttt{longest\_window}, \texttt{foldl},
\texttt{fromListWith}) serve a hundred problems.
\item \textbf{Data in, data out.} Prefer expressions that build new values to
statements that mutate old ones. Not because mutation is sinful -- the prelude
mutates inside its machines -- but because data-in/data-out functions compose,
test in one line, and survive a mutating interview prompt.
\item \textbf{Name the transformation.} A loop narrates \emph{how}; a
composition of named tools states \emph{what}. \texttt{unwords(reversed(%
words(s)))} is its own documentation. Vocabulary is leverage: every name you
own is a loop you never write again.
\end{itemize}

\section{Conventions you must internalise}
\textbf{None is Nothing.} Every partial prelude function signals absence with
\texttt{None}; \texttt{fromMaybe} lands defaults, \texttt{mapMaybe} and
\texttt{catMaybes} sweep Nothings out of lists, \texttt{isJust} and
\texttt{isNothing} test a single value. Never overload \texttt{0} or
\texttt{""} to mean ``missing''.

\textbf{Haskell names shadow builtins by design.} \texttt{id}, \texttt{map\_},
\texttt{filter\_}, \texttt{sort}: inside prelude code the Haskell vocabulary
wins. The trailing underscore marks the few collisions with keywords
(\texttt{not\_} being the survivor; \texttt{all}/\texttt{any} serve bare).

\textbf{Define before use.} \texttt{prelude.py} reads top to bottom; every name
is already defined when you meet it. This is also how you should
\emph{study} it: read a section, close the file, retype its one-liners from
memory, and check yourself. The definitions are short on purpose.
\chapter{The Recognition Drill}
Most timed failures are not coding failures; they are \emph{recognition}
failures -- ten minutes spent hand-rolling a loop that was a \texttt{scanl1}.
The drill below runs before any typing.

\section{Step 0: name three things}
From the statement, extract: the \textbf{input shape} (a list? two lists?
pairs? a grid? edges?); the \textbf{output shape} (one value? a list? a dict?
yes/no? \emph{all} solutions?); and the \textbf{quantity word} (the ``most'',
``fewest'', ``longest'', ``every'' that names what is asked). The output shape
is the strongest clue: one value smells like a fold, a dict smells like
\texttt{fromListWith}, ``all of them'' smells like enumeration or a tree.

\section{Step 1: walk the ladder}
Match the statement's words against the trigger column; the first row that
fits names your tools.

\begin{center}\small
\begin{tabular}{@{}p{0.14\textwidth}p{0.44\textwidth}p{0.32\textwidth}@{}}
\toprule
\textbf{Shape} & \textbf{Trigger words} & \textbf{Reach for} \\
\midrule
SORT FIRST & ``sorted'', ``k-th'', ``closest pair'', ``can they be arranged'', ``intervals'' & \texttt{sort, sortOn, bisect\_left, merge} \\
FOLD & ``total'', ``count'', ``largest'', ``best single \ldots'' & \texttt{foldl, minOn / maxOn} \\
SCAN & ``running'', ``so far'', ``at each step'', ``prefix sums'' & \texttt{scanl, scanl1} \\
DICT FOLD & ``group by'', ``frequency'', ``most common'', ``anagrams together'' & \texttt{fromListWith, Counter, unionWith} \\
WINDOW & ``substring'', ``contiguous'', ``longest run'', ``within the last k'' & \texttt{longest\_window, pairwise} \\
STACK & ``nested'', ``matching pairs'', ``most recent open'', ``next greater'' & \texttt{list as stack; foldl with stack acc} \\
UNFOLD & ``digits of'', ``split into pieces'', ``repeat until'', ``simulate'' & \texttt{unfoldr, iterate + take} \\
ENUMERATE & ``all subsets'', ``any combination'', ``every way to pick'' & \texttt{subsequences, replicateM, cross} \\
TREE / PATHS & ``nested'', ``prefix'' (trie), ``directories'', ``solutions as paths'' & \texttt{ITree, paths, setpath, tfold} \\
GRAPH & ``prerequisites'', ``network'', ``shortest route'', ``reachable'' & \texttt{Tree(1,list) adjacency, deque BFS, heap, until} \\
MEMOIZE & ``fewest'', ``max value'', ``count the ways'', ``can it be done'' & \texttt{memo over indices} \\
\bottomrule
\end{tabular}
\end{center}

Cross-cutting, any row: messy input $\rightarrow$ \texttt{mapMaybe}; absence
$\rightarrow$ \texttt{None} + \texttt{fromMaybe}; first hit $\rightarrow$
\texttt{find}; same-up-to-order $\rightarrow$ a canonical form (\texttt{sort}
or \texttt{Counter}); connectivity $\rightarrow$ \texttt{dsu}; repeated
best-next $\rightarrow$ the heap.

\section{Step 2: say it out loud}
State the shape, the tools, and the complexity before typing: ``group-by-key
with \texttt{fromListWith}, then best-by-key with \texttt{minOn},
$O(n)$.'' If two rows both apply, \textbf{the row that names the output wins}:
``most common word'' \emph{mentions} frequency (dict fold) but \emph{asks for}
one value (fold) -- so \texttt{Counter} feeds \texttt{minOn}.

Complexity vocabulary to keep loaded: $\sim$$10^7$ simple Python operations
per second is a safe planning number; $n \le 20$ licenses $2^n$ enumeration;
$n \le 10^4$ licenses $O(n^2)$; $n \le 10^6$ demands $O(n \log n)$ or better.

\section{Step 3: compose, don't invent}
Solutions in this course are pipelines of two or three prelude names plus one
lambda of your own. If you find yourself writing a raw loop, first ask which
combinator you are rebuilding. Sometimes the answer is ``none -- this loop is
the algorithm'' (BFS, a parser, a heap sift); then write the loop proudly.
\part{The Prelude, Section by Section}
\chapter{Functions About Functions}

Functional programming begins with one liberty: a function is a value. It can
be passed, returned, stored in a dict, and -- most usefully -- \emph{shaped} by
other functions. Section~1 is a kit of shapers. \texttt{flip} swaps an argument
order so a function fits a slot it was not written for; \texttt{curry} and
\texttt{uncurry} convert between two-argument functions and functions of pairs;
\texttt{partial} freezes the first arguments; \texttt{on} routes both arguments
of a comparison through a key first.

Two definitions here are load-bearing conventions rather than tools.
\texttt{NOTHING} is a private sentinel meaning ``no argument was supplied'' --
it lets \texttt{foldl} and \texttt{accumulate} distinguish an omitted initial
value from a legitimate \texttt{None}. And the file-wide contract \emph{None is
Nothing}: every partial function in the prelude signals absence with
\texttt{None}, \texttt{fromMaybe} is the standard way to land a default, and
\texttt{isJust}/\texttt{isNothing} name the None test itself -- a predicate
you can hand to \texttt{find}, \texttt{filter\_} or \texttt{all}, and one that
asks identity against \texttt{None} rather than truthiness, so \texttt{0} and
\texttt{""} count as present.
Section~2 adds the pair accessors -- \texttt{fst}, \texttt{snd}, \texttt{swap}
-- which read better than \texttt{[0]} and \texttt{[1]} in code built from
tuples.
\begin{lstlisting}[style=code]
# ============ 1. combinators ============ (functions about functions)

id        = lambda x: x                             # shadows builtin id() by design
const     = lambda x: lambda _: x
flip      = lambda f: (lambda x, y: f(y, x))        # flip f x y = f y x
curry     = lambda f: lambda x: lambda y: f(x, y)
uncurry   = lambda f: lambda p: f(*p)
partial   = lambda f, *bound: lambda *args: f(*bound, *args)
on        = lambda f, g: lambda x, y: f(g(x), g(y))  # Data.Function: combine via a key -- cmp `on` fst
NOTHING   = object()                                 # Maybe's Nothing: "no arg given"; test with `is`
fromMaybe = lambda d, x: d if x is None else x       # Data.Maybe: default for every None-returning tool
isJust    = lambda x: x is not None                  # Data.Maybe: None test as a handable predicate
isNothing = lambda x: x is None
\end{lstlisting}
\begin{lstlisting}[style=code]
# ============ 2. pairs ============

fst  = lambda p: p[0]
snd  = lambda p: p[1]
swap = lambda p: (p[1], p[0])               # Data.Tuple
\end{lstlisting}

\noindent\textit{At the REPL:}
\begin{lstlisting}[style=spec]
>>> add = lambda a, b: a + b
>>> flip(add)("world", "hello ")
'hello world'
>>> on(add, len)("ab", "c")            # combine THROUGH a key
3
>>> fromMaybe(0, None), fromMaybe(0, 5)
(0, 5)
>>> isJust(0), isNothing(None)         # identity, never truthiness
(True, True)
>>> swap((1, 2))
(2, 1)
\end{lstlisting}
\chapter{Arithmetic and Logic}

Small tools, chosen for their sharp edges. \texttt{signum} compresses a
three-way comparison into arithmetic. \texttt{div}/\texttt{mod} floor like
Haskell's \texttt{div}; \texttt{quot}/\texttt{rem} truncate toward zero --
knowing which one your language does to negative numbers has decided real
interviews. \texttt{gcd} recurses safely because Euclid's depth is tiny even
for 64-bit inputs (the worst case is consecutive Fibonaccis), and
\texttt{isqrt} is Newton's method on integers: floor square roots with no
floating point in sight.
\begin{lstlisting}[style=code]
# ============ 3. arithmetic & logic ============

succ      = lambda x: chr(ord(x) + 1) if isinstance(x, str) else x + 1   # Enum: numbers and Chars
pred      = lambda x: chr(ord(x) - 1) if isinstance(x, str) else x - 1
even      = lambda n: n % 2 == 0
odd       = lambda n: n % 2 == 1
not_      = lambda b: not b                                              # `not` is a Python keyword
otherwise = True
signum    = lambda x: (x > 0) - (x < 0)
div       = lambda a, b: a // b                                          # floors, like Haskell div
mod       = lambda a, b: a % b
quot      = lambda a, b: -(-a // b) if (a < 0) != (b < 0) else a // b    # truncates, like Haskell quot
rem       = lambda a, b: a - b * quot(a, b)
quotRem   = lambda a, b: (quot(a, b), rem(a, b))
gcd       = lambda a, b: abs(a) if b == 0 else gcd(b, a % b)             # >= 0 like Haskell; depth <= 90
lcm       = lambda a, b: abs(a // gcd(a, b) * b) if a and b else 0       # >= 0 like Haskell
hypot     = lambda x, y: (x*x + y*y) ** 0.5

def isqrt(n):  # floor sqrt without floats (Newton)
    if n < 0:
        raise ValueError("isqrt of negative")
    x = n
    y = (x + 1) // 2
    while y < x:
        x, y = y, (y + n // y) // 2
    return x if n else 0
\end{lstlisting}

\noindent\textit{At the REPL:}
\begin{lstlisting}[style=spec]
>>> signum(-7), signum(0), signum(3)
(-1, 0, 1)
>>> div(-7, 2), quot(-7, 2)            # floor vs truncate
(-4, -3)
>>> isqrt(10**18)
1000000000
\end{lstlisting}
\chapter{Lists I: The Basics}

The list vocabulary. \texttt{head}/\texttt{tail}/\texttt{init}/\texttt{last}
name the four ways to take a list apart at its ends; \texttt{null} asks if
there is anything left -- together they are the grammar of structural
recursion. Three definitions deserve special study. \texttt{nub} deduplicates
\emph{keeping first occurrences}, exploiting the guarantee that dicts remember
insertion order. \texttt{transpose} flips rows and columns, ragged-safe like
Haskell's -- short rows simply drop out of later columns. \texttt{pairwise}
zips a list against its own tail, giving every element
its right neighbour -- the tool for any rule about adjacency. The prefix family
(\texttt{isPrefixOf}, \texttt{isSuffixOf}, \texttt{stripPrefix}) works on
strings and lists alike, and \texttt{isSuffixOf}'s comment preserves a scar:
\texttt{xs[-0:]} is the \emph{whole} list, so never slice a suffix by
\texttt{-len(p)}.
\begin{lstlisting}[style=code]
# ============ 4. list basics ============

head        = lambda xs: xs[0]
tail        = lambda xs: xs[1:]
init        = lambda xs: xs[:-1]
last        = lambda xs: xs[-1]
null        = lambda xs: len(xs) == 0
elem        = lambda x, xs: x in xs
notElem     = lambda x, xs: x not in xs
replicate   = lambda n, x: [x] * n
drop        = lambda n, xs: list(xs)[max(n, 0):]                # finite lists only; n<0 drops nothing
splitAt     = lambda n, xs: (xs[:n], xs[n:]) if n >= 0 else (xs[:0], xs)   # n<0 splits at the front
nub         = lambda xs: list(dict.fromkeys(xs))    # unique, first occurrence wins (dicts keep order)
lookup      = lambda k, pairs: next((v for kk, v in pairs if kk == k), None)   # Nothing -> None
zip_        = lambda a, b: list(zip(a, b))          # shortest wins, any iterable
zip3        = lambda a, b, c: list(zip(a, b, c))
unzip       = lambda ps: tuple(map(list, zip(*ps))) if ps else ([], [])
enum        = lambda xs, start=0: zip_(list(range(start, start + len(xs))), xs)
pairwise    = lambda xs: list(zip(xs, xs[1:]))              # zip xs (tail xs)
isPrefixOf  = lambda p, xs: list(xs[:len(p)]) == list(p)             # Data.List; strings and lists alike
isSuffixOf  = lambda p, xs: list(xs[len(xs) - len(p):]) == list(p)   # NB not [-len(p):] -- [-0:] is ALL
stripPrefix = lambda p, xs: xs[len(p):] if isPrefixOf(p, xs) else None  # Just the rest, or Nothing
def transpose(xss):                    # Data.List: rows <-> cols, RAGGED-SAFE like Haskell
    xss = [list(xs) for xs in xss]     # any iterables in (rows may be one-shot); materialise once
    n = max(map(len, xss), default=0)  # short rows just drop out of later columns (no truncation)
    return [[xs[i] for xs in xss if i < len(xs)] for i in range(n)]
\end{lstlisting}

\noindent\textit{At the REPL:}
\begin{lstlisting}[style=spec]
>>> nub([3, 1, 3, 2, 1])
[3, 1, 2]
>>> transpose([[1, 2, 3], [4, 5, 6]])
[[1, 4], [2, 5], [3, 6]]
>>> pairwise([1, 4, 9])
[(1, 4), (4, 9)]
>>> stripPrefix("re", "rebuild")
'build'
\end{lstlisting}
\chapter{Lists II: Higher-Order}

Where loops go to be named. \texttt{map\_} transforms, \texttt{filter\_}
selects, \texttt{concat} flattens one level, and \texttt{concatMap} does both
at once -- it is the backbone of enumeration (``for every choice, produce all
results, pooled''). \texttt{find} is the early exit folds cannot perform: it is
lazy, works on infinite streams, and returns \texttt{None} on failure, feeding
\texttt{fromMaybe}. \texttt{mapMaybe} is the messy-input workhorse -- map a
parser that returns value-or-\texttt{None}, drop the Nothings, one pass. When
input is dirty, reach for it before you reach for a loop with an \texttt{if}
and a flag.
\begin{lstlisting}[style=code]
# ============ 5. higher-order lists ============

# Data.Maybe mapMaybe: map and drop the Nothings in ONE pass -- the parse-and-filter shape
mapMaybe  = lambda f, xs: [y for y in (f(x) for x in xs) if y is not None]
map_      = lambda f, xs: [f(x) for x in xs]
filter_   = lambda pred, xs: [x for x in xs if pred(x)]
# Data.List find: lazy first-match -- folds can't stop early, find can
find      = lambda pred, xs: next((x for x in xs if pred(x)), None)   # first match, else None
catMaybes = lambda xs: [x for x in xs if x is not None]                      # mapMaybe id
concat    = lambda xss: [x for xs in xss for x in xs]
concatMap = lambda f, xs: [y for x in xs for y in f(x)]
starmap   = lambda f, pairs: [f(*p) for p in pairs]   # map . uncurry
zipWith   = lambda f, a, b: [f(x, y) for x, y in zip(a, b)]
zipWith3  = lambda f, a, b, c: [f(x, y, z) for x, y, z in zip(a, b, c)]
cross     = lambda a, b: [(x, y) for x in a for y in b]     # cartesian (`product` names the fold)
def partition(pred, xs):                   # (keepers, rest) in ONE pass; pred called once per element
    yes, no = [], []
    for x in xs:
        (yes if pred(x) else no).append(x)
    return (yes, no)
\end{lstlisting}

\noindent\textit{At the REPL:}
\begin{lstlisting}[style=spec]
>>> concatMap(lambda w: [w, w.upper()], ["a", "b"])
['a', 'A', 'b', 'B']
>>> find(lambda x: x % 7 == 0, count(1))     # lazy: an infinite stream
7
>>> mapMaybe(lambda s: int(s) if s.isdigit() else None, ["3", "x", "7"])
[3, 7]
\end{lstlisting}
\chapter{Folds}

The fold is the master pattern: an accumulator, a combining function, one pass.
\texttt{foldl} is the definition -- modern Python has no builtin fold (its
\texttt{reduce} was banished to \texttt{functools}), so the prelude keeps only
the Haskell name -- functools can wait for days when imports return. Study
\texttt{foldl}'s handling of the missing initial value via
\texttt{NOTHING}. \texttt{foldr} is derived from
\texttt{foldl} by flipping and reversing -- an identity that is only valid
because Python is strict, and stated in the file header so you never forget the
caveat. \texttt{all}/\texttt{any} are the Boolean folds, used bare as builtins
over Booleans; \texttt{compose} folds a tuple of functions into one.

The section ends with the \emph{enumeration folds} -- brute-force licenses.
\texttt{subsequences} builds the powerset by doubling the accumulator per
element ($2^n$: fine to $n \approx 20$, and you should say that bound out
loud); \texttt{replicateM} builds every length-$n$ word over an alphabet
($|xs|^n$). When $n$ is small, honest enumeration beats clever and fragile.
\begin{lstlisting}[style=code]
# ============ 6. folds ============ (collapse a list to one value)

def compose(*fns): # foldr (.) id -- RIGHTMOST runs first; constant stack depth
    def composed(x):
        for f in reversed(fns):
            x = f(x)
        return x
    return composed

def foldl(f, xs, init=NOTHING):  # THE left fold; Python buried its own in functools as reduce
    it = iter(xs)
    if init is NOTHING:
        try:
            acc = next(it)
        except StopIteration:
            raise TypeError("fold of empty sequence with no initial value")
    else:
        acc = init
    for x in it:
        acc = f(acc, x)
    return acc

foldl1  = lambda f, xs: foldl(f, xs)
foldr   = lambda f, xs, init: foldl(flip(f), reversed(list(xs)), init)
foldr1  = lambda f, xs: foldl(flip(f), reversed(list(xs)))
product = lambda xs: foldl(lambda a, x: a * x, xs, 1)   # the numeric fold

# enumeration folds -- brute-force licenses for small n (say the bound out loud):
# Data.List subsequences: the powerset; each element DOUBLES acc -- 2^n, fine to n ~ 20
subsequences = lambda xs: foldl(lambda acc, x: acc + [s + [x] for s in acc], list(xs), [[]])
# Control.Monad replicateM: every length-n word over xs -- cross, n times; |xs|^n
replicateM   = lambda n, xs: foldl(lambda acc, _: [s + [x] for s in acc for x in xs], range(n), [[]])
\end{lstlisting}

\noindent\textit{At the REPL:}
\begin{lstlisting}[style=spec]
>>> foldl(lambda a, x: a * 10 + x, [1, 2, 3], 0)
123
>>> compose(str.strip, str.upper)("  hi  ")
'HI'
>>> subsequences([1, 2])
[[], [1], [2], [1, 2]]
>>> replicateM(2, "ab")
[['a', 'a'], ['a', 'b'], ['b', 'a'], ['b', 'b']]
\end{lstlisting}
\chapter{Scans, Streams, and Finite Views}

A scan is a fold that keeps its history: \texttt{scanl} hands back every
intermediate accumulator. Prefix sums are \texttt{scanl (+) 0}; running
minima are \texttt{scanl1 min}; the depth of nested parentheses is a scan over
$\pm1$. When you need ``the state at every step'', do not run $n$ folds -- run
one scan.

Section~8 turns to streams: generators that may never end.
\texttt{count}, \texttt{repeat}, \texttt{cycle}, \texttt{iterate} produce;
nothing is computed until something consumes. \texttt{unfoldr} is
\texttt{iterate}'s finite twin -- \texttt{f(seed)} returns \texttt{(value,
seed')} or \texttt{None} to stop -- and it ends the habit of threading
\texttt{(state, acc)} tuples through loops. Section~9 supplies the consumers
and finite views: \texttt{take}, \texttt{takeWhile}/\texttt{dropWhile},
\texttt{span}/\texttt{break\_} for splitting at a condition (one pass, safe
on generators), and the view
factories \texttt{inits}, \texttt{tails} (every prefix and suffix -- substring
problems start there) and \texttt{chunksOf}. Note the deliberate naming: these
carry Python's own \texttt{itertools} names with Haskell semantics, so both
vocabularies stay warm.
\begin{lstlisting}[style=code]
# ============ 7. scans ============ (a fold that keeps its history)

def accumulate(xs, f=None, initial=NOTHING):     # scanl / scanl1
    if f is None:
        f = lambda a, b: a + b
    it = iter(xs)
    if initial is NOTHING:
        try:
            acc = next(it)
        except StopIteration:
            return
    else:
        acc = initial
    yield acc
    for x in it:
        acc = f(acc, x)
        yield acc

scanl  = lambda f, z, xs: list(accumulate(xs, f, initial=z))
scanl1 = lambda f, xs: list(accumulate(xs, f))
scanr  = lambda f, z, xs: list(reversed(scanl(flip(f), z, list(reversed(list(xs))))))
scanr1 = lambda f, xs: list(reversed(scanl1(flip(f), list(reversed(list(xs))))))
\end{lstlisting}
\begin{lstlisting}[style=code]
# ============ 8. streams ============ (lazy, possibly infinite)

def count(start=0, step=1):                 # [start, start+step ..]
    n = start
    while True:
        yield n
        n += step

def repeat(x, n=None):
    if n is None:
        while True:
            yield x
    else:
        for _ in range(n):
            yield x

def cycle(xs):
    xs = list(xs)
    if not xs:
        raise ValueError("cycle: empty list")   # Haskell errors too; the alternative is a silent hang
    while True:
        yield from xs

def iterate(f, x):                          # iterate f x = [x, f x, f (f x), ..]
    while True:
        yield x
        x = f(x)

def unfoldr(f, seed):                       # Data.List unfoldr: f(seed) -> None | (value, seed')
    out = []                                # iterate's finite twin: grows a list until f says stop --
    step = f(seed)                          # no more threading (state, acc) tuples through until
    while step is not None:
        val, seed = step
        out.append(val)
        step = f(seed)
    return out

def chain(*iterables):                      # concat
    for it in iterables:
        yield from it
\end{lstlisting}
\begin{lstlisting}[style=code]
# ============ 9. slicing & spans ============ (finite views of streams)

def islice(iterable, stop):                 # take
    it = iter(iterable)
    for _ in range(stop):
        try:
            yield next(it)
        except StopIteration:
            return

take      = lambda n, xs: list(islice(iter(xs), n))           # works on infinite streams

def takewhile(pred, xs):
    for x in xs:
        if not pred(x):
            return
        yield x

def takeuntil(pred, xs):                    # like takewhile(not . pred), but INCLUDES the first hit
    for x in xs:
        yield x
        if pred(x):
            return

def dropwhile(pred, xs):
    it = iter(xs)
    for x in it:
        if not pred(x):
            yield x
            break
    yield from it

takeWhile = lambda p, xs: list(takewhile(p, xs))
dropWhile = lambda p, xs: list(dropwhile(p, xs))
def span(p, xs):                            # split at first failure, ONE pass; safe on one-shot iters
    xs = list(xs)
    i = next((i for i, x in enumerate(xs) if not p(x)), len(xs))
    return (xs[:i], xs[i:])
break_    = lambda p, xs: span(lambda x: not p(x), xs)
inits     = lambda xs: [xs[:i] for i in range(len(xs) + 1)]   # Data.List: every prefix, [] first
tails     = lambda xs: [xs[i:] for i in range(len(xs) + 1)]   # every suffix; substrings start here
chunksOf  = lambda n, xs: [xs[i:i + n] for i in range(0, len(xs), n)]   # Data.List.Split; short last ok
\end{lstlisting}

\noindent\textit{At the REPL:}
\begin{lstlisting}[style=spec]
>>> scanl(lambda a, b: a + b, 0, [3, 1, 4])
[0, 3, 4, 8]
>>> take(5, iterate(lambda x: 2 * x, 1))
[1, 2, 4, 8, 16]
>>> unfoldr(lambda m: None if m == 0 else (m % 10, m // 10), 407)
[7, 0, 4]
>>> tails("abc")
['abc', 'bc', 'c', '']
\end{lstlisting}
\chapter{Order and Text}

Sorting with keys is the single most-used line in any timed test.
\texttt{sortOn} evaluates the key once per element; tuple keys give multi-level
orderings, and the \texttt{(-count, name)} idiom gets ``count descending, ties
alphabetical'' in one shot -- something \texttt{reverse=True} cannot do.
\texttt{minOn}/\texttt{maxOn} answer ``best element by key'' in $O(n)$;
sorting first to take the head is paying $O(n \log n)$ for the same
information. \texttt{merge} joins two sorted lists stably -- the heart of merge
sort -- and the \texttt{bisect} pair performs binary search over sorted data:
\texttt{bisect\_left} finds the first slot $\ge x$, \texttt{bisect\_right} the
first $> x$; their difference counts occurrences. Section~11 is four honest
aliases (\texttt{words}, \texttt{unwords}, \texttt{lines}, \texttt{unlines})
that make string pipelines read like Haskell.
\begin{lstlisting}[style=code]
# ============ 10. sorting & searching ============

sortOn = lambda f, xs: sorted(xs, key=f)    # sortOn (schwartzian, f called once per element)
minOn  = lambda f, xs: min(xs, key=f)    # best by key in O(n) -- sortOn + head without the sort
maxOn  = lambda f, xs: max(xs, key=f)    # (minimumBy/maximumBy + comparing; first wins ties)

def cmp_to_key(cmp):                        # resurrect Python 2's cmp for sortBy
    class K:
        def __init__(self, x): self.x = x
        def __lt__(self, other): return cmp(self.x, other.x) < 0
    return K

sortBy = lambda cmp, xs: sorted(xs, key=cmp_to_key(cmp))   # comparator sort, Python 2 style

def merge(a, b):                            # merge two sorted lists, stable; the heart of merge sort
    out, i, j = [], 0, 0                    # (section 13's merge_lists is this same loop on ListNodes)
    while i < len(a) and j < len(b):
        if a[i] <= b[j]:
            out.append(a[i]); i += 1
        else:
            out.append(b[j]); j += 1
    return out + a[i:] + b[j:]

def bisect_left(a, x):                      # first index where a[i] >= x
    lo, hi = 0, len(a)
    while lo < hi:
        mid = (lo + hi) // 2
        if a[mid] < x:
            lo = mid + 1
        else:
            hi = mid
    return lo

def bisect_right(a, x):                     # first index where a[i] > x
    lo, hi = 0, len(a)
    while lo < hi:
        mid = (lo + hi) // 2
        if a[mid] <= x:
            lo = mid + 1
        else:
            hi = mid
    return lo
\end{lstlisting}
\begin{lstlisting}[style=code]
# ============ 11. strings ============

words   = str.split
unwords = " ".join
lines   = str.splitlines
unlines = lambda ls: "".join(l + "\n" for l in ls)
\end{lstlisting}

\noindent\textit{At the REPL:}
\begin{lstlisting}[style=spec]
>>> sortOn(lambda w: (len(w), w), ["bb", "a", "ccc", "aa"])
['a', 'aa', 'bb', 'ccc']
>>> minOn(abs, [-3, 2, 5])
2
>>> bisect_right([1, 3, 3, 5], 3) - bisect_left([1, 3, 3, 5], 3)
2
>>> unwords(words("  such   spacing  "))
'such spacing'
\end{lstlisting}
\chapter{Containers: Dicts, Bags, Trees, Heaps}

The longest section, and the richest. \texttt{defaultdict} rebuilds
autovivification from one dunder: touch a missing key and the factory fills it.
\texttt{Counter} is its \texttt{int} instance. Then the dict-building folds:
\texttt{insertWith} combines on collision, \texttt{fromListWith} folds a pair
list into a dict, combining collisions with f(new, old) -- Haskell's order, so
groups keep arrival order when you combine old + new -- and \texttt{unionWith}
merges two dicts under any combining function. \texttt{insertWith} is
persistent, returning a fresh dict like the immutable original. The bag
operations show the same shape specialised to
multisets -- and comparing \texttt{bag\_union} (\texttt{max}) with a summing
merge (\texttt{+}) is a first lesson in monoids.

\texttt{Tree} and \texttt{ITree} nest defaultdicts into tries;
\texttt{paths} folds a whole tree to \texttt{(keypath, leaf)} pairs;
\texttt{setpath} writes (autovivifying), \texttt{getpath} reads
\emph{without} autovivifying -- reading with \texttt{[]} would plant ghost
branches. The section closes with three machines built from scratch: a
two-stack \texttt{deque} with all four ends amortised $O(1)$, union--find with
path halving, and a binary heap as two functions on a plain list. Owning these
means never being helpless when a ``batteries included'' import is off the
table.
\begin{lstlisting}[style=code]
# ============ 12. containers ============

class defaultdict(dict):
    def __init__(self, factory=None, *args, **kw):
        super().__init__(*args, **kw)
        self.factory = factory
    def __missing__(self, key):
        if self.factory is None:
            raise KeyError(key)
        self[key] = self.factory()
        return self[key]

def Counter(xs):                            # count-by-value; a defaultdict(int) fold
    d = defaultdict(int)
    for x in xs:
        d[x] += 1
    return d

def insertWith(f, k, v, d):                 # Data.Map insertWith f k v m: PERSISTENT -- returns a new
    out = dict(d)                           # dict untouched; collision -> f(new, old), Haskell order
    out[k] = f(v, out[k]) if k in out else v   # order (new value first)
    return out

def fromListWith(f, pairs):                 # Data.Map fromListWith: THE dict-building fold; f(new, old)
    d = {}                                  # like insertWith -- so grouping with concat REVERSES each
    for k, v in pairs:                      # group (the classic Haskell gotcha); use add for counts
        d[k] = f(v, d[k]) if k in d else v
    return d                                # Counter == fromListWith(add) over (x, 1);
                                            # grouping == fromListWith(++) over (k, [v])

def unionWith(f, a, b):                     # Data.Map unionWith: merge two dicts, f on shared keys
    out = dict(a)                           # f(a's value, b's value) -- Haskell's left/right order
    for k, v in b.items():
        out[k] = f(out[k], v) if k in out else v
    return out                              # one combinator, many merges:
                                            # bag_union == unionWith(max), merge-sum == unionWith(add)

def groupBy(eq, xs):                        # Data.List groupBy: runs of CONSECUTIVE elements; each new
    out = []                                # element is compared via eq against the run's FIRST element,
    for x in xs:                            # exactly like Haskell (span-based) -- not its neighbor
        if out and eq(out[-1][0], x):
            out[-1].append(x)
        else:
            out.append([x])
    return out

group = lambda xs: groupBy(lambda a, b: a == b, xs)   # Data.List group: runs of equals, [[a]] -- no keys

# multiset ops on Counters (use .get, not [] -- indexing a Counter autovivifies zeros):
bag_union = lambda a, b: {k: max(a.get(k, 0), b.get(k, 0)) for k in a.keys() | b.keys()}
bag_inter = lambda a, b: {k: v for k in a.keys() & b.keys() if (v := min(a.get(k, 0), b.get(k, 0)))}
bag_diff  = lambda a, b: {k: a[k] - b[k] for k in a if a[k] > b.get(k, 0)}    # clipped at 0
bag_sub   = lambda a, b: all(b.get(k, 0) >= n for k, n in a.items())          # a <= b (inclusion)

Tree  = lambda depth, leaf: (defaultdict(leaf) if depth == 1
                             else defaultdict(lambda: Tree(depth - 1, leaf)))
ITree = lambda: defaultdict(ITree)

def paths(t):                               # cata for Tree/ITree: [(keypath, leaf)]
    # recursion depth = tree height (word length / len(nums)) -- well under the ~1000 limit
    if not isinstance(t, dict) or not t:    # leaf = non-dict value, or empty node
        return [([], t)]
    return [([k] + p, leaf) for k, sub in t.items() for p, leaf in paths(sub)]

leaves = lambda t: [l for _, l in paths(t)]

def setpath(t, ks, v):                      # write leaf v at key path ks (autovivifies interior)
    for k in ks[:-1]:
        t = t[k]
    t[ks[-1]] = v                           # NB: clobbers any subtree already at ks

def getpath(t, ks, default=None):           # read along ks WITHOUT autovivifying
    for k in ks:
        if not isinstance(t, dict) or k not in t:
            return default
        t = t[k]
    return t

class deque:                                # two-stack; all four ends amortized O(1)
    def __init__(self, xs=()):
        self._in, self._out = [], list(reversed(list(xs)))
    def append(self, x): self._in.append(x)
    def appendleft(self, x): self._out.append(x)
    def pop(self):
        if not self._in:
            self._in, self._out = list(reversed(self._out)), []
        return self._in.pop()
    def popleft(self):
        if not self._out:
            self._out, self._in = list(reversed(self._in)), []
        return self._out.pop()
    def __len__(self): return len(self._in) + len(self._out)
    def __iter__(self): return iter(self._out[::-1] + self._in)

def dsu(n):                                 # union-find, path halving; union -> False if already joined
    parent = list(range(n))
    def find(x):
        while parent[x] != x:
            parent[x] = parent[parent[x]]
            x = parent[x]
        return x
    def union(a, b):
        ra, rb = find(a), find(b)
        if ra == rb:
            return False
        parent[ra] = rb
        return True
    return find, union

def heappush(h, x):                         # min-heap on a plain list: sift-up
    h.append(x)
    i = len(h) - 1
    while i and h[(i - 1) // 2] > h[i]:
        h[(i - 1) // 2], h[i] = h[i], h[(i - 1) // 2]
        i = (i - 1) // 2

def heappop(h):                             # min-heap on a plain list: sift-down
    h[0], h[-1] = h[-1], h[0]
    top = h.pop()
    i, n = 0, len(h)
    while True:
        s, l, r = i, 2*i + 1, 2*i + 2
        if l < n and h[l] < h[s]: s = l
        if r < n and h[r] < h[s]: s = r
        if s == i:
            return top
        h[i], h[s] = h[s], h[i]
        i = s
\end{lstlisting}

\noindent\textit{At the REPL:}
\begin{lstlisting}[style=spec]
>>> fromListWith(lambda a, b: a + b, [("a", [1]), ("b", [2]), ("a", [3])])
{'a': [1, 3], 'b': [2]}
>>> unionWith(max, Counter("aab"), Counter("abb")) == {'a': 2, 'b': 2}
True
>>> t = ITree(); setpath(t, list("hi") + ["$"], "hi"); paths(t)
[(['h', 'i', '$'], 'hi')]
>>> h = []; heappush(h, 3); heappush(h, 1); heappop(h)
1
\end{lstlisting}
\chapter{Nodes, Grids, Windows, Control}

Section~13 covers the LeetCode givens. \texttt{TreeNode} with \texttt{tfold} --
a catamorphism -- makes every traversal a one-liner: in-, pre-, post-order,
depth and size are just different combining functions. \texttt{levelorder} is
BFS, the deque earning its keep. The linked-list block collects the pointer
idioms worth knowing cold: build by \texttt{foldr}, reverse with three
pointers, fast/slow for the middle, Floyd for cycles, dummy-head for merges.

Section~14 gives grids their neighbour generators (bounds-checked, so BFS
bodies stay clean) and \texttt{longest\_window}, the sliding-window skeleton:
you supply \texttt{add}, \texttt{rem} and \texttt{valid} as closures and it
does the two-pointer bookkeeping. Section~15 is control: \texttt{memo},
an unbounded memo in eight lines (with the cache exposed, so you can watch a DP
table fill), and \texttt{until}, the fixed-point loop -- iteration without
recursion depth limits.
\begin{lstlisting}[style=code]
# ============ 13. nodes ============ (the LeetCode givens: binary trees & linked lists)

class TreeNode:                             # the LeetCode binary tree
    def __init__(self, val=0, left=None, right=None):
        self.val, self.left, self.right = val, left, right

def tfold(f, z, t):                         # cata: f(val, l_acc, r_acc); depth = tree height
    return z if t is None else f(t.val, tfold(f, z, t.left), tfold(f, z, t.right))

inorder   = lambda t: tfold(lambda v, l, r: l + [v] + r, [], t)
preorder  = lambda t: tfold(lambda v, l, r: [v] + l + r, [], t)
postorder = lambda t: tfold(lambda v, l, r: l + r + [v], [], t)
tdepth    = lambda t: tfold(lambda _, l, r: 1 + max(l, r), 0, t)
tsize     = lambda t: tfold(lambda _, l, r: 1 + l + r, 0, t)

def levelorder(t):                          # BFS: the deque earning its keep
    out, q = [], deque([t] if t else [])
    while len(q):
        n = q.popleft()
        out.append(n.val)
        if n.left:  q.append(n.left)
        if n.right: q.append(n.right)
    return out

class ListNode:                             # the LeetCode linked list
    def __init__(self, val=0, next=None):
        self.val, self.next = val, next

def from_list(xs):                          # build; foldr of ListNode
    node = None
    for x in reversed(list(xs)):
        node = ListNode(x, node)
    return node

def to_list(node):                          # unfold back to a Python list
    out = []
    while node:
        out.append(node.val)
        node = node.next
    return out

def reverse_list(node):                     # three-pointer; prev is a fold accumulator
    prev = None
    while node:
        node.next, prev, node = prev, node, node.next
    return prev

def middle(node):                           # fast/slow: second middle for even lengths
    slow = fast = node
    while fast and fast.next:
        slow, fast = slow.next, fast.next.next
    return slow

def has_cycle(node):                        # Floyd: fast laps slow iff a cycle exists
    slow = fast = node
    while fast and fast.next:
        slow, fast = slow.next, fast.next.next
        if slow is fast:
            return True
    return False

def merge_lists(a, b):                      # merge two sorted lists; dummy-head idiom
    dummy = tail_ = ListNode()
    while a and b:
        if a.val <= b.val:
            tail_.next, a = a, a.next
        else:
            tail_.next, b = b, b.next
        tail_ = tail_.next
    tail_.next = a or b
    return dummy.next
\end{lstlisting}
\begin{lstlisting}[style=code]
# ============ 14. grids & windows ============

def neighbors4(r, c, R, C):                 # in-bounds orthogonal neighbors
    for dr, dc in ((1, 0), (-1, 0), (0, 1), (0, -1)):
        if 0 <= r + dr < R and 0 <= c + dc < C:
            yield r + dr, c + dc

def neighbors8(r, c, R, C):                 # ... plus diagonals
    for dr in (-1, 0, 1):
        for dc in (-1, 0, 1):
            if (dr or dc) and 0 <= r + dr < R and 0 <= c + dc < C:
                yield r + dr, c + dc

def longest_window(xs, valid, add, rem):    # sliding-window skeleton; state lives in the closures
    lo = best = 0
    for hi in range(len(xs)):
        add(xs[hi])
        while not valid():                  # shrink until legal again
            rem(xs[lo])
            lo += 1
        best = max(best, hi - lo + 1)
    return best
\end{lstlisting}
\begin{lstlisting}[style=code]
# ============ 15. control ============

def memo(f):                                # unbounded memoizer; enough for DP (no eviction by design)
    cache = {}
    def wrapped(*args, **kw):
        key = (args, frozenset(kw.items())) if kw else args
        if key not in cache:
            cache[key] = f(*args, **kw)
        return cache[key]
    wrapped.cache = cache                   # peek at the DP table if curious
    return wrapped

def until(p, f, x):                         # until p f x -- loop, not recursion (unbounded depth)
    while not p(x):
        x = f(x)
    return x
\end{lstlisting}

\noindent\textit{At the REPL:}
\begin{lstlisting}[style=spec]
>>> t = TreeNode(2, TreeNode(1), TreeNode(3))
>>> inorder(t), tdepth(t), tsize(t)
([1, 2, 3], 2, 3)
>>> to_list(reverse_list(from_list([1, 2, 3])))
[3, 2, 1]
>>> sorted(neighbors4(0, 0, 2, 2))
[(0, 1), (1, 0)]
>>> until(lambda x: x > 100, lambda x: x * 2, 1)
128
\end{lstlisting}
\chapter{Deep Dives}
The chapters above teach the prelude a section at a time; this one slows down
for the twenty-five names where the concept is the hard part. Each dive keeps
one shape: how the definition actually works (often with the evaluation
traced), why it earns its place, the Haskell connection where there is one,
and the sibling to contrast it with -- ending in one sentence built to be
remembered. Every example is a doctest from \texttt{prelude-doctests.py}, so
everything here runs.

\section{flip}
\noindent\textit{[combinators]} --- swaps the two arguments of a function. An adapter: the function you have takes (a, b), the slot you are filling supplies (b, a). This is how the prelude derives foldr from foldl.
\begin{lstlisting}[style=spec]
>>> sub = lambda a, b: a - b
>>> sub(10, 1), flip(sub)(10, 1)
(9, -9)
>>> flip(pow)(3, 2)                     # pow(2, 3)
8
\end{lstlisting}

How it works: flip(f) returns a new function that hands its two arguments to f in the opposite order. Nothing is computed at flip time -- it is pure plumbing, applied when the data arrives in the "wrong" order for the function you already have.

\begin{lstlisting}[style=spec]
>>> pair = lambda a, b: (a, b)
>>> flip(pair)(1, 2)
(2, 1)
\end{lstlisting}

Why it earns its place: the prelude's own foldr is the showcase. foldl's combiner takes (accumulator, element); a right fold's combiner takes (element, accumulator). Rather than build a second fold machine, foldr = foldl of the FLIPPED combiner over the reversed list -- one adapter turned an existing engine into its mirror image.

\begin{lstlisting}[style=spec]
>>> f = lambda x, acc: '(' + x + acc + ')'
>>> foldr(f, "abc", "*") == foldl(flip(f), list(reversed("abc")), "*")
True
\end{lstlisting}

Remember it as: when argument order is the only thing wrong, adapt -- don't rewrite.


Roll your own --- the definition:

\begin{lstlisting}[style=code]
flip = lambda f: (lambda x, y: f(y, x))  # flip f x y = f y x
\end{lstlisting}

\section{curry}
\noindent\textit{[combinators]} --- splits a two-argument function into two one-argument stages, so you can supply the first argument now and the second later.
\begin{lstlisting}[style=spec]
>>> add = lambda a, b: a + b
>>> curry(add)(1)(2)
3
>>> add_ten = curry(add)(10)
>>> add_ten(5)
15
\end{lstlisting}

How it works: read the three lambdas inside-out, then watch a call peel one layer at a time. Each call binds ONE argument by closure -- after curry(add)(10), the 10 lives in the returned function's captured environment, and that half-applied function is a value you can name, store, and pass.

\begin{lstlisting}[style=spec]
>>> curry(add)(10)(5)
15
>>> add_ten = curry(add)(10)
>>> add_ten(5), add_ten(90)
(15, 100)
\end{lstlisting}

Why it earns its place: the prelude's pipelines -- map\_, filter\_, find, compose -- all want ONE-argument functions in their function slot, while real logic usually takes two ("multiply by k", "is x greater than n?"). Currying is the adapter: close over the configuration argument now, hand the pipeline the unary stage it wants.

\begin{lstlisting}[style=spec]
>>> gt = curry(lambda n, x: x > n)
>>> filter_(gt(10), [4, 11, 8, 40])
[11, 40]
>>> scale = curry(lambda k, x: k * x)
>>> map_(scale(3), [1, 2, 3])
[3, 6, 9]
\end{lstlisting}

The Haskell connection: in Haskell EVERY function is curried automatically -- add :: Int -> Int -> Int really means Int -> (Int -> Int), so add 10 is ordinary code and map (add 10) xs falls out for free. Every arrow in a Haskell type is another layer of this doll.


Versus partial: both bind early. partial binds any number of arguments, all at once, when you have them; curry restructures the function so binding happens one call at a time, decided later by whoever holds it. In Python, partial is the cheaper everyday tool; curry is the pipeline-shaper and the reading key for Haskell signatures.


Remember it as: curry turns one two-argument function into two one-argument functions -- and one-argument functions are the only thing pipelines eat.


Roll your own: three lambdas, each returning the next -- every layer is a closure capturing one more argument until the innermost body finally has both. No state, no storage: the captured environment IS the data structure.

\begin{lstlisting}[style=code]
curry = lambda f: lambda x: lambda y: f(x, y)
\end{lstlisting}

\section{on}
\noindent\textit{[combinators]} --- builds a two-argument function that first sends BOTH arguments through a key function: on(f, key)(x, y) == f(key(x), key(y)). Its natural habitat is comparators: compare two records by one field.
\begin{lstlisting}[style=spec]
>>> on(add, len)("ab", "c")             # len("ab") + len("c")
3
>>> same_length = on(lambda a, b: a == b, len)
>>> same_length("abc", "xyz"), same_length("ab", "abc")
(True, False)
>>> sortBy(on(lambda a, b: a - b, snd), [(1, 9), (2, 3)])
[(2, 3), (1, 9)]
\end{lstlisting}

How it works: on(f, key) builds a two-argument function that routes BOTH inputs through key before combining them with f: on(f, key)(x, y) == f(key(x), key(y)). Combination logic and projection logic stay separate, so each is reusable on its own.

\begin{lstlisting}[style=spec]
>>> shorter = on(lambda a, b: a < b, len)
>>> shorter("ab", "abc"), shorter("abc", "ab")
(True, False)
\end{lstlisting}

Why it earns its place: comparator factories. Any "compare records by one field" collapses to on(compare, field) -- Haskell's sortBy (compare 'on' snd) is exactly sortBy(on(cmp, snd)) here, and reads as its own documentation.

\begin{lstlisting}[style=spec]
>>> sortBy(on(lambda a, b: a - b, snd), [('a', 9), ('b', 1), ('c', 5)])
[('b', 1), ('c', 5), ('a', 9)]
\end{lstlisting}

Versus sortOn: when a per-element key exists, sortOn(key, xs) is simpler and faster (the key runs once per element). on shines when an API demands a genuine two-argument comparator, or when the combining function is something other than comparison.


Remember it as: compare the keys, not the things.


Roll your own --- the definition:

\begin{lstlisting}[style=code]
on = lambda f, g: lambda x, y: f(g(x), g(y))  # Data.Function: combine via a key -- cmp `on` fst
\end{lstlisting}

\section{NOTHING}
\noindent\textit{[combinators]} --- a private sentinel object meaning "no argument was supplied". Why not None? Because None is a legitimate VALUE a caller might pass; NOTHING can only mean "omitted", since nobody else has a reference to it. Always test with 'is'. This is how foldl and accumulate tell an absent initial value from a real one.
\begin{lstlisting}[style=spec]
>>> x = NOTHING
>>> x is NOTHING
True
>>> def greet(name=NOTHING):
...     return "hello, stranger" if name is NOTHING else "hello, " + name
>>> greet()
'hello, stranger'
>>> greet("ada")
'hello, ada'
\end{lstlisting}

How it works: object() creates a fresh, featureless object whose only property is its IDENTITY -- nothing else in the program can ever be it. Testing x is NOTHING therefore asks exactly one question: "was I handed this precise sentinel?", which makes it a safe flag for "no argument was supplied".


Why not None, demonstrated: None is data (fromMaybe's whole job is handling it), so using it to also mean "argument omitted" would blur two different situations. The sentinel keeps them apart: None is Nothing in the DATA; NOTHING is Nothing in the ARGUMENT LIST.

\begin{lstlisting}[style=spec]
>>> def describe(x=NOTHING):
...     if x is NOTHING:
...         return "no argument"
...     return "got " + repr(x)
>>> describe()
'no argument'
>>> describe(None)
'got None'
\end{lstlisting}

Where the prelude leans on it: foldl and accumulate take an optional seed. With NOTHING as the default they can tell "seed omitted -- use the first element" from "the seed IS None", which a None default could not.


Remember it as: None is a value; NOTHING is the absence of one.


Roll your own --- the definition:

\begin{lstlisting}[style=code]
NOTHING = object()  # Maybe's Nothing: "no arg given"; test with `is`
\end{lstlisting}

\section{fromMaybe}
\noindent\textit{[combinators]} --- lands a default: fromMaybe(d, x) is x unless x is None, in which case it is d. The standard way to finish a pipeline built on the None-is-Nothing convention (find, lookup, getpath, stripPrefix all return None on failure).
\begin{lstlisting}[style=spec]
>>> fromMaybe(0, None), fromMaybe(0, 5)
(0, 5)
>>> fromMaybe('_', find(lambda c: c == 'q', "abc"))
'_'
>>> fromMaybe(99, lookup('b', [('a', 1), ('b', 2)]))
2
\end{lstlisting}

Start from a problem you already have: find looks for something that might not be there, and when it is not, you are left holding None -- which you cannot print to a user, add, or put in a report. So you write a backup plan:

\begin{lstlisting}[style=spec]
>>> result = find(even, [3, 5, 7])
>>> if result is None:
...     result = 0
>>> result
0
\end{lstlisting}

fromMaybe is those three lines as one word. Argument order: the backup plan first, then the risky thing -- said aloud it parses: "use 0 if this comes back empty".

\begin{lstlisting}[style=spec]
>>> fromMaybe(0, find(even, [3, 5, 7]))
0
>>> fromMaybe(0, find(even, [3, 5, 8]))
8
\end{lstlisting}

Versus x or d: Python's or falls back on ANYTHING falsy -- 0, "", [], False -- not just None, so it silently destroys legitimate zero-ish answers. fromMaybe asks the only correct question: is this the absence marker?

\begin{lstlisting}[style=spec]
>>> fromMaybe(99, 0), (0 or 99)
(0, 99)
\end{lstlisting}

Remember it as: ifNoneUse -- the three-line backup plan as one word; or-with-a-conscience.


Roll your own --- the definition:

\begin{lstlisting}[style=code]
fromMaybe = lambda d, x: d if x is None else x  # Data.Maybe: default for every None-returning tool
\end{lstlisting}

\section{isNothing}
\noindent\textit{[combinators]} --- the mirror: True only for None. Counting failures reads as one word.
\begin{lstlisting}[style=spec]
>>> isNothing(None), isNothing(0)
(True, False)
>>> len(filter_(isNothing, [None, 3, None]))
2
\end{lstlisting}

Together they finish the Maybe family: fromMaybe LANDS an optional value, mapMaybe/catMaybes FILTER optional values, isJust/isNothing TEST one. Haskell's Data.Maybe has all five under the same names.

\begin{lstlisting}[style=spec]
>>> readings = [(1, None), (2, 7.5), (3, None)]
>>> fromMaybe(-1, find(isJust, map(snd, readings)))
7.5
\end{lstlisting}

Roll your own --- the definition:

\begin{lstlisting}[style=code]
isNothing = lambda x: x is None
\end{lstlisting}

\section{isqrt}
\noindent\textit{[arithmetic \& logic]} --- floor square root using Newton's method on INTEGERS -- no floats, so no precision loss on huge numbers; a negative input raises rather than lying.
\begin{lstlisting}[style=spec]
>>> isqrt(16), isqrt(17)
(4, 4)
>>> isqrt(10**18)
1000000000
\end{lstlisting}

Roll your own: Newton's iteration on integers -- repeatedly replace x with the average of x and n // x. It converges from above, so loop while the new estimate still shrinks; all arithmetic is //, so no float (and no float rounding lie) ever appears.

\begin{lstlisting}[style=code]
def isqrt(n):  # floor sqrt without floats (Newton)
    if n < 0:
        raise ValueError("isqrt of negative")
    x = n
    y = (x + 1) // 2
    while y < x:
        x, y = y, (y + n // y) // 2
    return x if n else 0
\end{lstlisting}

\section{transpose}
\noindent\textit{[list basics]} --- rows become columns, RAGGED-SAFE exactly like Haskell's: short rows simply drop out of later columns; nothing is silently truncated.
\begin{lstlisting}[style=spec]
>>> transpose([[1, 2, 3], [4, 5, 6]])
[[1, 4], [2, 5], [3, 6]]
>>> transpose([[1, 2, 3], [4, 5]])
[[1, 4], [2, 5], [3]]
>>> transpose([])
[]
\end{lstlisting}

Rows may be any iterables, even one-shot ones -- the rows are materialised first, so a generator of iterators transposes like a list of lists (without that step the generator is spent by the length scan):

\begin{lstlisting}[style=spec]
>>> transpose(iter([1, 2]) for _ in range(2))
[[1, 1], [2, 2]]
>>> transpose(["abc", "de"])
[['a', 'd'], ['b', 'e'], ['c']]
\end{lstlisting}

Roll your own: materialise the rows, find the longest, then for each column index keep xs[i] only from rows that reach it -- the 'if i < len(xs)' filter is precisely where short rows drop out of later columns instead of truncating everyone.

\begin{lstlisting}[style=code]
def transpose(xss):  # Data.List: rows <-> cols, RAGGED-SAFE like Haskell
    xss = [list(xs) for xs in xss]  # any iterables in (rows may be one-shot); materialise once
    n = max(map(len, xss), default=0)  # short rows just drop out of later columns (no truncation)
    return [[xs[i] for xs in xss if i < len(xs)] for i in range(n)]
\end{lstlisting}

\section{find}
\noindent\textit{[higher-order lists]} --- the first element satisfying the predicate, else None. Lazy: it stops at the first hit, which folds cannot do, and therefore works even on infinite streams. Finish with fromMaybe if you need a default.
\begin{lstlisting}[style=spec]
>>> find(even, [1, 3, 4, 6])
4
>>> find(even, [1, 3]) is None
True
>>> find(lambda x: x % 7 == 0, count(1))
7
\end{lstlisting}

How it works: find wraps a generator expression in next() with a None default. The generator is the point -- elements are tested one at a time, on demand, and the search STOPS at the first hit. Nothing after the match is ever examined.


Why it earns its place: folds and filter\_ consume the whole list; they cannot stop early. find can -- which makes it the right tool the moment the word "first" appears in a statement, and the only tool that survives an infinite stream. Land defaults with fromMaybe.

\begin{lstlisting}[style=spec]
>>> find(lambda x: x * x > 50, count(1))
8
>>> fromMaybe('?', find(str.isdigit, "abc7de"))
'7'
\end{lstlisting}

Versus filter\_: filter\_ answers "which ones?"; find answers "who was first, if anyone?". Taking filter\_(p, xs)[0] instead does wasted work and explodes when nothing matches; find does neither.


Remember it as: find is the early exit that folds cannot perform.


Roll your own: a generator expression handed to next() with a default. The generator is what makes it lazy -- elements are produced one at a time, so the first hit stops everything -- and the default is what makes it total.

\begin{lstlisting}[style=code]
find = lambda pred, xs: next((x for x in xs if pred(x)), None)  # first match, else None
\end{lstlisting}

\section{partition}
\noindent\textit{[higher-order lists]} --- split into (satisfy, don't) in ONE pass, both sides in original order; the predicate runs exactly once per element, so it may be expensive or even stateful.
\begin{lstlisting}[style=spec]
>>> partition(even, [1, 2, 3, 4])
([2, 4], [1, 3])
>>> partition(str.isalpha, "a1b2")
(['a', 'b'], ['1', '2'])
\end{lstlisting}

Roll your own: two buckets and one conditional append -- (yes if pred(x) else no).append(x). The conditional expression picks the LIST, then the append mutates it; one pass, exactly one predicate call per element.

\begin{lstlisting}[style=code]
def partition(pred, xs):  # (keepers, rest) in ONE pass; pred called once per element
    yes, no = [], []
    for x in xs:
        (yes if pred(x) else no).append(x)
    return (yes, no)
\end{lstlisting}

\section{mapMaybe}
\noindent\textit{[higher-order lists]} --- map a function that returns a value or None, and keep only the values: parse-and-filter in ONE pass. The shape for messy input -- write a parser that answers None for garbage, then mapMaybe it over everything. dict.get is already such a function.
\begin{lstlisting}[style=spec]
>>> mapMaybe(lambda s: int(s) if s.isdigit() else None, ["3", "x", "7"])
[3, 7]
>>> mapMaybe({'a': 1, 'b': 2}.get, ['a', 'x', 'b'])
[1, 2]
\end{lstlisting}

How it works: run f over every element; f answers a VALUE when it can make sense of the input and None when it cannot; keep only the values. Map and filter fused into one pass, with the parser itself deciding what survives.

\begin{lstlisting}[style=spec]
>>> parse_int = lambda s: int(s) if s.lstrip('-').isdigit() else None
>>> mapMaybe(parse_int, ["12", "x", "-3", ""])
[12, -3]
\end{lstlisting}

Why it earns its place: messy input. The alternative is a loop with an if, a flag and an append -- three chances for a bug. Here the contract is crisp: write parse(item) -> value-or-None once, and mapMaybe handles every item. dict.get is already such a function, so lookups that may miss compose directly.

\begin{lstlisting}[style=spec]
>>> legend = {'a': 1, 'b': 2}
>>> mapMaybe(legend.get, "cabbage")
[1, 2, 2, 1]
\end{lstlisting}

The Haskell connection: Data.Maybe's mapMaybe, with None standing in for Nothing -- the "parse, don't validate" discipline: turn questionable data into clean data at the boundary, and everything downstream stops worrying.


Remember it as: map with a parser, drop the Nothings, one pass.


Roll your own --- the definition:

\begin{lstlisting}[style=code]
mapMaybe = lambda f, xs: [y for y in (f(x) for x in xs) if y is not None]
\end{lstlisting}

\section{foldl}
\noindent\textit{[folds]} --- THE left fold, and the prelude's most important definition. It is exactly this loop: acc = init; for x in xs: acc = f(acc, x); return acc. The combiner takes (ACCUMULATOR, element) -- accumulator first. Argument order of foldl itself: (f, xs, init), with init optional -- omitted, the first element seeds the accumulator (and an empty list is then an error).
\begin{lstlisting}[style=spec]
>>> foldl(lambda acc, x: acc * 10 + x, [1, 2, 3], 0)
123
>>> foldl(max, [3, 9, 2])               # no init: 3 seeds it
9
>>> foldl(add, [], 0)                   # empty + init: the init comes back
0
>>> foldl(lambda acc, w: acc + len(w), ["ab", "c"], 0)
3
\end{lstlisting}

How it works: foldl is precisely this loop, given a name and a contract:

\begin{lstlisting}[style=spec]
>>> def foldl_spelled_out(f, xs, init):
...     acc = init
...     for x in xs:
...         acc = f(acc, x)
...     return acc
>>> foldl_spelled_out(add, [1, 2, 3], 0) == foldl(add, [1, 2, 3], 0)
True
\end{lstlisting}

Two disciplines hide in the signature. The combiner takes (ACCUMULATOR, element) -- accumulator first, always; reversing them is the classic fold bug. And the optional init is guarded by NOTHING, so the fold can tell "no seed -- use the first element" from "the seed is None".


Why it earns its place: the fold is the universal list consumer. sum, max, reverse, "build a dict", "run a state machine over events" -- each is one choice of combiner and seed. When you can say "carry THIS state, update it THIS way per element", the loop is already written.

\begin{lstlisting}[style=spec]
>>> foldl(lambda acc, x: [x] + acc, [1, 2, 3], [])       # reverse is a fold
[3, 2, 1]
>>> foldl(lambda acc, w: acc + len(w), ["ab", "cde"], 0)
5
\end{lstlisting}

Remember it as: a fold is a for-loop whose state has been given a contract.


Roll your own: seed the accumulator -- init if one was given, else the first element (the NOTHING guard is what lets None be a legitimate seed) -- then a single for-loop folding acc = f(acc, x). The entire contract is which side the accumulator sits on.

\begin{lstlisting}[style=code]
def foldl(f, xs, init=NOTHING):  # THE left fold; Python buried its own in functools as reduce
    it = iter(xs)
    if init is NOTHING:
        try:
            acc = next(it)
        except StopIteration:
            raise TypeError("fold of empty sequence with no initial value")
    else:
        acc = init
    for x in it:
        acc = f(acc, x)
    return acc
\end{lstlisting}

\section{foldr}
\noindent\textit{[folds]} --- the RIGHT fold: nests from the right, f takes (element, ACCUMULATOR) -- element first, mirror of foldl. Implemented as foldl of the flipped function over the reversed list: iterative and stack-safe, an identity that only holds because Python is strict. The parenthesisation example makes the direction visible.
\begin{lstlisting}[style=spec]
>>> foldr(lambda x, acc: '(' + x + '+' + acc + ')', "abc", "z")
'(a+(b+(c+z)))'
>>> foldl(lambda acc, x: '(' + acc + '+' + x + ')', "abc", "z")
'(((z+a)+b)+c)'
>>> foldr(lambda x, acc: [x] + acc, [1, 2, 3], [])
[1, 2, 3]
\end{lstlisting}

How it works: the mirror image -- the combiner takes (element, ACCUMULATOR) and the nesting grows from the RIGHT: f(a, f(b, f(c, z))). The prelude does not write a second recursion: foldr = foldl(flip(f), reversed(xs), init). Flip fixes the argument order, reversing fixes the traversal, and the existing engine runs backwards.


The honesty clause: that identity only holds because Python is STRICT. Haskell's foldr can process infinite lists -- laziness lets f decide whether the rest is ever needed; ours materialises and reverses the list first, so it cannot. Same name, same finite behaviour, different superpower.


Why it earns its place: building right-to-left. Consing onto a front is natural from the right -- section 13's from\_list IS a foldr of ListNode, and here it is, derived by hand:

\begin{lstlisting}[style=spec]
>>> foldr(lambda x, acc: [x] + acc, [1, 2, 3], [])       # rebuild, order kept
[1, 2, 3]
>>> to_list(foldr(ListNode, [1, 2, 3], None))            # from_list, spelled out
[1, 2, 3]
\end{lstlisting}

Remember it as: foldl eats left-to-right; foldr thinks right-to-left -- and in strict Python it is foldl wearing flip and a reversal.


Roll your own --- the definition:

\begin{lstlisting}[style=code]
foldr = lambda f, xs, init: foldl(flip(f), reversed(list(xs)), init)
\end{lstlisting}

\section{compose}
\noindent\textit{[folds]} --- fold any number of functions into one; the RIGHTMOST runs first, like mathematical composition f(g(x)).
\begin{lstlisting}[style=spec]
>>> compose(str.strip, str.upper)("  hi  ")
'HI'
>>> f = compose(succ, succ, succ)
>>> f(0)
3
\end{lstlisting}

How it works: a fold over the functions themselves -- neighbours merge into lambda x: f(g(x)), seeded with the identity, so compose() with no arguments is id. The RIGHTMOST function runs first, matching mathematical (f . g)(x) = f(g(x)).

\begin{lstlisting}[style=spec]
>>> compose()(42) == id(42)
True
>>> compose(len, str.strip)("  hi  ")     # strip first, then len
2
\end{lstlisting}

Why it earns its place: it turns a pipeline into a VALUE. A composed function can be named, mapped, stored in a rules table -- processing steps become data you can shuffle. The constraint is the one curry answers: every stage must be unary, which is why compose-heavy code leans on curry and partial to shape its stages.

\begin{lstlisting}[style=spec]
>>> normalize = compose(str.lower, str.strip)
>>> map_(normalize, ["  Hi ", " YO "])
['hi', 'yo']
\end{lstlisting}

Remember it as: foldr (.) id -- a pipeline you can hold in your hand.


Roll your own: capture the functions once, and have the returned function thread x through reversed(fns) in a plain loop. The loop (rather than nesting lambdas pairwise) keeps stack depth constant no matter how many stages you compose.

\begin{lstlisting}[style=code]
def compose(*fns): # foldr (.) id -- RIGHTMOST runs first; constant stack depth
    def composed(x):
        for f in reversed(fns):
            x = f(x)
        return x
    return composed
\end{lstlisting}

\section{subsequences}
\noindent\textit{[folds]} --- every subset of the elements (the powerset), built by a fold in which each element DOUBLES the accumulator: keep every existing subset, and every existing subset extended by the newcomer. 2\^{}n results -- a brute-force license for n up to about 20, and you should say that bound out loud before using it.
\begin{lstlisting}[style=spec]
>>> subsequences([1, 2])
[[], [1], [2], [1, 2]]
>>> subsequences("ab")
[[], ['a'], ['b'], ['a', 'b']]
>>> len(subsequences("abcd"))
16
\end{lstlisting}

How it works: watch the accumulator double. Start from [[]] -- one subset, the empty one. Each element x rewrites acc as acc + [s + [x] for s in acc]: every existing subset survives (x left out), and every existing subset reappears extended by x (x taken). n elements, n doublings, 2\^{}n subsets, in a stable order.

\begin{lstlisting}[style=spec]
>>> acc = [[]]
>>> acc = acc + [s + [1] for s in acc]; acc
[[], [1]]
>>> acc = acc + [s + [2] for s in acc]; acc
[[], [1], [2], [1, 2]]
\end{lstlisting}

Why it earns its place: it is a brute-force LICENSE. At n <= 20, checking every subset is a few million cheap operations -- honest, exhaustive, and unbeatable to debug. Say the bound out loud ("2\^{}15 is 32768 -- fine") and enumerate; cleverness can come later if the bound grows.

\begin{lstlisting}[style=spec]
>>> best = maxOn(sum, [s for s in subsequences([3, -1, 4, -2]) if sum(s) <= 6])
>>> sum(best)
6
\end{lstlisting}

Remember it as: the powerset fold -- each element doubles the world.


Roll your own --- the definition:

\begin{lstlisting}[style=code]
subsequences = lambda xs: foldl(lambda acc, x: acc + [s + [x] for s in acc], list(xs), [[]])
\end{lstlisting}

\section{accumulate}
\noindent\textit{[scans]} --- the generator behind the scans: yields the running combination of everything seen so far. With no function it sums; 'initial' seeds it. This is Python's itertools name with Haskell's scan semantics.
\begin{lstlisting}[style=spec]
>>> list(accumulate([1, 2, 3]))
[1, 3, 6]
>>> list(accumulate([1, 2], initial=10))
[10, 11, 13]
>>> list(accumulate([3, 1, 2], min))
[3, 1, 1]
\end{lstlisting}

Roll your own: the same skeleton as foldl, but a GENERATOR -- yield the accumulator once before the loop, then yield again after every combine. The early bare return on an empty, seedless input is what makes an empty scan come back empty instead of raising.

\begin{lstlisting}[style=code]
def accumulate(xs, f=None, initial=NOTHING):  # scanl / scanl1
    if f is None:
        f = lambda a, b: a + b
    it = iter(xs)
    if initial is NOTHING:
        try:
            acc = next(it)
        except StopIteration:
            return
    else:
        acc = initial
    yield acc
    for x in it:
        acc = f(acc, x)
        yield acc
\end{lstlisting}

\section{scanl}
\noindent\textit{[scans]} --- a fold that KEEPS ITS HISTORY: returns every intermediate accumulator, starting with the seed, so the result has len(xs)+1 elements and its last element equals the foldl. Argument order (f, z, xs) -- seed in the middle, Haskell's order, unlike foldl's trailing init. Prefix sums are the canonical scan.
\begin{lstlisting}[style=spec]
>>> scanl(add, 0, [3, 1, 4])
[0, 3, 4, 8]
>>> last(scanl(add, 0, [1, 2, 3])) == foldl(add, [1, 2, 3], 0)
True
\end{lstlisting}

The DP secret: many celebrated one-pass algorithms are scans in a trench coat, because a scan IS a dynamic-programming table -- entry i holds the answer for the prefix ending at i. Kadane's famous maximum-subarray "trick" is one scanl1 and a max:

\begin{lstlisting}[style=spec]
>>> xs = [-2, 1, -3, 4, -1, 2, 1, -5, 4]
>>> max(scanl1(lambda best, x: max(x, best + x), xs))
6
\end{lstlisting}

Prefix sums, the other superpower: scan once with (+) and every range sum becomes two reads, pref[j] - pref[i] -- O(n) window questions collapse to O(1).

\begin{lstlisting}[style=spec]
>>> pref = scanl(add, 0, [3, 1, 4, 1, 5])
>>> pref[4] - pref[1]                     # sum of elements 1..3
6
\end{lstlisting}

Remember it as: when you need the state at EVERY step, one scan replaces n folds.


Roll your own --- the definition:

\begin{lstlisting}[style=code]
scanl = lambda f, z, xs: list(accumulate(xs, f, initial=z))
\end{lstlisting}

\section{unfoldr}
\noindent\textit{[streams]} --- iterate's FINITE twin, and the cure for threading (state, acc) tuples through loops. You supply f(seed) returning either (value, seed') to emit a value and continue, or None to stop; unfoldr collects the values. Anything that "peels" a structure -- digits, chunks, parent chains -- is an unfoldr.
\begin{lstlisting}[style=spec]
>>> unfoldr(lambda m: None if m == 0 else (m % 10, m // 10), 407)
[7, 0, 4]
>>> unfoldr(lambda n: None if n == 0 else (n, n - 1), 5)
[5, 4, 3, 2, 1]
>>> unfoldr(const(None), "seed")        # stop immediately: empty
[]
\end{lstlisting}

How it works: the step function carries a two-outcome contract -- f(seed) returns (value, next\_seed) to emit and continue, or None to stop -- and unfoldr collects what was emitted. It is the fold REVERSED: a fold consumes a list into a summary; an unfold grows a list out of a seed. (Category theory says anamorphism; interviewers say "generate the sequence".)


Why it earns its place: it retires the (state, accumulator) tuple you would otherwise thread through a while loop. Anything that PEELS -- digits off a number, chunks off a list, parents up a tree, a Collatz trajectory -- states its entire logic in the step function.

\begin{lstlisting}[style=spec]
>>> unfoldr(lambda n: None if n == 1 else (n, n // 2 if even(n) else 3 * n + 1), 6)
[6, 3, 10, 5, 16, 8, 4, 2]
\end{lstlisting}

Versus iterate: iterate is the infinite cousin -- an unfold whose step never says None, consumed with take. Reach for iterate when the stream is conceptually endless, unfoldr when the seed itself knows where the end is.


Remember it as: fold tears down; unfold builds up; the seed carries the future.


More unfolds, same two-outcome contract -- peeling chunks off a list, bits off a number, walking a parent chain to its root:

\begin{lstlisting}[style=spec]
>>> unfoldr(lambda xs: None if xs == [] else (xs[:2], xs[2:]), [1, 2, 3, 4, 5])
[[1, 2], [3, 4], [5]]
>>> unfoldr(lambda m: None if m == 0 else (m % 2, m // 2), 13)       # bits, LSB first
[1, 0, 1, 1]
>>> parents = {'c': 'b', 'b': 'a', 'a': None}
>>> unfoldr(lambda n: None if n is None else (n, parents[n]), 'c')   # ancestry
['c', 'b', 'a']
\end{lstlisting}

And the seed can carry several facts at once -- here (current, next, remaining) generates Fibonacci with its own countdown built in:

\begin{lstlisting}[style=spec]
>>> unfoldr(lambda s: None if s[2] == 0 else (s[0], (s[1], s[0] + s[1], s[2] - 1)), (0, 1, 8))
[0, 1, 1, 2, 3, 5, 8, 13]
\end{lstlisting}

Roll your own: prime the pump with step = f(seed), then loop while step is not None -- append the value, feed the new seed back to f. The two-outcome contract lives entirely in f; unfoldr is just the loop that honours it, plus the list that collects what was emitted.

\begin{lstlisting}[style=code]
def unfoldr(f, seed):  # Data.List unfoldr: f(seed) -> None | (value, seed')
    out = []  # iterate's finite twin: grows a list until f says stop --
    step = f(seed)  # no more threading (state, acc) tuples through until
    while step is not None:
        val, seed = step
        out.append(val)
        step = f(seed)
    return out
\end{lstlisting}

\section{span}
\noindent\textit{[slicing \& spans]} --- split at the first failure: (longest satisfying prefix, the rest), computed in ONE pass so it is safe on one-shot iterators (the old two-scan version silently broke on generators -- this is the fix). A lexer in one call.
\begin{lstlisting}[style=spec]
>>> span(lambda x: x < 3, [1, 2, 3, 4, 1])
([1, 2], [3, 4, 1])
>>> span(str.isdigit, "123abc")
(['1', '2', '3'], ['a', 'b', 'c'])
\end{lstlisting}

Roll your own: materialise xs once (that is the one-pass guarantee), locate the first failing index with next() over enumerate -- defaulting to len(xs) when the predicate never fails -- and slice twice. The previous two-scan version ran takewhile and dropwhile separately and silently broke on one-shot generators; this is the repair.

\begin{lstlisting}[style=code]
def span(p, xs):  # split at first failure, ONE pass; safe on one-shot iters
    xs = list(xs)
    i = next((i for i, x in enumerate(xs) if not p(x)), len(xs))
    return (xs[:i], xs[i:])
\end{lstlisting}

\section{cmp\_to\_key}
\noindent\textit{[sorting \& searching]} --- adapts an old-style comparator (returning negative, zero or positive) into the key= protocol. Reach for it only when a rule genuinely compares two elements against each other and no per- element key exists.
\begin{lstlisting}[style=spec]
>>> sorted(["bb", "a"], key=cmp_to_key(lambda a, b: len(a) - len(b)))
['a', 'bb']
\end{lstlisting}

Roll your own: a tiny wrapper class implementing ONE dunder -- \_\_lt\_\_ defers to cmp(a, b) < 0. Python's sort only ever asks "is a less than b?", so that single method is enough to smuggle any comparator through the key= interface.

\begin{lstlisting}[style=code]
def cmp_to_key(cmp):  # resurrect Python 2's cmp for sortBy
    class K:
        def __init__(self, x): self.x = x
        def __lt__(self, other): return cmp(self.x, other.x) < 0
    return K
\end{lstlisting}

\section{merge}
\noindent\textit{[sorting \& searching]} --- join two ALREADY-SORTED lists into one sorted list, stably, in O(n+m): the two-pointer loop at the heart of merge sort.
\begin{lstlisting}[style=spec]
>>> merge([1, 4, 6], [2, 3, 7])
[1, 2, 3, 4, 6, 7]
>>> merge([], [1, 2])
[1, 2]
\end{lstlisting}

Roll your own: two cursors, always copying the smaller head forward; the <= (rather than <) is what makes the merge STABLE, preferring the left list on ties. When either side runs dry, glue both tails on -- one of them is empty, so the + is harmless.

\begin{lstlisting}[style=code]
def merge(a, b):  # merge two sorted lists, stable; the heart of merge sort
    out, i, j = [], 0, 0  # (section 13's merge_lists is this same loop on ListNodes)
    while i < len(a) and j < len(b):
        if a[i] <= b[j]:
            out.append(a[i]); i += 1
        else:
            out.append(b[j]); j += 1
    return out + a[i:] + b[j:]
\end{lstlisting}

\section{bisect\_left}
\noindent\textit{[sorting \& searching]} --- binary search over a SORTED list: the first index whose element is >= x; equivalently, where x would insert to the left of its equals. If a[i] != x there, x is absent -- that check is binary search.
\begin{lstlisting}[style=spec]
>>> bisect_left([10, 20, 20, 30], 20)
1
>>> bisect_left([10, 20, 30], 25)       # insertion point for a missing value
2
\end{lstlisting}

How it works: on a sorted list the predicate "element < x" reads True for a prefix and False for the rest -- TTTTFFFF. Binary search does not hunt for x; it finds that BOUNDARY, halving the interval per probe. bisect\_left returns the first index >= x, bisect\_right the first > x, and from those two facts three idioms fall out.

\begin{lstlisting}[style=spec]
>>> xs = [10, 20, 20, 30]
>>> bisect_left(xs, 20), bisect_right(xs, 20)     # the run of 20s occupies [1, 3)
(1, 3)
\end{lstlisting}

Idiom one, membership: x is present iff the slot bisect\_left names actually holds it. Idiom two, insertion point: that same index is where x belongs to keep the list sorted. Idiom three, counting: right minus left counts occurrences -- and with two different probes, the values inside any range.

\begin{lstlisting}[style=spec]
>>> i = bisect_left(xs, 25)
>>> i, (i < len(xs) and xs[i] == 25)              # 25 would insert at 3; absent
(3, False)
\end{lstlisting}

Why it earns its place: every threshold ladder is secretly a sorted list. Grading bands, tax brackets, rate tiers -- one bisect replaces the if/elif staircase:

\begin{lstlisting}[style=spec]
>>> grade = lambda score: "FDCBA"[bisect_right([60, 70, 80, 90], score)]
>>> grade(59), grade(60), grade(95)
('F', 'D', 'A')
\end{lstlisting}

Remember it as: bisect finds the boundary in a sea of Trues and Falses; everything else is reading that boundary three ways.


Roll your own: the lo/hi halving loop over an invariant -- everything left of lo is < x, everything from hi rightwards is >= x. One comparison, a[mid] < x, decides which half survives; when lo meets hi, that meeting point IS the boundary. bisect\_right is the identical loop with <= in place of < -- one character moves the boundary to the far side of the equals.

\begin{lstlisting}[style=code]
def bisect_left(a, x):  # first index where a[i] >= x
    lo, hi = 0, len(a)
    while lo < hi:
        mid = (lo + hi) // 2
        if a[mid] < x:
            lo = mid + 1
        else:
            hi = mid
    return lo
\end{lstlisting}

\section{defaultdict}
\noindent\textit{[containers]} --- a dict whose missing keys create themselves via a factory: touch d[k] and factory() is installed there. This "autovivification" makes grouping and counting one-liners. Gotcha: READING a missing key also plants it -- use 'in' or .get when you only want to look.
\begin{lstlisting}[style=spec]
>>> d = defaultdict(list)
>>> d['x'].append(1)
>>> d
{'x': [1]}
>>> d2 = defaultdict(int)
>>> d2['ghost']                         # a mere read...
0
>>> 'ghost' in d2                       # ...planted the key
True
\end{lstlisting}

How it works: one dunder does everything. When d[k] misses, Python calls d.\_\_missing\_\_(k); this subclass responds by calling the stored factory, installing the result, and returning it. The factory runs only on a miss -- and it takes NO arguments, which is why the leaf recipes in Tree/ITree are zero-argument lambdas.


Why it earns its place: it deletes the check-then-create dance from every grouping and counting loop. The write path becomes a bare append or +=, and the structure assembles itself on touch -- autovivification, which section 12's trie tools then apply recursively.


The discipline it demands: a mere READ also plants the key, so split your paths -- write with d[k], but ask questions with 'in' or .get. The prelude's bag operations read Counters with .get for exactly this reason.

\begin{lstlisting}[style=spec]
>>> d = defaultdict(int)
>>> d['hits'] += 1                      # write path: indexing is the point
>>> d.get('misses', 0), 'misses' in d   # read path: no ghost key planted
(0, False)
\end{lstlisting}

Remember it as: index to build, .get to ask.


Roll your own: subclass dict and implement one dunder -- \_\_missing\_\_, which Python calls when [k] fails. Call the factory, install the result, hand it back; iteration, repr and everything else is inherited. Note the factory takes no arguments, which is why leaf recipes are zero-argument lambdas.

\begin{lstlisting}[style=code]
class defaultdict(dict):
    def __init__(self, factory=None, *args, **kw):
        super().__init__(*args, **kw)
        self.factory = factory
    def __missing__(self, key):
        if self.factory is None:
            raise KeyError(key)
        self[key] = self.factory()
        return self[key]
\end{lstlisting}

\section{Counter}
\noindent\textit{[containers]} --- count occurrences: a defaultdict(int) fold over the input. Missing keys read as 0 (mind the autovivification gotcha above); rank the .items() with sortOn.
\begin{lstlisting}[style=spec]
>>> c = Counter("abracadabra")
>>> c['a'], c['z']
(5, 0)
>>> sortOn(lambda kv: -kv[1], list(Counter("aab").items()))
[('a', 2), ('b', 1)]
\end{lstlisting}

Roll your own: a defaultdict(int) plus one fold -- d[x] += 1 per element. The += on a missing key is autovivification doing all the work.

\begin{lstlisting}[style=code]
def Counter(xs):  # count-by-value; a defaultdict(int) fold
    d = defaultdict(int)
    for x in xs:
        d[x] += 1
    return d
\end{lstlisting}

\section{insertWith}
\noindent\textit{[containers]} --- Data.Map's insertWith f k v m: insert (k, v), combining with f(new, old) on collision -- Haskell's argument order, new value first. PERSISTENT, also like Haskell: it returns a fresh dict and leaves the argument untouched.
\begin{lstlisting}[style=spec]
>>> insertWith(add, 'x', 4, {'x': 1})
{'x': 5}
>>> d = {'y': 2}
>>> insertWith(add, 'x', 4, d)
{'y': 2, 'x': 4}
>>> d                                    # the original survives
{'y': 2}
\end{lstlisting}

Roll your own: copy the dict first -- dict(d) IS the persistence -- then one conditional: f(v, out[k]) on collision (new value first, Haskell's order) or a plain insert. Compare fromListWith, which inlines the same rule but mutates, because there the dict is its own private accumulator.

\begin{lstlisting}[style=code]
def insertWith(f, k, v, d):  # Data.Map insertWith f k v m: PERSISTENT -- returns a new
    out = dict(d)  # dict untouched; collision -> f(new, old), Haskell order
    out[k] = f(v, out[k]) if k in out else v  # order (new value first)
    return out
\end{lstlisting}

\section{fromListWith}
\noindent\textit{[containers]} --- THE dict-building fold: pour (key, value) pairs into a dict, combining collisions with f(new, old) -- Haskell's order, like insertWith. Keys keep first-seen order.
\begin{lstlisting}[style=spec]
>>> fromListWith(add, [(c, 1) for c in "aab"])
{'a': 2, 'b': 1}
>>> fromListWith(lambda new, old: old + new, [('a', [1]), ('b', [2]), ('a', [3])])
{'a': [1, 3], 'b': [2]}
\end{lstlisting}

The order gotcha, faithfully Haskell's: because f receives (new, old), grouping with a bare concat REVERSES each group -- the classic Data.Map surprise. Combine as old + new to keep arrival order, or use plain addition where order cannot matter (counts).

\begin{lstlisting}[style=spec]
>>> fromListWith(add, [(len(w), [w]) for w in ["hi", "ox", "sun"]])
{2: ['ox', 'hi'], 3: ['sun']}
>>> fromListWith(lambda new, old: old + new, [(len(w), [w]) for w in ["hi", "ox", "sun"]])
{2: ['hi', 'ox'], 3: ['sun']}
\end{lstlisting}

Why it earns its place -- one function, three monoids: choosing f is choosing what the dict MEANS. Addition over (k, 1) pairs counts; old + new over (k, [v]) pairs groups; max merges bags. That is the whole family: insertWith is one collision, fromListWith a list of them, unionWith two dicts' worth (where f receives (a's value, b's value)).


The recognition cue: whenever a loop builds a dict with an if-else on key existence, you are hand-rolling this fold. Name the pairs, name the monoid, and the loop disappears.


Remember it as: grouping, counting, and merging are one fold -- whose combiner hears the NEW value first.


Roll your own: an empty dict and one loop applying the collision rule per pair -- d[k] = f(v, d[k]) if the key exists, else v. It is insertWith inlined and IN-PLACE: mutation is fine here because the dict being built is the fold's own accumulator, invisible until returned.

\begin{lstlisting}[style=code]
def fromListWith(f, pairs):  # Data.Map fromListWith: THE dict-building fold; f(new, old)
    d = {}  # like insertWith -- so grouping with concat REVERSES each
    for k, v in pairs:  # group (the classic Haskell gotcha); use add for counts
        d[k] = f(v, d[k]) if k in d else v
    return d  # Counter == fromListWith(add) over (x, 1);
  # grouping == fromListWith(++) over (k, [v])
\end{lstlisting}

\section{unionWith}
\noindent\textit{[containers]} --- merge two dicts, combining SHARED keys with f; a's keys keep their order, b's new keys follow. Choosing f is choosing a monoid: max gives bag union, + gives a summing merge.
\begin{lstlisting}[style=spec]
>>> unionWith(max, {'a': 1, 'b': 5}, {'b': 2, 'c': 3})
{'a': 1, 'b': 5, 'c': 3}
>>> unionWith(add, {'a': 1, 'b': 2}, {'b': 3, 'c': 4})
{'a': 1, 'b': 5, 'c': 4}
\end{lstlisting}

Roll your own: start from a copy of a, then pour b's items through the collision rule f(a's value, b's value). The asymmetry is deliberate: a's keys keep their positions, and only keys new to b append at the end.

\begin{lstlisting}[style=code]
def unionWith(f, a, b):  # Data.Map unionWith: merge two dicts, f on shared keys
    out = dict(a)  # f(a's value, b's value) -- Haskell's left/right order
    for k, v in b.items():
        out[k] = f(out[k], v) if k in out else v
    return out  # one combinator, many merges:
  # bag_union == unionWith(max), merge-sum == unionWith(add)
\end{lstlisting}

\section{groupBy}
\noindent\textit{[containers]} --- runs by an equivalence predicate; each newcomer is compared with eq against the run's FIRST element (span-based, exactly Haskell's groupBy), not against its neighbour -- a subtle and deliberate fidelity.
\begin{lstlisting}[style=spec]
>>> groupBy(lambda a, b: a == b, [1, 1, 2])
[[1, 1], [2]]
>>> groupBy(lambda a, b: b - a <= 1, [1, 2, 3, 10, 11])
[[1, 2], [3], [10, 11]]
\end{lstlisting}

The word that matters is CONSECUTIVE: a new run starts whenever eq against the run's first element fails, so 'abab' yields four runs. Not a weakness -- a different question: runs, streaks and compression care about position, and these are the only tools asking about it.

\begin{lstlisting}[style=spec]
>>> [(run[0], len(run)) for run in group("aaabbbbcc")]     # run lengths
[('a', 3), ('b', 4), ('c', 2)]
\end{lstlisting}

For CATEGORIES -- position irrelevant -- either sort first, so equals become adjacent, or use fromListWith, which never cared about position at all. Choosing between them is stating what your key means.

\begin{lstlisting}[style=spec]
>>> [(run[0], len(run)) for run in group(sorted("abab"))]
[('a', 2), ('b', 2)]
\end{lstlisting}

Remember it as: group answers "how long are the runs?"; fromListWith answers "what are the piles?".


Roll your own: one pass in which each element either joins the current run or starts a fresh one. The test is eq(out[-1][0], x) -- against the run's FIRST element, not its last -- which is exactly how Haskell's span-based groupBy behaves, and the detail that makes non-transitive predicates (like "within 1 of") agree with it.

\begin{lstlisting}[style=code]
def groupBy(eq, xs):  # Data.List groupBy: runs of CONSECUTIVE elements; each new
    out = []  # element is compared via eq against the run's FIRST element,
    for x in xs:  # exactly like Haskell (span-based) -- not its neighbor
        if out and eq(out[-1][0], x):
            out[-1].append(x)
        else:
            out.append([x])
    return out
\end{lstlisting}

\section{ITree}
\noindent\textit{[containers]} --- the INFINITE tree: every branch is another ITree, so any depth autovivifies on touch. This is the trie: index letter by letter and the branch builds itself.
\begin{lstlisting}[style=spec]
>>> t = ITree()
>>> t['x']['y']['z'] = 9
>>> t
{'x': {'y': {'z': 9}}}
\end{lstlisting}

How it works: the definition is a knot -- ITree = lambda: defaultdict(ITree). Every missing key manufactures another ITree, which will do the same, without end. Indexing therefore GROWS the tree: touching t['c']['a']['t'] conjures the whole branch. A trie in one line, built from nothing but \_\_missing\_\_.

\begin{lstlisting}[style=spec]
>>> t = ITree()
>>> for w in ["cat", "car"]: setpath(t, list(w) + ['$'], w)
>>> sorted([leaf for p, leaf in paths(t) if p[-1] == '$'])
['car', 'cat']
\end{lstlisting}

Why it earns its place: hierarchies -- tries, directory trees, nested groupings -- stop needing setup code. Writing is indexing (or setpath, which also plants a leaf); harvesting is paths. The one edged rule: reads autovivify too, so QUESTIONS go through getpath, never bare indexing -- getpath's entry demonstrates the ghost branch this avoids.


Versus Tree(depth, leaf): Tree is the finite cousin -- fixed depth, with a real leaf factory (list, int) at the bottom instead of more tree. Tree(1, list) is the grouping dict; ITree is for depth you cannot know in advance, like words.


Remember it as: a dict that dreams up more of itself on demand.


Roll your own: ITree = lambda: defaultdict(ITree) mentions itself -- legal, because the lambda body only runs on a MISS, by which time the name is bound. Tree bottoms the same recursion out at a fixed depth by handing the last level a real leaf factory instead of more tree.

\begin{lstlisting}[style=code]
ITree = lambda: defaultdict(ITree)
\end{lstlisting}

\section{paths}
\noindent\textit{[containers]} --- fold a whole Tree/ITree into a flat list of (keypath, leaf) pairs, in insertion order. The read- everything companion to setpath's write-one-path.
\begin{lstlisting}[style=spec]
>>> t = ITree()
>>> setpath(t, ['a', 'b'], 1); setpath(t, ['a', 'c'], 2)
>>> paths(t)
[(['a', 'b'], 1), (['a', 'c'], 2)]
\end{lstlisting}

How it works: structural recursion with two arms, exactly the case-arms-as-fold story. A non-dict (or empty node) is a LEAF: report the single row ([], leaf). A dict is a BRANCH: recurse into every subtree and prepend the branching key to each keypath that comes back. The whole tree flattens into (keypath, leaf) rows, insertion-ordered, losslessly.


Why it earns its place: building a trie is the easy half; the value is HARVESTING it. autocomplete filters rows under a prefix, deepest\_directory takes maxOn of keypath length, leaves reads the second column. One cata, many harvests -- you never write tree-walking code again.

\begin{lstlisting}[style=spec]
>>> t = ITree()
>>> setpath(t, ['a', 'b'], 1); setpath(t, ['c'], 2)
>>> paths(t)
[(['a', 'b'], 1), (['c'], 2)]
>>> maxOn(lambda row: len(row[0]), paths(t))
(['a', 'b'], 1)
\end{lstlisting}

Remember it as: paths turns a tree back into rows; setpath turns rows into a tree.


Roll your own: recursion on one question -- is this still a dict? A leaf reports a single row with an empty path; a branch recurses into every child and conses its key onto every path that comes back. Depth equals tree height, comfortably inside Python's limit for sane tries.

\begin{lstlisting}[style=code]
def paths(t):  # cata for Tree/ITree: [(keypath, leaf)]
  # recursion depth = tree height (word length / len(nums)) -- well under the ~1000 limit
    if not isinstance(t, dict) or not t:  # leaf = non-dict value, or empty node
        return [([], t)]
    return [([k] + p, leaf) for k, sub in t.items() for p, leaf in paths(sub)]
\end{lstlisting}

\section{setpath}
\noindent\textit{[containers]} --- write a leaf value at the end of a key path, autovivifying the branch on the way. Caution: writing at an interior key CLOBBERS the subtree below it.
\begin{lstlisting}[style=spec]
>>> t = ITree()
>>> setpath(t, ['a', 'b'], 7)
>>> t
{'a': {'b': 7}}
>>> setpath(t, ['a'], 0)                # clobbers {'b': 7}
>>> t
{'a': 0}
\end{lstlisting}

Roll your own: walk all but the last key by PLAIN INDEXING -- on an ITree that autovivifies the branch as you go -- then assign at the final key, the single write. getpath is the mirror image with the safety catch on: isinstance and 'in' checks instead of indexing, so reading plants nothing.

\begin{lstlisting}[style=code]
def setpath(t, ks, v):  # write leaf v at key path ks (autovivifies interior)
    for k in ks[:-1]:
        t = t[k]
    t[ks[-1]] = v  # NB: clobbers any subtree already at ks
\end{lstlisting}

\section{deque}
\noindent\textit{[containers]} --- a double-ended queue built from two stacks, all four end operations amortised O(1) (elements migrate between stacks only when one runs dry). This is the BFS frontier; a plain list's pop(0) is O(n).
\begin{lstlisting}[style=spec]
>>> d = deque([1, 2, 3])
>>> d.appendleft(0); d.append(4)
>>> d.popleft(), d.pop()
(0, 4)
>>> len(d), list(d)
(3, [1, 2, 3])
\end{lstlisting}

How it works: two plain lists, nose to nose -- \_out holds the front (reversed), \_in holds the back. append pushes onto \_in, appendleft onto \_out, both O(1). The trick is popping an EMPTY side: the other list is reversed across in one O(n) migration, and popping is O(1) again.


Why that still counts as O(1): amortization. Each element migrates at most once between being pushed and popped, so n operations cost O(n) in total -- constant per operation on average, which is what the workload's big-O cares about. A bare Python list cannot say this: pop(0) shifts every element, every single time.

\begin{lstlisting}[style=spec]
>>> q = deque()
>>> for x in [1, 2, 3]: q.append(x)
>>> q.popleft(), q.popleft()            # first pop pays the migration; the next is free
(1, 2)
\end{lstlisting}

Why it earns its place: BFS. The frontier must yield its OLDEST member (popleft) while newcomers join at the back (append) -- levelorder and grid\_path are this class doing its one job.


Remember it as: two stacks facing each other make a queue; the reversal bill is paid once per element.


Roll your own: two plain lists nose to nose, each end appending and popping its own list at O(1). All the cleverness is in the pop guard: when your side is empty, reverse the OTHER list across in one go. Each element makes that crossing at most once between entering and leaving, so the occasional O(n) reversal amortises to O(1) per operation.

\begin{lstlisting}[style=code]
class deque:  # two-stack; all four ends amortized O(1)
    def __init__(self, xs=()):
        self._in, self._out = [], list(reversed(list(xs)))
    def append(self, x): self._in.append(x)
    def appendleft(self, x): self._out.append(x)
    def pop(self):
        if not self._in:
            self._in, self._out = list(reversed(self._out)), []
        return self._in.pop()
    def popleft(self):
        if not self._out:
            self._out, self._in = list(reversed(self._in)), []
        return self._out.pop()
    def __len__(self): return len(self._in) + len(self._out)
    def __iter__(self): return iter(self._out[::-1] + self._in)
\end{lstlisting}

\section{dsu}
\noindent\textit{[containers]} --- disjoint-set union (union-find) over n elements, with path halving. find gives a set's representative; union merges and returns False when the two were ALREADY connected -- which is exactly a cycle detector (Kruskal's test).
\begin{lstlisting}[style=spec]
>>> find, union = dsu(4)
>>> union(0, 1), union(1, 2)
(True, True)
>>> union(0, 2)                         # already connected: a cycle
False
>>> find(0) == find(2), find(0) == find(3)
(True, False)
\end{lstlisting}

How it works: parent[] links every element toward its set's ROOT, and find follows the chain up. The subtle line is parent[x] = parent[parent[x]] -- path HALVING: every find re-points each visited node at its grandparent, so chains shrink as a side effect of being walked. That one line makes operations effectively constant time (inverse Ackermann, for the curious).


Why it earns its place: dynamic connectivity. "Are these two in the same group yet?" under a stream of merges is awkward for every other structure and trivial here. And union returning False -- already connected -- IS cycle detection: Kruskal's spanning tree is built by skipping every False.

\begin{lstlisting}[style=spec]
>>> find, union = dsu(5)
>>> [union(a, b) for a, b in [(0, 1), (1, 2), (3, 4), (0, 2)]]
[True, True, True, False]
>>> len({find(i) for i in range(5)})   # two components remain
2
\end{lstlisting}

Remember it as: find walks to the root, halving as it goes; a False union is a cycle caught.


Roll your own: parent[] as the forest and two closures instead of a class. find's loop carries the entire trick in one line -- parent[x] = parent[parent[x]] re-points every visited node at its grandparent (path halving), flattening chains as a side effect of walking them. union is find twice plus one root re-point, answering False when the roots already agree.

\begin{lstlisting}[style=code]
def dsu(n):  # union-find, path halving; union -> False if already joined
    parent = list(range(n))
    def find(x):
        while parent[x] != x:
            parent[x] = parent[parent[x]]
            x = parent[x]
        return x
    def union(a, b):
        ra, rb = find(a), find(b)
        if ra == rb:
            return False
        parent[ra] = rb
        return True
    return find, union
\end{lstlisting}

\section{heappush}
\noindent\textit{[containers]} --- push onto a min-heap kept in a plain list (sift-up). The smallest element is always h[0]. For a MAX-heap, push negated values.
\begin{lstlisting}[style=spec]
>>> h = []
>>> for x in [5, 1, 8]: heappush(h, x)
>>> h[0]
1
\end{lstlisting}

How it works: the flat list IS a binary tree -- h[i]'s children live at 2i+1 and 2i+2 -- kept under one invariant: every parent <= its children. heappush appends at the first free leaf and sifts UP, swapping with its parent while smaller; heappop moves the last leaf to the root and sifts DOWN toward the smaller child. Each walks one root-to-leaf path: O(log n).


Why the invariant is enough: nothing is fully sorted -- siblings sit in any order -- yet h[0] is ALWAYS the minimum, because the root beats its children, who beat theirs. You pay log n per operation for a permanently readable minimum; draining the heap is heapsort; pushing negatives turns min into max.

\begin{lstlisting}[style=spec]
>>> h = []
>>> for x in [7, 2, 9, 4]: heappush(h, x)
>>> h[0]
2
>>> [heappop(h) for _ in range(4)]
[2, 4, 7, 9]
\end{lstlisting}

Why it earns its place: "repeatedly take the best next thing" -- Dijkstra's frontier, running\_median's two halves, any streaming top-k. When next-best is needed many times while the data keeps changing, the heap beats re-sorting every time.


Remember it as: a flat list, a parent-beats-children promise, and log-n repairs.


Roll your own: append at the first free leaf, then sift UP -- while the parent at (i - 1) // 2 is larger, swap and climb. The index arithmetic IS the tree; watch the array after four pushes:

\begin{lstlisting}[style=spec]
>>> h = []
>>> for x in [7, 2, 9, 4]: heappush(h, x)
>>> h                       # root 2; children 4 and 9; 7 pushed down under 4
[2, 4, 9, 7]
\end{lstlisting}
\begin{lstlisting}[style=code]
def heappush(h, x):  # min-heap on a plain list: sift-up
    h.append(x)
    i = len(h) - 1
    while i and h[(i - 1) // 2] > h[i]:
        h[(i - 1) // 2], h[i] = h[i], h[(i - 1) // 2]
        i = (i - 1) // 2
\end{lstlisting}

\section{heappop}
\noindent\textit{[containers]} --- remove and return the minimum (sift-down). Popping repeatedly drains in sorted order -- that is heapsort.
\begin{lstlisting}[style=spec]
>>> h = []
>>> for x in [5, 1, 8]: heappush(h, x)
>>> [heappop(h) for _ in range(3)]
[1, 5, 8]
>>> g = []
>>> for x in [3, 9, 5]: heappush(g, -x)
>>> -heappop(g)                         # max-heap via negation
9
\end{lstlisting}

Roll your own: swap the root with the last leaf, pop it off the end, then sift DOWN -- repeatedly swap with the SMALLER of the two children (choosing the smaller is the classic bug site) until neither child is smaller. One root-to-leaf path each way: O(log n).

\begin{lstlisting}[style=code]
def heappop(h):  # min-heap on a plain list: sift-down
    h[0], h[-1] = h[-1], h[0]
    top = h.pop()
    i, n = 0, len(h)
    while True:
        s, l, r = i, 2*i + 1, 2*i + 2
        if l < n and h[l] < h[s]: s = l
        if r < n and h[r] < h[s]: s = r
        if s == i:
            return top
        h[i], h[s] = h[s], h[i]
        i = s
\end{lstlisting}

\section{tfold}
\noindent\textit{[nodes]} --- the tree catamorphism: recursively combines each node's value with the already-folded left and right results, z standing in for empty subtrees. EVERY structural tree computation is one choice of f and z -- the five traversal/measure functions below are all tfold one-liners.
\begin{lstlisting}[style=spec]
>>> t = TreeNode(2, TreeNode(1), TreeNode(3))
>>> tfold(lambda v, l, r: v + l + r, 0, t)      # sum
6
>>> tfold(lambda v, l, r: max(v, l, r), 0, t)   # max
3
\end{lstlisting}

How it works: the tree catamorphism -- structural recursion with a contract. Two arms, exactly like a case expression: the empty tree answers z; a node answers f(value, left\_result, right\_result), both results computed by the same recursion. Hand it (f, z) and any question about a tree becomes one line -- the five traversals and both measures in section 13 are all tfold with different arms.


Why it earns its place: it separates WALKING from DECIDING. The recursion is written once, correct forever; you only ever author the arms. Even tree TRANSFORMATION fits -- mirror a tree by swapping the child results while rebuilding:

\begin{lstlisting}[style=spec]
>>> t = TreeNode(2, TreeNode(1), TreeNode(3))
>>> mirror = tfold(lambda v, l, r: TreeNode(v, r, l), None, t)
>>> inorder(mirror)
[3, 2, 1]
\end{lstlisting}

The one caveat: recursion depth equals tree height -- fine for balanced and interview-sized trees; say a sentence out loud about adversarially skewed ones.


Remember it as: z is the empty-tree arm, f is the node arm -- the fold walks, you decide.


Roll your own: one line of structural recursion -- z when the tree is None, else f over the node's value and the two recursive results. Every traversal and measure in this section is that line with a different f plugged in.

\begin{lstlisting}[style=code]
def tfold(f, z, t):  # cata: f(val, l_acc, r_acc); depth = tree height
    return z if t is None else f(t.val, tfold(f, z, t.left), tfold(f, z, t.right))
\end{lstlisting}

\section{levelorder}
\noindent\textit{[nodes]} --- breadth-first values, level by level: the deque earning its keep. Compare preorder on the same tree to see DFS vs BFS.
\begin{lstlisting}[style=spec]
>>> t = TreeNode(1, TreeNode(2, TreeNode(4)), TreeNode(3))
>>> levelorder(t)
[1, 2, 3, 4]
>>> preorder(t)
[1, 2, 4, 3]
\end{lstlisting}

Roll your own: seed a deque with the root, then loop -- popleft, record the value, append the children that exist. FIFO order IS breadth-first; swap the deque for a list used as a stack and the identical loop turns into depth-first.

\begin{lstlisting}[style=code]
def levelorder(t):  # BFS: the deque earning its keep
    out, q = [], deque([t] if t else [])
    while len(q):
        n = q.popleft()
        out.append(n.val)
        if n.left:  q.append(n.left)
        if n.right: q.append(n.right)
    return out
\end{lstlisting}

\section{reverse\_list}
\noindent\textit{[nodes]} --- reverse in place with the three-pointer walk; prev is a fold accumulator wearing pointers.
\begin{lstlisting}[style=spec]
>>> to_list(reverse_list(from_list([1, 2, 3])))
[3, 2, 1]
\end{lstlisting}

Roll your own: the three-pointer walk compressed into one simultaneous assignment -- node.next, prev, node = prev, node, node.next. The right-hand side is evaluated in full before any assignment lands, which is what saves the pointer you are about to overwrite.

\begin{lstlisting}[style=code]
def reverse_list(node):  # three-pointer; prev is a fold accumulator
    prev = None
    while node:
        node.next, prev, node = prev, node, node.next
    return prev
\end{lstlisting}

\section{middle}
\noindent\textit{[nodes]} --- the middle node by fast/slow pointers (fast moves twice per step); on even lengths, the SECOND middle.
\begin{lstlisting}[style=spec]
>>> middle(from_list([1, 2, 3])).val
2
>>> middle(from_list([1, 2, 3, 4])).val
3
\end{lstlisting}

Roll your own: fast advances two hops for slow's one, so when fast falls off the end, slow stands at the middle. The loop condition (fast and fast.next) is precisely what selects the SECOND middle on even lengths.

\begin{lstlisting}[style=code]
def middle(node):  # fast/slow: second middle for even lengths
    slow = fast = node
    while fast and fast.next:
        slow, fast = slow.next, fast.next.next
    return slow
\end{lstlisting}

\section{has\_cycle}
\noindent\textit{[nodes]} --- Floyd's tortoise and hare: the fast pointer laps the slow one if and only if a cycle exists; O(1) space.
\begin{lstlisting}[style=spec]
>>> has_cycle(from_list([1, 2, 3]))
False
>>> n = ListNode(1); n.next = n         # a self-loop
>>> has_cycle(n)
True
\end{lstlisting}

Roll your own: the same fast/slow walk, but watching for the pointers to MEET -- on a cycle the fast runner laps the slow one; on a straight list it simply falls off. O(1) space, no visited set.

\begin{lstlisting}[style=code]
def has_cycle(node):  # Floyd: fast laps slow iff a cycle exists
    slow = fast = node
    while fast and fast.next:
        slow, fast = slow.next, fast.next.next
        if slow is fast:
            return True
    return False
\end{lstlisting}

\section{merge\_lists}
\noindent\textit{[nodes]} --- merge two SORTED linked lists, using the dummy-head idiom to avoid special-casing the first node.
\begin{lstlisting}[style=spec]
>>> to_list(merge_lists(from_list([1, 3]), from_list([2, 4])))
[1, 2, 3, 4]
\end{lstlisting}

Roll your own: the dummy-head idiom -- a throwaway node to hang the result from, so appending the first real node needs no special case; walk both lists copying the smaller head, glue the survivor, return dummy.next.

\begin{lstlisting}[style=code]
def merge_lists(a, b):  # merge two sorted lists; dummy-head idiom
    dummy = tail_ = ListNode()
    while a and b:
        if a.val <= b.val:
            tail_.next, a = a, a.next
        else:
            tail_.next, b = b, b.next
        tail_ = tail_.next
    tail_.next = a or b
    return dummy.next
\end{lstlisting}

\section{neighbors4}
\noindent\textit{[grids \& windows]} --- the in-bounds orthogonal neighbours of cell (r, c) in an R-by-C grid. Bounds checking lives HERE, so BFS bodies stay clean.
\begin{lstlisting}[style=spec]
>>> sorted(neighbors4(0, 0, 2, 2))      # corner: two neighbours
[(0, 1), (1, 0)]
>>> sorted(neighbors4(1, 1, 3, 3))      # centre: four
[(0, 1), (1, 0), (1, 2), (2, 1)]
\end{lstlisting}

Roll your own: yield each in-bounds delta of (r, c) -- the bounds check lives HERE so every caller's BFS body stays clean. neighbors8 is the same generator over a 3x3 delta grid, with (dr or dc) excluding the centre cell.

\begin{lstlisting}[style=code]
def neighbors4(r, c, R, C):  # in-bounds orthogonal neighbors
    for dr, dc in ((1, 0), (-1, 0), (0, 1), (0, -1)):
        if 0 <= r + dr < R and 0 <= c + dc < C:
            yield r + dr, c + dc
\end{lstlisting}

\section{longest\_window}
\noindent\textit{[grids \& windows]} --- the sliding-window skeleton: it grows the right edge one element at a time (calling add), shrinks the left edge while your valid() says the window is illegal (calling rem), and tracks the best length. YOU supply the state as closures; the skeleton does the two-pointer bookkeeping. Longest- substring-without-repeats and longest-run-under-a- budget are both three closures away.
\begin{lstlisting}[style=spec]
>>> seen = []
>>> longest_window("abcabb", lambda: len(seen) == len(set(seen)),
...                seen.append, lambda c: seen.pop(0))
3
>>> w = []
>>> longest_window([2, 1, 3, 4], lambda: sum(w) <= 6,
...                w.append, lambda x: w.pop(0))
3
\end{lstlisting}

How it works: inversion of control. The skeleton owns the two pointers -- it grows hi one element per step (calling add), then shrinks lo while your valid() reports the window broken (calling rem per evicted element) -- and it tracks the best length ever seen legal. You own only the STATE, held in closures over variables sitting next to the call.


Designing the closures is answering one question: "what makes a window ILLEGAL, and what must I track to know?" Distinct characters -> a duplicate count. A budget -> a running sum. At most k of something -> a counter. add and rem are that state's increment and decrement; valid is the legality test.

\begin{lstlisting}[style=spec]
>>> w = []
>>> longest_window("aabcb", lambda: len(w) == len(set(w)), w.append, lambda c: w.pop(0))
3
\end{lstlisting}

Why closures and not a class: the state lives exactly as long as the call, sits in scope where you can read it, and needs no ceremony. Functions carrying state -- the same muscle as curry's captured argument, here working a two-pointer job.


Remember it as: the skeleton slides the window; your three closures say what legal means.


Roll your own: hi sweeps forward exactly once, add() greeting each newcomer; an inner while shrinks lo, rem() by rem(), until valid() approves again. Every element is added once and removed at most once -- the entire two-pointer O(n) argument, visible in two loops.

\begin{lstlisting}[style=code]
def longest_window(xs, valid, add, rem):  # sliding-window skeleton; state lives in the closures
    lo = best = 0
    for hi in range(len(xs)):
        add(xs[hi])
        while not valid():  # shrink until legal again
            rem(xs[lo])
            lo += 1
        best = max(best, hi - lo + 1)
    return best
\end{lstlisting}

\section{memo}
\noindent\textit{[control]} --- an unbounded memoizer as a plain decorator (no parentheses): results are cached by the argument tuple, so each distinct call computes once. Keyword arguments fold into the key; the cache is exposed as .cache.
\begin{lstlisting}[style=spec]
>>> @memo
... def fib(n): return n if n < 2 else fib(n - 1) + fib(n - 2)
>>> fib(30)
832040
>>> fib.cache[(10,)]
55
>>> len(fib.cache)                       # one entry per distinct argument
31
\end{lstlisting}

How it works: the decorator wraps f with a dict keyed on the arguments. First call computes and stores; every repeat is a lookup. That is the entire mechanism -- and why arguments must be HASHABLE: lists become tuples before they can be keys.


Why it earns its place: memoized recursion IS dynamic programming. Write the honest recurrence -- the answer for state (i, j) in terms of smaller states -- decorate it, and the cache is your DP table, filled in dependency order with no index bookkeeping. The exposed .cache lets you SEE the table: its size is the state count, exactly the number complexity analysis wants said out loud.

\begin{lstlisting}[style=spec]
>>> @memo
... def ways(r, c):                       # lattice paths to (r, c), moving right/down
...     return 1 if r == 0 or c == 0 else ways(r - 1, c) + ways(r, c - 1)
>>> ways(3, 3)
20
>>> len(ways.cache)                       # the 4x4 grid minus (0,0), never needed
15
\end{lstlisting}

Versus functools.lru\_cache: the standard library's version adds eviction (maxsize); memo is deliberately unbounded, because DP wants every state kept -- and it is nine lines you can retype in any interview that bans imports.


Remember it as: recursion states the truth; the cache makes it affordable; .cache is the table.


Roll your own: a dict captured in a closure, keyed by the argument tuple (keyword arguments folded in as a frozenset when present -- tuples and frozensets because keys must be hashable). Compute on miss, look up on hit, and hang the dict on the wrapper as .cache so the DP table stays inspectable.

\begin{lstlisting}[style=code]
def memo(f):  # unbounded memoizer; enough for DP (no eviction by design)
    cache = {}
    def wrapped(*args, **kw):
        key = (args, frozenset(kw.items())) if kw else args
        if key not in cache:
            cache[key] = f(*args, **kw)
        return cache[key]
    wrapped.cache = cache  # peek at the DP table if curious
    return wrapped
\end{lstlisting}

\section{until}
\noindent\textit{[control]} --- iterate f from x until the predicate holds, returning the first satisfying value: a fixed-point loop with no recursion-depth limit. The state can be a tuple carrying whatever the loop must remember.
\begin{lstlisting}[style=spec]
>>> until(lambda x: x > 100, lambda x: x * 2, 1)
128
>>> collatz = lambda s: (s[0] // 2 if even(s[0]) else 3 * s[0] + 1, s[1] + 1)
>>> until(lambda s: s[0] == 1, collatz, (6, 0))     # 8 steps to reach 1
(1, 8)
\end{lstlisting}

How it works: keep applying f until p approves -- iteration to a goal or a FIXED POINT, as a loop rather than recursion, so depth is unbounded. The state can be anything; a tuple lets one loop carry several facts at once, like the collatz example's (value, steps).


Why it earns its place: some processes have no list to fold over -- they run "until stable". Newton's isqrt converges; bellman\_ford relaxes every edge until the distance map stops changing (the textbook fixed point; g7's hint calls it by name). until states such loops declaratively: here is the step, here is done.

\begin{lstlisting}[style=spec]
>>> until(lambda n: n * n >= 30, succ, 0)      # smallest n with n squared >= 30
6
>>> until(lambda xs: xs == sorted(xs), sorted, [3, 1, 2])
[1, 2, 3]
\end{lstlisting}

Remember it as: while-not-done, promoted to an expression that returns its answer.


Roll your own: while not p(x): x = f(x); return x. Iteration promoted to a value-returning expression -- and a loop rather than recursion on purpose, because a fixed point may take unboundedly many steps.

\begin{lstlisting}[style=code]
def until(p, f, x):  # until p f x -- loop, not recursion (unbounded depth)
    while not p(x):
        x = f(x)
    return x
\end{lstlisting}
\part{The Nine Problem Classes}
\chapter{Class I --- Arithmetic and Unfolds}

The warm-up class, and the home of the \emph{unfold}: producing a sequence from
a seed instead of consuming one. Digits, base renderings and numeral systems all
peel a value apart step by step -- \texttt{unfoldr} when you want the pieces as
a list, \texttt{until} when you only want the final state, \texttt{foldl} when
a value table drives the steps. Watch how often a produced list is immediately
folded back into a value: unfold-then-fold is a pipeline, not a coincidence.
\medskip\noindent\textit{Practice files live in }\texttt{g1\_arith\_and\_unfolds/}\textit{;
each carries the prelude snippets it needs inline. Attempt every problem there
before reading on -- the doctests are the referee.}

\section{sign: Sign of a number}
\begin{lstlisting}[style=spec]
Return -1 for negative, 0 for zero, 1 for positive.

Contract:
    sign(x: int | float) -> int

Hint:
    The prelude ships this exact function: signum, one comparison
    subtraction, no branches. Recognize it, don't reinvent it.
\end{lstlisting}
\noindent\textit{Doctests:}
\begin{lstlisting}[style=spec]
>>> sign(-7)
-1
>>> sign(0)
0
>>> sign(3)
1
>>> sign(-0.5)
-1
\end{lstlisting}
\noindent\textbf{Solution.}
\begin{lstlisting}[style=code]
sign = signum
\end{lstlisting}

\section{rect\_area: Rectangle area}
\begin{lstlisting}[style=spec]
Contract:
    rect_area(w, h) -> w * h  (result type follows the inputs)

Hint:
    Deliberate overkill: area is the numeric fold, product([w, h]).
    See the smallest possible fold once; recognize every fold after.
\end{lstlisting}
\noindent\textit{Doctests:}
\begin{lstlisting}[style=spec]
>>> rect_area(3, 4)
12
>>> rect_area(2.5, 4)
10.0
>>> rect_area(0, 99)
0
\end{lstlisting}
\noindent\textbf{Solution.}
\begin{lstlisting}[style=code]
rect_area = lambda w, h: product([w, h])
\end{lstlisting}

\section{reverse\_digits: Reverse an integer's digits}
\begin{lstlisting}[style=spec]
Negatives keep their sign; trailing zeros vanish.

Contract:
    reverse_digits(n: int) -> int

Hint:
    unfoldr peels the digits with (m % 10, m // 10) -- least
    significant first, which IS reversed order. Fold them back into a
    number with acc * 10 + d; signum reapplies the sign.
\end{lstlisting}
\noindent\textit{Doctests:}
\begin{lstlisting}[style=spec]
>>> reverse_digits(123)
321
>>> reverse_digits(-45)
-54
>>> reverse_digits(120)
21
>>> reverse_digits(7)
7
\end{lstlisting}
\noindent\textbf{Solution.}
\begin{lstlisting}[style=code]
def reverse_digits(n):
    digits = unfoldr(lambda m: None if m == 0 else (m % 10, m // 10), abs(n))
    return signum(n) * foldl(lambda a, d: a * 10 + d, digits, 0)
\end{lstlisting}

\noindent\textit{Remark.} The state never needed to be a tuple threaded through
\texttt{until} -- \texttt{unfoldr} peels the digits into a list (already reversed,
because \texttt{m \% 10} yields the low digit first), and a fold reassembles them.
Unfold, then fold: when you see that pair, you are looking at a pipeline.

\section{to\_base: Render n in base b}
\begin{lstlisting}[style=spec]
Non-negative n, base 2..16, uppercase digits.

Contract:
    to_base(n: int, b: int) -> str

Hint:
    unfoldr emits digits low-to-high -- "0123456789ABCDEF"[m % b],
    seed m // b, stop at 0. Reverse and join; `or "0"` covers zero
    (its unfold is empty).
\end{lstlisting}
\noindent\textit{Doctests:}
\begin{lstlisting}[style=spec]
>>> to_base(255, 16)
'FF'
>>> to_base(5, 2)
'101'
>>> to_base(0, 8)
'0'
>>> to_base(100, 10)
'100'
>>> to_base(31, 16)
'1F'
\end{lstlisting}
\noindent\textbf{Solution.}
\begin{lstlisting}[style=code]
DIGITS = "0123456789ABCDEF"

def to_base(n, b):
    ds = unfoldr(lambda m: None if m == 0 else (DIGITS[m % b], m // b), n)
    return "".join(reversed(ds)) or "0"
\end{lstlisting}

\section{to\_roman: Integer to Roman numeral, 1 <= n <= 3999}
\begin{lstlisting}[style=spec]
Contract:
    to_roman(n: int) -> str

Hint:
    A left fold over the value table (subtractive pairs included:
    900 'CM', 40 'XL', ...) carrying (remaining, out): each step
    appends sym * (remaining // val) and keeps remaining % val.
\end{lstlisting}
\noindent\textit{Doctests:}
\begin{lstlisting}[style=spec]
>>> to_roman(1994)
'MCMXCIV'
>>> to_roman(9)
'IX'
>>> to_roman(58)
'LVIII'
>>> to_roman(3000)
'MMM'
>>> to_roman(1)
'I'
\end{lstlisting}
\noindent\textbf{Solution.}
\begin{lstlisting}[style=code]
TABLE = [(1000, "M"), (900, "CM"), (500, "D"), (400, "CD"), (100, "C"),
         (90, "XC"), (50, "L"), (40, "XL"), (10, "X"), (9, "IX"),
         (5, "V"), (4, "IV"), (1, "I")]

def to_roman(n):
    step = lambda acc, vs: (acc[0] % vs[0], acc[1] + vs[1] * (acc[0] // vs[0]))
    return foldl(step, TABLE, (n, ""))[1]
\end{lstlisting}

\section{from\_roman: Parse a valid Roman numeral}
\begin{lstlisting}[style=spec]
Contract:
    from_roman(s: str) -> int

Hint:
    pairwise gives each symbol its right neighbor; a value counts
    NEGATIVE when its neighbor is bigger ('IV' = -1 + 5). Pad the
    value list with a trailing 0 so the last symbol gets a neighbor
    too, then sum.
\end{lstlisting}
\noindent\textit{Doctests:}
\begin{lstlisting}[style=spec]
>>> from_roman('MCMXCIV')
1994
>>> from_roman('IX')
9
>>> from_roman('LVIII')
58
>>> from_roman('III')
3
>>> from_roman('MMM')
3000
\end{lstlisting}
\noindent\textbf{Solution.}
\begin{lstlisting}[style=code]
VAL = {"I": 1, "V": 5, "X": 10, "L": 50, "C": 100, "D": 500, "M": 1000}

def from_roman(s):
    vs = map_(lambda c: VAL[c], s) + [0]
    return sum(map_(lambda p: -p[0] if p[0] < p[1] else p[0], pairwise(vs)))
\end{lstlisting}

\noindent\textit{Remark.} Every subtractive rule in Roman numerals is a statement about
a symbol and its \emph{right neighbour}. \texttt{pairwise} materialises exactly
that relationship, and the padding zero gives the final symbol a neighbour too.
When a rule talks about adjacent elements, reach for \texttt{pairwise} before you
reach for indices.
\chapter{Class II --- Lists and Strings}

Bread and butter: \texttt{map\_}, \texttt{filter\_}, \texttt{concat},
\texttt{concatMap}, and the string pair \texttt{words}/\texttt{unwords}. Two
problems here are deliberate one-liners (\texttt{flatten} \emph{is}
\texttt{concat}, \texttt{transpose\_matrix} \emph{is} \texttt{transpose}):
recognising an existing tool is a skill worth isolating and drilling. The
harder members (\texttt{chunks}, \texttt{permutations}) preview unfolds and
backtracking enumeration.
\medskip\noindent\textit{Practice files live in }\texttt{g2\_lists\_and\_strings/}\textit{;
each carries the prelude snippets it needs inline. Attempt every problem there
before reading on -- the doctests are the referee.}

\section{greet: Greeting string}
\begin{lstlisting}[style=spec]
Contract:
    greet(name: str) -> 'Hello, <name>!'

Hint:
    unwords joins words with single spaces; the bang concatenates.
    Trivial on purpose -- the warmup proves the pipeline.
\end{lstlisting}
\noindent\textit{Doctests:}
\begin{lstlisting}[style=spec]
>>> greet("Ada")
'Hello, Ada!'
>>> greet("world")
'Hello, world!'
\end{lstlisting}
\noindent\textbf{Solution.}
\begin{lstlisting}[style=code]
greet = lambda name: unwords(["Hello,", name]) + "!"
\end{lstlisting}

\section{count\_vowels: Count vowels, case-insensitive}
\begin{lstlisting}[style=spec]
Contract:
    count_vowels(s: str) -> int

Hint:
    len . filter_ (in "aeiou") . lower -- a filter consumed by len is
    the counting idiom.
\end{lstlisting}
\noindent\textit{Doctests:}
\begin{lstlisting}[style=spec]
>>> count_vowels("Programming")
3
>>> count_vowels("AEIOU")
5
>>> count_vowels("xyz")
0
>>> count_vowels("")
0
\end{lstlisting}
\noindent\textbf{Solution.}
\begin{lstlisting}[style=code]
count_vowels = lambda s: len(filter_(lambda c: c in "aeiou", s.lower()))
\end{lstlisting}

\section{fizzbuzz: FizzBuzz as data}
\begin{lstlisting}[style=spec]
For 1..n: multiples of 3 -> "Fizz", of 5 -> "Buzz", of both ->
"FizzBuzz", else the number as a string.

Contract:
    fizzbuzz(n: int) -> list[str]

Hint:
    Solve it BOTH ways and compare -- the contrast is the lesson:
    1. COMPOSITION OF ARMS: map_ one per-number rule over
       range(1, n+1), building the word by concatenating guards --
       "Fizz" * (i % 3 == 0) + "Buzz" * (i % 5 == 0) -- with
       `or str(i)` catching the empty case. Overlaps compose:
       15 earns "FizzBuzz" from two rules for free.
    2. FIRST-MATCH SELECTION: a RULES table of (predicate, result)
       pairs, find picking the first hit; const(otherwise) is the
       catch-all arm. Overlaps need their own row (15 goes first),
       but the ladder is now DATA -- appendable, testable rule by
       rule. Which generalizes better depends on whether cases
       overlap.
\end{lstlisting}
\noindent\textit{Doctests:}
\begin{lstlisting}[style=spec]
>>> fizzbuzz(5)
['1', '2', 'Fizz', '4', 'Buzz']
>>> fizzbuzz(15)[-1]
'FizzBuzz'
>>> fizzbuzz(0)
[]
\end{lstlisting}
\noindent\textbf{Solution.}
\begin{lstlisting}[style=code]
# way 1 -- composition of arms: overlapping guards concatenate
fizzbuzz = lambda n: map_(
    lambda i: ("Fizz" * (i % 3 == 0) + "Buzz" * (i % 5 == 0)) or str(i),
    range(1, n + 1))

# way 2 -- first-match selection: the ladder as a RULES table (this one wins)
RULES = [(lambda i: i % 15 == 0, const("FizzBuzz")),
         (lambda i: i % 3 == 0,  const("Fizz")),
         (lambda i: i % 5 == 0,  const("Buzz")),
         (const(otherwise),      str)]
speak = lambda i: snd(find(lambda rule: fst(rule)(i), RULES))(i)
fizzbuzz = lambda n: map_(speak, range(1, n + 1))
\end{lstlisting}

\section{word\_lengths: Map each word to its length}
\begin{lstlisting}[style=spec]
Contract:
    word_lengths(sentence: str) -> dict[str, int]
    First-appearance key order; repeats collapse harmlessly.

Hint:
    words splits; zip_ the words against map_ len of them; dict() the
    pairs. Duplicate keys overwrite with the same value -- free dedup.
\end{lstlisting}
\noindent\textit{Doctests:}
\begin{lstlisting}[style=spec]
>>> word_lengths("the quick fox")
{'the': 3, 'quick': 5, 'fox': 3}
>>> word_lengths("go go go")
{'go': 2}
>>> word_lengths("")
{}
\end{lstlisting}
\noindent\textbf{Solution.}
\begin{lstlisting}[style=code]
word_lengths = lambda s: dict(zip_(words(s), map_(len, words(s))))
\end{lstlisting}

\section{flatten: Flatten one level of nesting}
\begin{lstlisting}[style=spec]
Contract:
    flatten(lists: list[list]) -> list

Hint:
    This IS the prelude's concat. One name, zero work.
\end{lstlisting}
\noindent\textit{Doctests:}
\begin{lstlisting}[style=spec]
>>> flatten([[1, 2], [3], [4, 5]])
[1, 2, 3, 4, 5]
>>> flatten([[], [1], []])
[1]
>>> flatten([])
[]
\end{lstlisting}
\noindent\textbf{Solution.}
\begin{lstlisting}[style=code]
flatten = concat
\end{lstlisting}

\section{caesar: Caesar cipher, letters only, case preserved}
\begin{lstlisting}[style=spec]
Contract:
    caesar(s: str, k: int) -> str

Hint:
    map_ a per-char shift over the string; the wrap is mod 26 on
    ord() offsets; join the pieces. One helper, one map.
\end{lstlisting}
\noindent\textit{Doctests:}
\begin{lstlisting}[style=spec]
>>> caesar("abc", 2)
'cde'
>>> caesar("XYZ", 3)
'ABC'
>>> caesar("a-b", 1)
'b-c'
>>> caesar("Hello, World!", 13)
'Uryyb, Jbeyq!'
>>> caesar("abc", 0)
'abc'
\end{lstlisting}
\noindent\textbf{Solution.}
\begin{lstlisting}[style=code]
def caesar(s, k):
    def shift(c):
        if not c.isalpha():
            return c
        base = ord("A") if c.isupper() else ord("a")
        return chr(base + (ord(c) - base + k) % 26)
    return "".join(map_(shift, s))
\end{lstlisting}

\section{contains\_duplicate: NeetCode: Contains Duplicate}
\begin{lstlisting}[style=spec]
Contract:
    contains_duplicate(nums: list[int]) -> bool
    True iff any value appears at least twice.

Hint:
    nub keeps first occurrences; duplicates exist iff nub shrinks the
    list.
\end{lstlisting}
\noindent\textit{Doctests:}
\begin{lstlisting}[style=spec]
>>> contains_duplicate([1, 2, 3, 1])
True
>>> contains_duplicate([1, 2, 3, 4])
False
>>> contains_duplicate([])
False
\end{lstlisting}
\noindent\textbf{Solution.}
\begin{lstlisting}[style=code]
contains_duplicate = lambda nums: len(nub(nums)) != len(nums)
\end{lstlisting}

\section{chunks: Split a list into pieces of length k}
\begin{lstlisting}[style=spec]
Last piece may be shorter.

Contract:
    chunks(xs: list, k: int) -> list[list]

Hint:
    splitAt already returns (front, rest) == (value, next_seed), so
    chunks IS unfoldr of splitAt; null(rest) is the stop. One line.
    (You are deriving the prelude's chunksOf -- deriving a tool once
    is why it sticks.)
\end{lstlisting}
\noindent\textit{Doctests:}
\begin{lstlisting}[style=spec]
>>> chunks([1, 2, 3, 4, 5], 2)
[[1, 2], [3, 4], [5]]
>>> chunks([1, 2, 3, 4], 2)
[[1, 2], [3, 4]]
>>> chunks([], 3)
[]
>>> chunks([1], 5)
[[1]]
\end{lstlisting}
\noindent\textbf{Solution.}
\begin{lstlisting}[style=code]
chunks = lambda xs, k: unfoldr(
    lambda rest: None if null(rest) else splitAt(k, rest), list(xs))
\end{lstlisting}

\noindent\textit{Remark.} \texttt{splitAt} already returns \texttt{(front, rest)} -- which
is precisely the \texttt{(value, next\_seed)} contract of \texttt{unfoldr}. The
whole solution is recognising that two signatures line up. You are deriving the
prelude's own \texttt{chunksOf}; deriving a tool once is why it sticks.

\section{permutations: All distinct orderings, sorted}
\begin{lstlisting}[style=spec]
Contract:
    permutations(s: str) -> list[str]

Hint:
    concatMap over the choice of first character: pick s[i], prepend
    it to every permutation of the rest. nub kills duplicates from
    repeated letters; sort at the end.
\end{lstlisting}
\noindent\textit{Doctests:}
\begin{lstlisting}[style=spec]
>>> permutations("abc")
['abc', 'acb', 'bac', 'bca', 'cab', 'cba']
>>> permutations("aab")
['aab', 'aba', 'baa']
>>> permutations("x")
['x']
>>> permutations("")
['']
\end{lstlisting}
\noindent\textbf{Solution.}
\begin{lstlisting}[style=code]
def permutations(s):
    if len(s) <= 1:
        return [s]
    opts = concatMap(
        lambda i: [s[i] + rest for rest in permutations(s[:i] + s[i + 1:])],
        range(len(s)))
    return sorted(nub(opts))
\end{lstlisting}

\noindent\textit{Remark.} \texttt{concatMap} is the shape of ``for every choice,
recurse, and pool the results.'' Any backtracking enumeration can be written this
way; \texttt{nub} then absorbs duplicate letters without a special case.

\section{reverse\_words: Reverse the word order}
\begin{lstlisting}[style=spec]
Runs of whitespace collapse (that's what a bare split does).

Contract:
    reverse_words(s: str) -> str

Hint:
    An unwords . reverse . words sandwich -- three names, no loop.
\end{lstlisting}
\noindent\textit{Doctests:}
\begin{lstlisting}[style=spec]
>>> reverse_words("the quick brown fox")
'fox brown quick the'
>>> reverse_words("  hello   world  ")
'world hello'
>>> reverse_words("solo")
'solo'
\end{lstlisting}
\noindent\textbf{Solution.}
\begin{lstlisting}[style=code]
reverse_words = lambda s: unwords(list(reversed(words(s))))
\end{lstlisting}

\section{transpose\_matrix: Flip rows and columns}
\begin{lstlisting}[style=spec]
Contract:
    transpose_matrix(rows: list[list]) -> list[list]

Hint:
    This IS the prelude's transpose -- like flatten was concat.
    One name, zero work; know WHY zip(*rows) does it.
\end{lstlisting}
\noindent\textit{Doctests:}
\begin{lstlisting}[style=spec]
>>> transpose_matrix([[1, 2, 3], [4, 5, 6]])
[[1, 4], [2, 5], [3, 6]]
>>> transpose_matrix([[5]])
[[5]]
>>> transpose_matrix([[7], [8], [9]])
[[7, 8, 9]]
\end{lstlisting}
\noindent\textbf{Solution.}
\begin{lstlisting}[style=code]
transpose_matrix = transpose
\end{lstlisting}

\section{rotate\_right: Rotate a list right by k}
\begin{lstlisting}[style=spec]
k may be 0 or larger than the list.

Contract:
    rotate_right(xs: list, k: int) -> list

Hint:
    One splitAt at len - k%len, then glue the halves swapped.
    The k=0 trap: xs[-0:] thinking sneaks in if you slice by -k.
\end{lstlisting}
\noindent\textit{Doctests:}
\begin{lstlisting}[style=spec]
>>> rotate_right([1, 2, 3, 4, 5], 2)
[4, 5, 1, 2, 3]
>>> rotate_right([1, 2], 0)
[1, 2]
>>> rotate_right([1, 2, 3], 5)
[2, 3, 1]
>>> rotate_right([], 3)
[]
\end{lstlisting}
\noindent\textbf{Solution.}
\begin{lstlisting}[style=code]
def rotate_right(xs, k):
    if not xs:
        return []
    a, b = splitAt(len(xs) - k % len(xs), list(xs))
    return b + a
\end{lstlisting}
\chapter{Class III --- Folds, Scans and Streams}

The heart of functional problem solving. A fold collapses a list to one value;
a scan is a fold that keeps its history; a stream generates lazily until
\texttt{take} says stop. The famous surprises live here: Kadane's algorithm is
a \texttt{scanl1}, Fibonacci is \texttt{iterate} on a pair, prefix sums turn
every window-sum question into two slices and a \texttt{zipWith}. When a
statement says ``running'', ``so far'' or ``at each step'', you are being told
the answer is a scan.
\medskip\noindent\textit{Practice files live in }\texttt{g3\_folds\_scans\_streams/}\textit{;
each carries the prelude snippets it needs inline. Attempt every problem there
before reading on -- the doctests are the referee.}

\section{two\_sum: NeetCode: Two Sum}
\begin{lstlisting}[style=spec]
Contract:
    two_sum(nums: list[int], target: int) -> tuple[int, int]
    Return indices (i, j), i < j, with nums[i] + nums[j] == target.
    Exactly one solution exists. len(nums) >= 2.

Hint:
    A fold over enum(nums) carrying a {value: index} dict.
    The future needs to know: which values have I seen, and where?
\end{lstlisting}
\noindent\textit{Doctests:}
\begin{lstlisting}[style=spec]
>>> two_sum([2, 7, 11, 15], 9)
(0, 1)
>>> two_sum([3, 2, 4], 6)
(1, 2)
>>> two_sum([3, 3], 6)
(0, 1)
\end{lstlisting}
\noindent\textbf{Solution.}
\begin{lstlisting}[style=code]
def two_sum(nums, target):
    seen = {}
    for i, x in enum(nums):
        if target - x in seen:
            return (seen[target - x], i)
        seen[x] = i
\end{lstlisting}

\section{best\_time\_stock: NeetCode: Best Time to Buy and Sell Stock}
\begin{lstlisting}[style=spec]
Contract:
    max_profit(prices: list[int]) -> int
    One buy, one later sell. Best achievable profit, 0 if none.
    len(prices) >= 1.

Hint:
    scanl1 min gives the running minimum -- the best buy so far.
    zipWith (-) prices against it, then maximum. A scan, not a loop.
\end{lstlisting}
\noindent\textit{Doctests:}
\begin{lstlisting}[style=spec]
>>> max_profit([7, 1, 5, 3, 6, 4])
5
>>> max_profit([7, 6, 4, 3, 1])
0
>>> max_profit([5])
0
\end{lstlisting}
\noindent\textbf{Solution.}
\begin{lstlisting}[style=code]
def max_profit(prices):
    mins = scanl1(min, prices)
    return max(zipWith(lambda p, m: p - m, prices, mins))
\end{lstlisting}

\section{maximum\_subarray: NeetCode: Maximum Subarray (Kadane)}
\begin{lstlisting}[style=spec]
Contract:
    max_subarray(nums: list[int]) -> int
    Largest sum of a non-empty contiguous subarray. len(nums) >= 1.

Hint:
    Kadane IS scanl1: best-ending-here = max(x, best_prev + x).
    The scan is the DP table; maximum reads the answer off it.
\end{lstlisting}
\noindent\textit{Doctests:}
\begin{lstlisting}[style=spec]
>>> max_subarray([-2, 1, -3, 4, -1, 2, 1, -5, 4])
6
>>> max_subarray([1])
1
>>> max_subarray([5, 4, -1, 7, 8])
23
>>> max_subarray([-3, -1, -2])
-1
\end{lstlisting}
\noindent\textbf{Solution.}
\begin{lstlisting}[style=code]
def max_subarray(nums):
    return max(scanl1(lambda best, x: max(x, best + x), nums))
\end{lstlisting}

\noindent\textit{Remark.} Kadane's algorithm is famous as a clever trick; seen
functionally it is one \texttt{scanl1}. The scan \emph{is} the DP table
(best sum ending at each index), and \texttt{max} reads the answer off it. Many
celebrated one-pass algorithms are scans wearing a trench coat.

\section{product\_except\_self: NeetCode: Product of Array Except Self}
\begin{lstlisting}[style=spec]
Contract:
    product_except_self(nums: list[int]) -> list[int]
    out[i] == product of all nums except nums[i]. No division. O(n).
    len(nums) >= 2.

Hint:
    Prefix products = scanl (*) 1, suffix products = scanr (*) 1.
    out = zipWith (*) prefixes suffixes -- offset by one on each side.
\end{lstlisting}
\noindent\textit{Doctests:}
\begin{lstlisting}[style=spec]
>>> product_except_self([1, 2, 3, 4])
[24, 12, 8, 6]
>>> product_except_self([-1, 1, 0, -3, 3])
[0, 0, 9, 0, 0]
\end{lstlisting}
\noindent\textbf{Solution.}
\begin{lstlisting}[style=code]
def product_except_self(nums):
    mul = lambda a, b: a * b
    pre = scanl(mul, 1, nums)
    suf = scanr(mul, 1, nums)
    return zipWith(mul, pre[:-1], suf[1:])
\end{lstlisting}

\noindent\textit{Remark.} One scan from the left, one from the right, and a
\texttt{zipWith} to marry them. The offset slicing (\texttt{pre[:-1]},
\texttt{suf[1:]}) is where the ``except self'' lives -- each position sees the
product of everything strictly before and strictly after it.

\section{moving\_average: Average of every k-window}
\begin{lstlisting}[style=spec]
Contract:
    moving_average(xs: list, k: int) -> list[float]
    len(xs) >= k >= 1.

Hint:
    Prefix sums = scanl (+) 0 -- one scan, then every window sum is
    pref[i+k] - pref[i]: zipWith (-) the two offset views of the same
    scan, divide by k.
\end{lstlisting}
\noindent\textit{Doctests:}
\begin{lstlisting}[style=spec]
>>> moving_average([1, 2, 3, 4], 2)
[1.5, 2.5, 3.5]
>>> moving_average([1, 1, 1], 3)
[1.0]
>>> moving_average([5], 1)
[5.0]
\end{lstlisting}
\noindent\textbf{Solution.}
\begin{lstlisting}[style=code]
def moving_average(xs, k):
    pref = scanl(lambda a, b: a + b, 0, xs)
    return zipWith(lambda hi, lo: (hi - lo) / k, pref[k:], pref[:-k])
\end{lstlisting}

\section{pascal: First n rows of Pascal's triangle}
\begin{lstlisting}[style=spec]
Contract:
    pascal(n: int) -> list[list[int]]

Hint:
    The triangle is a stream: iterate the row step from [1], where
    next = zipWith (+) (0:row) (row:0) -- the row against itself,
    shifted. take n consumes the stream.
\end{lstlisting}
\noindent\textit{Doctests:}
\begin{lstlisting}[style=spec]
>>> pascal(4)
[[1], [1, 1], [1, 2, 1], [1, 3, 3, 1]]
>>> pascal(1)
[[1]]
>>> pascal(0)
[]
>>> pascal(6)[-1]
[1, 5, 10, 10, 5, 1]
\end{lstlisting}
\noindent\textbf{Solution.}
\begin{lstlisting}[style=code]
def pascal(n):
    step = lambda row: zipWith(lambda a, b: a + b, [0] + row, row + [0])
    return take(n, iterate(step, [1]))
\end{lstlisting}

\section{climbing\_stairs: NeetCode: Climbing Stairs}
\begin{lstlisting}[style=spec]
Contract:
    climb_stairs(n: int) -> int
    Ways to reach step n taking 1 or 2 steps. n >= 1.

Hint:
    Fibonacci in disguise. iterate the pair step, take the nth --
    an unfold consumed by indexing, no table needed.
\end{lstlisting}
\noindent\textit{Doctests:}
\begin{lstlisting}[style=spec]
>>> climb_stairs(1)
1
>>> climb_stairs(2)
2
>>> climb_stairs(5)
8
>>> climb_stairs(45)
1836311903
\end{lstlisting}
\noindent\textbf{Solution.}
\begin{lstlisting}[style=code]
def climb_stairs(n):
    step = lambda p: (snd(p), fst(p) + snd(p))
    return fst(last(take(n + 1, iterate(step, (1, 1)))))
\end{lstlisting}

\section{primes\_up\_to: All primes <= n, ascending}
\begin{lstlisting}[style=spec]
Contract:
    primes_up_to(n: int) -> list[int]

Hint:
    Streams, Haskell style: candidates = takeWhile (<= n) count(2);
    keep p when no divisor in takeWhile (<= isqrt p) count(2) divides
    it -- all over the remainders. (The array sieve is faster; this
    version is about reading streams.)
\end{lstlisting}
\noindent\textit{Doctests:}
\begin{lstlisting}[style=spec]
>>> primes_up_to(20)
[2, 3, 5, 7, 11, 13, 17, 19]
>>> primes_up_to(2)
[2]
>>> primes_up_to(1)
[]
>>> len(primes_up_to(100))
25
\end{lstlisting}
\noindent\textbf{Solution.}
\begin{lstlisting}[style=code]
def primes_up_to(n):
    cands = takeWhile(lambda x: x <= n, count(2))
    is_prime = lambda p: all(
        p % d for d in takeWhile(lambda d: d <= isqrt(p), count(2)))
    return filter_(is_prime, cands)
\end{lstlisting}

\section{totals\_and\_deltas: A scan and its inverse}
\begin{lstlisting}[style=spec]
Contract:
    running_totals(xs: list) -> list   # cumulative sums
    deltas(xs: list) -> list           # each element minus the previous

Hint:
    running_totals IS scanl1 (+). deltas is the pairwise differences.
    They undo each other: prepend the first element to
    deltas(running_totals(xs)) and xs comes back -- say why.
\end{lstlisting}
\noindent\textit{Doctests:}
\begin{lstlisting}[style=spec]
>>> running_totals([3, 1, 4])
[3, 4, 8]
>>> running_totals([])
[]
>>> deltas([5, 3, 8])
[-2, 5]
>>> deltas([7])
[]
>>> ts = running_totals([2, 5, 1]); [2] + deltas(ts) == [2, 5, 1]
True
\end{lstlisting}
\noindent\textbf{Solution.}
\begin{lstlisting}[style=code]
running_totals = lambda xs: scanl1(lambda a, b: a + b, xs)
deltas = lambda xs: map_(lambda p: p[1] - p[0], pairwise(xs))
\end{lstlisting}

\noindent\textit{Remark.} \texttt{scanl1 (+)} and pairwise differences are
inverses -- discrete integration and differentiation. Recognising an operation's
inverse is a testing superpower: the round-trip doctest here is a property test
in miniature.

\section{paren\_depth: Deepest paren nesting}
\begin{lstlisting}[style=spec]
Input is balanced, only '(' and ')'.

Contract:
    max_depth(s: str) -> int

Hint:
    Depth is a running total: map each paren to +1/-1, scanl (+) 0
    is the entire depth trace, maximum reads off the answer. No
    stack needed -- the scan IS the stack height's history.
\end{lstlisting}
\noindent\textit{Doctests:}
\begin{lstlisting}[style=spec]
>>> max_depth("(()(()))")
3
>>> max_depth("()()")
1
>>> max_depth("")
0
\end{lstlisting}
\noindent\textbf{Solution.}
\begin{lstlisting}[style=code]
def max_depth(s):
    steps = map_(lambda c: 1 if c == "(" else -1, s)
    return max(scanl(lambda a, b: a + b, 0, steps))
\end{lstlisting}
\chapter{Class IV --- Sorting and Searching}

``Does order unlock it?'' is the first question of the recognition ladder for a
reason: sorting linearises a problem so that greedy passes, binary searches and
pairwise arguments become available. The exchange argument (\texttt{photo\_%
lineup}), the tie-break trick (\texttt{russian\_doll}) and patience sorting
(\texttt{lis}) are all consequences of choosing the \emph{right} sort key --
which is where the thinking happens. \texttt{sensor\_pairing} closes the class
with the interview's own domain problem.
\medskip\noindent\textit{Practice files live in }\texttt{g4\_sorting\_and\_searching/}\textit{;
each carries the prelude snippets it needs inline. Attempt every problem there
before reading on -- the doctests are the referee.}

\section{is\_anagram\_sorted: Anagrams via canonical form}
\begin{lstlisting}[style=spec]
True when a and b use exactly the same letters, ignoring case and
non-letters. (g5 solves the same problem with bags -- do both, argue
which you'd ship.)

Contract:
    is_anagram(a: str, b: str) -> bool

Hint:
    A canonical form makes equality trivial: sort of the cleaned,
    lowered characters. filter_ str.isalpha does the cleaning.
\end{lstlisting}
\noindent\textit{Doctests:}
\begin{lstlisting}[style=spec]
>>> is_anagram("Listen", "Silent")
True
>>> is_anagram("hello", "world")
False
>>> is_anagram("a gentleman", "elegant man")
True
>>> is_anagram("ab", "abb")
False
\end{lstlisting}
\noindent\textbf{Solution.}
\begin{lstlisting}[style=code]
canon = lambda s: sorted(filter_(str.isalpha, s.lower()))
is_anagram = lambda a, b: canon(a) == canon(b)
\end{lstlisting}

\section{second\_largest: Second largest distinct value}
\begin{lstlisting}[style=spec]
At least two distinct values guaranteed.

Contract:
    second_largest(xs: list) -> value

Hint:
    nub then sort; the answer is one from the end.
\end{lstlisting}
\noindent\textit{Doctests:}
\begin{lstlisting}[style=spec]
>>> second_largest([1, 5, 3])
3
>>> second_largest([5, 1, 5, 3])
3
>>> second_largest([2, 1])
1
>>> second_largest([-1, -5, -3])
-3
\end{lstlisting}
\noindent\textbf{Solution.}
\begin{lstlisting}[style=code]
second_largest = lambda xs: sorted(nub(xs))[-2]
\end{lstlisting}

\section{binary\_search: Index of target in a sorted list, else -1}
\begin{lstlisting}[style=spec]
Contract:
    binary_search(xs: list, target) -> int

Hint:
    bisect_left finds the leftmost slot where the target COULD be;
    one bounds-check plus one equality test finishes the job. Know
    why bisect_left(xs, x) == len(xs) is possible.
\end{lstlisting}
\noindent\textit{Doctests:}
\begin{lstlisting}[style=spec]
>>> binary_search([1, 3, 5, 7], 5)
2
>>> binary_search([1, 3, 5, 7], 1)
0
>>> binary_search([1, 3, 5, 7], 7)
3
>>> binary_search([1, 3], 2)
-1
>>> binary_search([], 1)
-1
\end{lstlisting}
\noindent\textbf{Solution.}
\begin{lstlisting}[style=code]
def binary_search(xs, x):
    i = bisect_left(xs, x)
    return i if i < len(xs) and xs[i] == x else -1
\end{lstlisting}

\section{merge\_intervals: Merge overlapping (and touching) intervals}
\begin{lstlisting}[style=spec]
Intervals are inclusive [s, e]: [1,2] and [2,3] merge to [1,3].

Contract:
    merge_intervals(intervals: list[list[int]]) -> list[list[int]]
    Result sorted by start.

Hint:
    sortOn fst, then one fold: each interval either extends the
    accumulator's last entry or starts a new one.
\end{lstlisting}
\noindent\textit{Doctests:}
\begin{lstlisting}[style=spec]
>>> merge_intervals([[1, 3], [2, 6], [8, 10], [15, 18]])
[[1, 6], [8, 10], [15, 18]]
>>> merge_intervals([[1, 2], [2, 3]])
[[1, 3]]
>>> merge_intervals([[5, 7], [1, 2], [3, 6]])
[[1, 2], [3, 7]]
>>> merge_intervals([[4, 4]])
[[4, 4]]
\end{lstlisting}
\noindent\textbf{Solution.}
\begin{lstlisting}[style=code]
def merge_intervals(intervals):
    def step(acc, iv):
        s, e = iv
        if acc and s <= acc[-1][1]:
            acc[-1][1] = max(acc[-1][1], e)
        else:
            acc.append([s, e])
        return acc
    return foldl(step, sortOn(fst, [list(iv) for iv in intervals]), [])
\end{lstlisting}

\section{photo\_lineup: The reported Anduril prompt, prelude-sized}
\begin{lstlisting}[style=spec]
Back row and front row, same length; can the columns be arranged so
every back person is STRICTLY taller than the front person?

Contract:
    can_lineup(back: list[int], front: list[int]) -> bool

Hint:
    Sort both rows; all over zipWith (>). Be ready to argue why
    sorted-vs-sorted decides it (the exchange argument).
\end{lstlisting}
\noindent\textit{Doctests:}
\begin{lstlisting}[style=spec]
>>> can_lineup([5, 1, 3], [2, 4, 0])
True
>>> can_lineup([1, 2, 3], [1, 2, 3])
False
>>> can_lineup([10, 9], [1, 2])
True
>>> can_lineup([2], [3])
False
\end{lstlisting}
\noindent\textbf{Solution.}
\begin{lstlisting}[style=code]
can_lineup = lambda back, front: all(
    zipWith(lambda b, f: b > f, sorted(back), sorted(front)))
\end{lstlisting}

\section{merge\_sort: A real O(n log n) sort, by hand}
\begin{lstlisting}[style=spec]
No sorted(), no .sorted().

Contract:
    merge_sorted(xs: list) -> list  (new list; stable)

Hint:
    splitAt the midpoint, recurse on both halves, and the prelude's
    merge does the two-pointer join. Also retype merge itself from
    memory once -- it's the loop interviews ask for.
\end{lstlisting}
\noindent\textit{Doctests:}
\begin{lstlisting}[style=spec]
>>> merge_sorted([5, 2, 8, 1])
[1, 2, 5, 8]
>>> merge_sorted([])
[]
>>> merge_sorted([3, 3, 1])
[1, 3, 3]
>>> merge_sorted([1])
[1]
\end{lstlisting}
\noindent\textbf{Solution.}
\begin{lstlisting}[style=code]
def merge_sorted(xs):
    if len(xs) <= 1:
        return list(xs)
    a, b = splitAt(len(xs) // 2, list(xs))
    return merge(merge_sorted(a), merge_sorted(b))
\end{lstlisting}

\section{lis: Longest strictly increasing subsequence (length)}
\begin{lstlisting}[style=spec]
Contract:
    lis(xs: list[int]) -> int

Hint:
    Patience sorting: fold the sequence into `tails`, where tails[i]
    is the smallest possible tail of an increasing run of length i+1.
    bisect_left says which pile each element lands on; a landing past
    the end grows the answer. O(n log n) where the obvious DP is
    O(n^2) -- say both.
\end{lstlisting}
\noindent\textit{Doctests:}
\begin{lstlisting}[style=spec]
>>> lis([10, 9, 2, 5, 3, 7, 101, 18])
4
>>> lis([5, 4, 3])
1
>>> lis([1, 2, 3])
3
>>> lis([])
0
>>> lis([7, 7, 7])
1
\end{lstlisting}
\noindent\textbf{Solution.}
\begin{lstlisting}[style=code]
def lis(xs):
    def place(tails, x):
        i = bisect_left(tails, x)
        return tails + [x] if i == len(tails) else tails[:i] + [x] + tails[i + 1:]
    return len(foldl(place, xs, []))
\end{lstlisting}

\noindent\textit{Remark.} Patience sorting: \texttt{tails[i]} is the smallest possible tail
of an increasing run of length $i{+}1$; \texttt{bisect\_left} says which pile
each element lands on. Landing past the end grows the answer. Know the
$O(n^2)$ DP too, and say both complexities out loud.

\section{russian\_doll: LC 354: Russian Doll Envelopes}
\begin{lstlisting}[style=spec]
Envelope (w, h) fits inside another only when BOTH are strictly
smaller. No rotation. Longest nesting chain.

Contract:
    max_dolls(envelopes: list[tuple[int, int]]) -> int

Hint:
    Sort by (width asc, height DESC), then plain LIS on the heights
    -- g4_07 verbatim. The descending tie-break is the whole trick:
    equal-width envelopes arrive tallest-first, so no two of them can
    ever chain in a strictly-increasing height run. Sorting
    linearized one dimension; LIS handles the other. O(n log n).
\end{lstlisting}
\noindent\textit{Doctests:}
\begin{lstlisting}[style=spec]
>>> max_dolls([(5, 4), (6, 4), (6, 7), (2, 3)])
3
>>> max_dolls([(1, 1), (1, 1), (1, 1)])
1
>>> max_dolls([(3, 4), (3, 5), (3, 6)])
1
>>> max_dolls([])
0
>>> max_dolls([(2, 3)])
1
\end{lstlisting}
\noindent\textbf{Solution.}
\begin{lstlisting}[style=code]
def max_dolls(envelopes):
    hs = [h for _, h in sortOn(lambda e: (e[0], -e[1]), envelopes)]
    def place(tails, x):
        i = bisect_left(tails, x)
        return tails + [x] if i == len(tails) else tails[:i] + [x] + tails[i + 1:]
    return len(foldl(place, hs, []))
\end{lstlisting}

\noindent\textit{Remark.} Sorting by \texttt{(width asc, height desc)} is the whole
problem. The descending tie-break guarantees equal-width envelopes can never
chain, which collapses a 2-D dominance question into a 1-D LIS on heights.
Sorting linearised one dimension; the tails trick handles the other.

\section{pairs\_within\_budget: Count pairs whose sum fits a budget}
\begin{lstlisting}[style=spec]
Pairs are (i, j) with i < j; values may be negative; the count does
not depend on input order -- that sentence licenses sorting.

Contract:
    pairs_within_budget(t: int, xs: list[int]) -> int

Hint:
    Sort first. Then each element asks: how many LATER elements keep
    the sum <= t? bisect_right(ys, t - x) answers in O(log n) --
    subtract the i+1 elements at or before x itself, clip at 0.
\end{lstlisting}
\noindent\textit{Doctests:}
\begin{lstlisting}[style=spec]
>>> pairs_within_budget(5, [3, 1, 4, 2])
4
>>> pairs_within_budget(-1, [-2, -3, 5])
1
>>> pairs_within_budget(100, [1, 2, 3])
3
>>> pairs_within_budget(0, [1, 2])
0
\end{lstlisting}
\noindent\textbf{Solution.}
\begin{lstlisting}[style=code]
def pairs_within_budget(t, xs):
    ys = sorted(xs)
    return sum(max(0, bisect_right(ys, t - x) - (i + 1)) for i, x in enum(ys))
\end{lstlisting}

\section{sensor\_pairing: The screen's domain problem, function form}
\begin{lstlisting}[style=spec]
Two sensor feeds watch the same airspace. A detection is
(id, t, bearing_degrees); ids are sortable strings. Detections a, b
are COMPATIBLE when |a.t - b.t| <= t_tol and their bearing difference
is <= b_tol -- and bearings wrap: 359.5 and 0.3 are 0.8 apart.

Contract:
    pair_detections(a, b, t_tol, b_tol)
        -> (pairs, unpaired_a_ids, unpaired_b_ids)
    Greedy: repeatedly take the lowest-cost compatible pair of still-
    unused detections, cost = |dt| + wrapped bearing diff (exact ties
    by ids). pairs as (a_id, b_id); all three outputs sorted.

Hint:
    Build the compatible candidates as one comprehension over both
    feeds (cost, a_id, b_id, i, j), sort -- the fold that follows
    takes each candidate whose indices are still unused. The wrap is
    min(d, 360 - d) in ONE helper. This is mock_phased.py distilled.
\end{lstlisting}
\noindent\textit{Doctests:}
\begin{lstlisting}[style=spec]
>>> a = [('A0', 10.0, 90.0), ('A1', 20.0, 180.0)]
>>> b = [('B0', 10.2, 91.0), ('B1', 20.1, 179.5)]
>>> pair_detections(a, b, 0.5, 2.0)
([('A0', 'B0'), ('A1', 'B1')], [], [])
>>> pair_detections([('A0', 10.0, 359.5)], [('B0', 10.1, 0.3)], 0.5, 2.0)
([('A0', 'B0')], [], [])
>>> pair_detections([('A0', 0.0, 10.0)], [('B0', 900.0, 200.0)], 0.5, 2.0)
([], ['A0'], ['B0'])
>>> a = [('A0', 10.0, 45.0), ('A1', 30.0, 90.0)]
>>> b = [('B0', 10.5, 44.5), ('B1', 10.6, 46.0), ('B2', 50.0, 200.0)]
>>> pair_detections(a, b, 1.0, 2.0)
([('A0', 'B0')], ['A1'], ['B1', 'B2'])
\end{lstlisting}
\noindent\textbf{Solution.}
\begin{lstlisting}[style=code]
def pair_detections(a, b, t_tol, b_tol):
    wrapd = lambda x, y: min(abs(x - y) % 360, 360 - abs(x - y) % 360)
    cands = sorted([(abs(ta - tb) + wrapd(ba, bb), ai, bi, i, j)
                  for i, (ai, ta, ba) in enum(a)
                  for j, (bi, tb, bb) in enum(b)
                  if abs(ta - tb) <= t_tol and wrapd(ba, bb) <= b_tol])
    used_a, used_b, pairs = set(), set(), []
    for _cost, ai, bi, i, j in cands:
        if i not in used_a and j not in used_b:
            used_a.add(i)
            used_b.add(j)
            pairs.append((ai, bi))
    ua = [d[0] for i, d in enum(a) if i not in used_a]
    ub = [d[0] for j, d in enum(b) if j not in used_b]
    return sorted(pairs), sorted(ua), sorted(ub)
\end{lstlisting}

\noindent\textit{Remark.} The interview capstone, distilled: build every
compatible candidate, sort by cost, greedily take what is still unused. The
wrap-around rule lives in \emph{one} helper, which is why phase mutations cost
one line here. Name the better algorithm (optimal assignment, Hungarian,
$O(n^3)$) and deliberately choose the simpler one -- that sentence is a senior
signal.
\chapter{Class V --- Bags, Dicts and Grouping}

Aggregating by key. \texttt{Counter} counts, \texttt{fromListWith} is the
general dict-building fold, \texttt{unionWith} merges two dicts under any
combining function -- and choosing that function is choosing a monoid
(\texttt{max} gives bag union, \texttt{+} gives a summing merge). The class
also carries the Maybe pipeline (\texttt{first\_unique}, \texttt{log\_report}):
parse to value-or-\texttt{None}, then let the Maybe tools do the bookkeeping.
\medskip\noindent\textit{Practice files live in }\texttt{g5\_bags\_and\_grouping/}\textit{;
each carries the prelude snippets it needs inline. Attempt every problem there
before reading on -- the doctests are the referee.}

\section{valid\_anagram: NeetCode: Valid Anagram}
\begin{lstlisting}[style=spec]
Contract:
    is_anagram(s: str, t: str) -> bool
    True iff t is a permutation of s. Lowercase ASCII.

Hint:
    Two Counter folds; anagram == equal multisets. (Compare with
    g4's sorted-canonical version: same problem, different algebra.)
\end{lstlisting}
\noindent\textit{Doctests:}
\begin{lstlisting}[style=spec]
>>> is_anagram("anagram", "nagaram")
True
>>> is_anagram("rat", "car")
False
>>> is_anagram("", "")
True
\end{lstlisting}
\noindent\textbf{Solution.}
\begin{lstlisting}[style=code]
is_anagram = lambda s, t: Counter(s) == Counter(t)
\end{lstlisting}

\section{ransom\_note: HackerRank: Ransom Note}
\begin{lstlisting}[style=spec]
Contract:
    can_make(note: list[str], magazine: list[str]) -> bool
    True iff every word in note is available in magazine with at
    least the required multiplicity. Case-sensitive whole words.

Hint:
    Multiset inclusion: Counter(note) <= Counter(magazine), spelled
    as an `all` over the note's counts.
\end{lstlisting}
\noindent\textit{Doctests:}
\begin{lstlisting}[style=spec]
>>> can_make("give me one grand today night".split(), "give me one grand today".split())
False
>>> can_make("two times three is not four".split(), "two times three is not four".split())
True
>>> can_make(["give"], ["Give"])
False
\end{lstlisting}
\noindent\textbf{Solution.}
\begin{lstlisting}[style=code]
def can_make(note, magazine):
    have = Counter(magazine)
    need = Counter(note)
    return all(have[w] >= n for w, n in need.items())
\end{lstlisting}

\section{most\_frequent: Most frequent value; ties -> smallest}
\begin{lstlisting}[style=spec]
Contract:
    most_frequent(xs: list) -> value

Hint:
    Counter, then minOn the tuple key (-count, value). The
    two-level tie-break lives entirely in the key, and minOn is
    O(n) where sortOn + head pays O(n log n) for the same answer.
\end{lstlisting}
\noindent\textit{Doctests:}
\begin{lstlisting}[style=spec]
>>> most_frequent([1, 2, 2, 3, 3])
2
>>> most_frequent([5, 5, 1])
5
>>> most_frequent([3, 1, 2])
1
>>> most_frequent(["b", "a", "b", "a"])
'a'
\end{lstlisting}
\noindent\textbf{Solution.}
\begin{lstlisting}[style=code]
most_frequent = lambda xs: minOn(
    lambda kv: (-kv[1], kv[0]), list(Counter(xs).items()))[0]
\end{lstlisting}

\section{top\_k\_words: The k most frequent words}
\begin{lstlisting}[style=spec]
Ordered by count descending, ties by word ascending. At least k
distinct words guaranteed.

Contract:
    top_k_words(k: int, text: str) -> list[tuple[str, int]]

Hint:
    Counter over the words, sortOn (-count, word), take k. The whole
    pipeline is three prelude names deep.
\end{lstlisting}
\noindent\textit{Doctests:}
\begin{lstlisting}[style=spec]
>>> top_k_words(2, "the cat the dog the cat bird")
[('the', 3), ('cat', 2)]
>>> top_k_words(3, "b a a b c c")
[('a', 2), ('b', 2), ('c', 2)]
\end{lstlisting}
\noindent\textbf{Solution.}
\begin{lstlisting}[style=code]
top_k_words = lambda k, text: take(
    k, sortOn(lambda kv: (-kv[1], kv[0]), list(Counter(text.split()).items())))
\end{lstlisting}

\section{group\_words: Group words by first letter}
\begin{lstlisting}[style=spec]
First-appearance order for both keys and members.

Contract:
    group_words(words: list[str]) -> dict[str, list[str]]

Hint:
    THE fromListWith showcase: pairs (first letter, [word]). The
    combiner receives f(new, old) -- Haskell's order -- so KEEP
    arrival order by combining old + new; a bare (+) would reverse
    each group (the classic Haskell gotcha). Keys keep first-seen
    order.
\end{lstlisting}
\noindent\textit{Doctests:}
\begin{lstlisting}[style=spec]
>>> group_words(["ant", "bee", "ape"])
{'a': ['ant', 'ape'], 'b': ['bee']}
>>> group_words(["cat"])
{'c': ['cat']}
>>> group_words([])
{}
\end{lstlisting}
\noindent\textbf{Solution.}
\begin{lstlisting}[style=code]
group_words = lambda ws: fromListWith(
    lambda new, old: old + new, [(w[0], [w]) for w in ws])
\end{lstlisting}

\section{group\_anagrams: NeetCode: Group Anagrams}
\begin{lstlisting}[style=spec]
Contract:
    group_anagrams(words: list[str]) -> list[list[str]]
    Partition words into anagram classes; classes in first-seen order,
    members in input order.

Hint:
    Key = canonical form (sorted letters). Two prelude routes, do
    both: Tree(1, list) with a bare append (autovivification aids the
    write-heavy access), or one fromListWith over (key, [w]) pairs --
    combining old + new, since f receives (new, old); a bare (+)
    would reverse each group. Compare.
\end{lstlisting}
\noindent\textit{Doctests:}
\begin{lstlisting}[style=spec]
>>> group_anagrams(["eat", "tea", "tan", "ate", "nat", "bat"])
[['eat', 'tea', 'ate'], ['tan', 'nat'], ['bat']]
>>> group_anagrams([""])
[['']]
>>> group_anagrams(["a"])
[['a']]
\end{lstlisting}
\noindent\textbf{Solution.}
\begin{lstlisting}[style=code]
def group_anagrams(words):
    return list(fromListWith(lambda new, old: old + new,
                             [("".join(sorted(w)), [w]) for w in words]).values())
\end{lstlisting}

\noindent\textit{Remark.} Two respectable routes: \texttt{Tree(1, list)} with a
bare append (autovivification suits write-heavy access) or one
\texttt{fromListWith} over \texttt{(key, [w])} pairs, combining old + new
(the combiner hears the new value first). Solve it both ways once; afterwards
you will see every grouping problem as a one-liner.

\section{merge\_sum: Merge two dicts of numbers, summing shared keys}
\begin{lstlisting}[style=spec]
Don't mutate the inputs. Key order: a's keys, then b's new ones.

Contract:
    merge_sum(a: dict, b: dict) -> dict

Hint:
    This IS unionWith with (+): merge, combining shared keys by
    addition. The prelude's bag_union is the SAME merge with max --
    one combinator, different monoid. Spot the difference and say it.
    (fromListWith over the two items() lists concatenated gets there
    too; decide which reads better.)
\end{lstlisting}
\noindent\textit{Doctests:}
\begin{lstlisting}[style=spec]
>>> merge_sum({'a': 1, 'b': 2}, {'b': 3, 'c': 4})
{'a': 1, 'b': 5, 'c': 4}
>>> merge_sum({}, {'x': 1})
{'x': 1}
>>> merge_sum({'x': 1}, {})
{'x': 1}
\end{lstlisting}
\noindent\textbf{Solution.}
\begin{lstlisting}[style=code]
merge_sum = lambda a, b: unionWith(lambda x, y: x + y, a, b)
\end{lstlisting}

\section{rle\_encode: Run-length encode}
\begin{lstlisting}[style=spec]
"aaabbc" -> "a3b2c1"; every run gets a count, even runs of one.

Contract:
    rle_encode(s: str) -> str

Hint:
    The prelude's group gathers runs of CONSECUTIVE equals as bare
    lists. map_ each run to run[0] + str(len(run)), join.
\end{lstlisting}
\noindent\textit{Doctests:}
\begin{lstlisting}[style=spec]
>>> rle_encode("aaabbc")
'a3b2c1'
>>> rle_encode("abc")
'a1b1c1'
>>> rle_encode("aaaa")
'a4'
>>> rle_encode("")
''
\end{lstlisting}
\noindent\textbf{Solution.}
\begin{lstlisting}[style=code]
rle_encode = lambda s: "".join(
    map_(lambda run: run[0] + str(len(run)), group(s)))
\end{lstlisting}

\section{lru\_class: An LRU cache class, dict-only}
\begin{lstlisting}[style=spec]
Capacity at construction; put(key, value); get(key) -> value or None.
Both get and put count as use; puts beyond capacity evict the least
recently used entry.

Contract:
    class LRU(capacity)
        .put(key, value) -> None
        .get(key) -> value | None

Hint:
    The container lesson with no snippet to lean on: dicts remember
    insertion order, so pop + re-insert IS move-to-most-recent, and
    next(iter(d)) is the least recent. That's the whole machine.
\end{lstlisting}
\noindent\textit{Doctests:}
\begin{lstlisting}[style=spec]
>>> c = LRU(2)
>>> c.put('a', 1); c.put('b', 2)
>>> c.get('a')
1
>>> c.put('c', 3)
>>> c.get('b') is None
True
>>> c.get('c')
3
>>> c.get('a')
1
>>> c.put('a', 99); c.get('a')
99
\end{lstlisting}
\noindent\textbf{Solution.}
\begin{lstlisting}[style=code]
class LRU:
    def __init__(self, capacity):
        self.capacity = capacity
        self.data = {}
    def get(self, key):
        if key not in self.data:
            return None
        val = self.data.pop(key)
        self.data[key] = val
        return val
    def put(self, key, value):
        if key in self.data:
            self.data.pop(key)
        elif len(self.data) >= self.capacity:
            self.data.pop(next(iter(self.data)))
        self.data[key] = value
\end{lstlisting}

\section{first\_unique: First non-repeating character}
\begin{lstlisting}[style=spec]
Contract:
    first_unique(s: str) -> str   # the char, or '_' if none

Hint:
    Three prelude names, one line: Counter the string, find the first
    char whose count is 1, fromMaybe the '_' default. Two passes,
    O(n) -- and find stops at the first hit.
\end{lstlisting}
\noindent\textit{Doctests:}
\begin{lstlisting}[style=spec]
>>> first_unique("swiss")
'w'
>>> first_unique("aabb")
'_'
>>> first_unique("")
'_'
\end{lstlisting}
\noindent\textbf{Solution.}
\begin{lstlisting}[style=code]
def first_unique(s):
    counts = Counter(s)
    return fromMaybe('_', find(lambda c: counts[c] == 1, s))
\end{lstlisting}

\noindent\textit{Remark.} Three prelude names, one line: \texttt{Counter} builds the
evidence, \texttt{find} stops at the first witness, \texttt{fromMaybe} supplies
the default. This is what vocabulary buys: the solution is its own
specification.

\section{log\_report: Error report from messy logs}
\begin{lstlisting}[style=spec]
A well-formed line is "<ts> <LEVEL> <service>: <message...>": LEVEL
in {INFO, WARN, ERROR}, service nonempty and glued to its colon,
message may be empty. Blank lines are ignored; any other non-blank
line is malformed.

Contract:
    error_report(loglines: list[str]) -> (report, malformed_count)
    report = [(service, error_count)] for services with >= 1 ERROR,
    sorted by count desc, then name asc.

Hint:
    THE mapMaybe shape: parse(line) -> (level, service) or None.
    The Nones ARE the malformed count; catMaybes + a filter feed
    Counter; sortOn (-count, name) ranks. Messy input, no raw loops.
\end{lstlisting}
\noindent\textit{Doctests:}
\begin{lstlisting}[style=spec]
>>> logs = ["t1 ERROR api: down", "t2 INFO web: ok", "garbage line",
...         "t3 ERROR api: db", "t4 ERROR web: 500", "t5 WARN api: slow",
...         "t6 DEBUG api: no", "t7 ERROR db missing"]
>>> error_report(logs)
([('api', 2), ('web', 1)], 3)
>>> error_report(["t1 INFO web: ok", "half a line"])
([], 1)
>>> error_report([])
([], 0)
\end{lstlisting}
\noindent\textbf{Solution.}
\begin{lstlisting}[style=code]
LEVELS = {"INFO", "WARN", "ERROR"}

def parse(line):                            # (level, service) or None
    parts = line.split(maxsplit=3)
    if (len(parts) < 3 or parts[1] not in LEVELS
            or not parts[2].endswith(":") or parts[2] == ":"):
        return None
    return (parts[1], parts[2][:-1])

def error_report(loglines):
    parsed = map_(parse, filter_(lambda l: l.strip(), loglines))
    malformed = len(filter_(lambda p: p is None, parsed))
    errs = [svc for lv, svc in catMaybes(parsed) if lv == "ERROR"]
    report = sortOn(lambda kv: (-kv[1], kv[0]), list(Counter(errs).items()))
    return report, malformed
\end{lstlisting}

\noindent\textit{Remark.} The messy-input shape: \texttt{parse} returns a value or
\texttt{None}, and everything downstream is bookkeeping over Maybes --
\texttt{catMaybes} for the survivors, a filter for the casualty count,
\texttt{Counter} and \texttt{sortOn} for the report. No raw loops, no flags.

\section{invert\_dict: Swap a dict's keys and values}
\begin{lstlisting}[style=spec]
Values are unique (so the inversion is lossless).

Contract:
    invert(d: dict) -> dict

Hint:
    A dict is a list of pairs wearing a hat: items(), swap each pair,
    dict() the result. The swap showcase.
\end{lstlisting}
\noindent\textit{Doctests:}
\begin{lstlisting}[style=spec]
>>> invert({'a': 1, 'b': 2})
{1: 'a', 2: 'b'}
>>> invert({})
{}
\end{lstlisting}
\noindent\textbf{Solution.}
\begin{lstlisting}[style=code]
invert = lambda d: dict(map_(swap, list(d.items())))
\end{lstlisting}
\chapter{Class VI --- Tries and Paths}

Hierarchy as nested dicts. \texttt{ITree} autovivifies branches on touch;
\texttt{setpath}/\texttt{getpath} write and read key paths; \texttt{paths}
folds a whole tree back into a flat list of (keypath, leaf). The subset-sum
pair is the deep exhibit: the same search recorded first as a decision tree
(a multiset of solutions) and then as a trie (a set) -- the data structure
changed the semantics.
\medskip\noindent\textit{Practice files live in }\texttt{g6\_tries\_and\_paths/}\textit{;
each carries the prelude snippets it needs inline. Attempt every problem there
before reading on -- the doctests are the referee.}

\section{deepest\_directory: Challenge: deepest path in a directory forest}
\begin{lstlisting}[style=spec]
Contract:
    deepest(dirs: list[str]) -> str
    Each string is a '/'-separated path ("a/b/c"). Return the deepest
    full path (unique deepest guaranteed). Paths may share prefixes.

Hint:
    Building IS walking: foldl indexing over the parts autovivifies the
    whole branch -- no setpath needed since there is no leaf value.
    Then one paths() + maxOn keypath length. Note why setpath would
    be WRONG here: writing "a/b" after "a/b/c" would clobber the
    subtree.
\end{lstlisting}
\noindent\textit{Doctests:}
\begin{lstlisting}[style=spec]
>>> deepest(["a/b/c", "a/d", "e"])
'a/b/c'
>>> deepest(["x"])
'x'
>>> deepest(["a/b", "a/b/c/d", "a/b/c"])
'a/b/c/d'
\end{lstlisting}
\noindent\textbf{Solution.}
\begin{lstlisting}[style=code]
def deepest(dirs):
    t = ITree()
    for d in dirs:
        foldl(lambda node, k: node[k], d.split('/'), t)
    return '/'.join(maxOn(lambda pl: len(pl[0]), paths(t))[0])
\end{lstlisting}

\noindent\textit{Remark.} Walking \emph{is} building: folding indexing over the
path parts autovivifies the whole branch. Note why \texttt{setpath} would be
wrong -- writing \texttt{a/b} after \texttt{a/b/c} would clobber the subtree.
Choosing the writing primitive is part of the design.

\section{autocomplete: Challenge: prefix trie with completion}
\begin{lstlisting}[style=spec]
Contract:
    make_trie(words: list[str]) -> ITree
        Letter-by-letter trie; each word terminated by '$' whose leaf
        is the word itself.
    complete(trie, prefix: str) -> list[str]
        All stored words starting with prefix, sorted. A word counts
        as its own completion.

Hint:
    setpath writes each word; getpath descends to the prefix's subtree
    WITHOUT autovivifying (reading an ITree with [] would plant ghost
    branches); paths() harvests every '$' leaf below.
    (No-trie baseline: filter_ + isPrefixOf over the word list, O(n*m)
    per query -- the trie is what repeated queries buy.)
\end{lstlisting}
\noindent\textit{Doctests:}
\begin{lstlisting}[style=spec]
>>> make_trie(["at"]) == {'a': {'t': {'$': 'at'}}}
True
>>> t = make_trie(["cat", "car", "card", "dog", "do"])
>>> complete(t, "ca")
['car', 'card', 'cat']
>>> complete(t, "do")
['do', 'dog']
>>> complete(t, "z")
[]
>>> 'z' in t
False
\end{lstlisting}
\noindent\textbf{Solution.}
\begin{lstlisting}[style=code]
def make_trie(words):
    t = ITree()
    for w in words:
        setpath(t, list(w) + ['$'], w)
    return t

def complete(trie, prefix):
    sub = getpath(trie, list(prefix), default={})
    return sorted([leaf for p, leaf in paths(sub) if p and p[-1] == '$'])
\end{lstlisting}

\section{subset\_sum\_paths: Challenge: all paths in the decision tree summing to a total}
\begin{lstlisting}[style=spec]
Contract:
    build_tree(nums: list[int], total: int) -> tree
        The take/skip decision tree over nums. Node = dict with keys
        ('take', x) and ('skip', x); leaf = bool (True iff the takes
        along that path sum to total).
    tree_paths(tree) -> list[list[int]]
        The taken values along every path ending in True, take-first
        DFS order.
    sum_paths(nums: list[int], total: int) -> list[list[int]]
        Every subsequence of nums summing to total. Negatives allowed,
        so a path only ends at full depth -- rem == 0 mid-path proves
        nothing.

Hint:
    An unfold builds the tree from seed (nums, remaining); a fold over
    the tree collects paths ending in True. unfold + fold = hylomorphism:
    fuse them and the tree never exists -- but here it is the exhibit.
    2^n leaves: every subset gets judged, exactly once.
    (The prelude now also ships subsequences -- filtering the powerset
    by sum is the same 2^n with the tree left implicit.)
\end{lstlisting}
\noindent\textit{Doctests:}
\begin{lstlisting}[style=spec]
>>> build_tree([], 0)
True
>>> build_tree([5], 5) == {('take', 5): True, ('skip', 5): False}
True
>>> tree_paths(True)
[[]]
>>> tree_paths(False)
[]
>>> tree_paths(build_tree([1, 2], 3))
[[1, 2]]
>>> sum_paths([1, 2, 3], 3)
[[1, 2], [3]]
>>> sum_paths([2, 4, 6, 10], 16)
[[2, 4, 10], [6, 10]]
>>> sum_paths([1, -1, 2], 2)
[[1, -1, 2], [2]]
>>> sum_paths([], 0)
[[]]
>>> sum_paths([1, 1], 2)
[[1, 1]]
\end{lstlisting}
\noindent\textbf{Solution.}
\begin{lstlisting}[style=code]
def build_tree(nums, total):
    if null(nums):
        return total == 0
    x = head(nums)
    return {('take', x): build_tree(tail(nums), total - x),
            ('skip', x): build_tree(tail(nums), total)}

def tree_paths(tree):
    if tree is True:
        return [[]]
    if tree is False:
        return []
    return concat([[[x] + p if tag == 'take' else p for p in tree_paths(sub)]
                   for (tag, x), sub in tree.items()])

def sum_paths(nums, total):
    return tree_paths(build_tree(nums, total))
\end{lstlisting}

\section{subset\_sum\_trie: Challenge: subset-sum, winners recorded in a solution trie}
\begin{lstlisting}[style=spec]
Contract:
    solution_trie(nums: list[int], total: int) -> ITree
        Search the take/skip space; write only the WINNING paths (the
        taken values, terminated by '$') into an ITree.
    sum_paths(nums: list[int], total: int) -> list[list[int]]
        The winning paths read back off the trie via paths().
        NB: a trie stores a SET of solutions -- duplicates arising from
        different index choices (e.g. [1, 1] with total 1) collapse to one.
        That is a semantic change from 03's multiset.

Hint:
    The full trio: unfold searches, ITree records, paths() reads.
    Why the '$' terminator? One solution can be a prefix of another
    ([2] inside [2, 4, -4]) -- without an end marker the shorter one
    vanishes into the longer one's interior.
\end{lstlisting}
\noindent\textit{Doctests:}
\begin{lstlisting}[style=spec]
>>> solution_trie([1, 2], 3) == {1: {2: {'$': True}}}
True
>>> sum_paths([1, 2, 3], 3)
[[1, 2], [3]]
>>> sum_paths([2, 4, -4], 2)
[[2, 4, -4], [2]]
>>> sum_paths([1, 1], 1)
[[1]]
>>> sum_paths([], 0)
[[]]
\end{lstlisting}
\noindent\textbf{Solution.}
\begin{lstlisting}[style=code]
def solution_trie(nums, total):
    t = ITree()
    def go(i, rem, picked):
        if i == len(nums):
            if rem == 0:
                setpath(t, picked + ['$'], True)
            return
        go(i + 1, rem - nums[i], picked + [nums[i]])
        go(i + 1, rem, picked)
    go(0, total, [])
    return t

def sum_paths(nums, total):
    return [p[:-1] for p, _ in paths(solution_trie(nums, total)) if p[-1] == '$']
\end{lstlisting}

\noindent\textit{Remark.} The \texttt{\$} terminator earns its place: one
solution can be a prefix of another, and without an end marker the shorter one
vanishes into the longer one's interior. Also note the semantic shift from the
decision-tree version: a trie stores a \emph{set} of solutions, so duplicate
index-choices collapse.
\chapter{Class VII --- Graphs, Grids, Windows and Heaps}

Things pointing at other things. Adjacency is a \texttt{Tree(1, list)} fold
over the edges; BFS is the \texttt{deque} earning its keep; Dijkstra adds a
heap and lazy deletion; Kahn's topological sort turns cycle detection into a
counting argument; Bellman--Ford is \texttt{until} iterating to a fixed point.
The window problems ride along because their state discipline (add, remove,
valid) is the same closure-shaped thinking on a 1-D track.
\medskip\noindent\textit{Practice files live in }\texttt{g7\_graphs\_grids\_heaps/}\textit{;
each carries the prelude snippets it needs inline. Attempt every problem there
before reading on -- the doctests are the referee.}

\section{longest\_unique: Longest substring without repeats}
\begin{lstlisting}[style=spec]
Contract:
    longest_unique(s: str) -> int

Hint:
    The longest_window skeleton does the two-pointer bookkeeping; you
    supply the state as closures: per-char counts (Tree(1, int)) plus
    a duplicate tally. valid() == "no duplicates in the window".
\end{lstlisting}
\noindent\textit{Doctests:}
\begin{lstlisting}[style=spec]
>>> longest_unique("abcabcbb")
3
>>> longest_unique("bbbbb")
1
>>> longest_unique("")
0
>>> longest_unique("abba")
2
\end{lstlisting}
\noindent\textbf{Solution.}
\begin{lstlisting}[style=code]
def longest_unique(s):
    counts, dups = Tree(1, int), [0]
    def add(c):
        counts[c] += 1
        dups[0] += counts[c] == 2
    def rem(c):
        counts[c] -= 1
        dups[0] -= counts[c] == 1
    return longest_window(s, lambda: dups[0] == 0, add, rem)
\end{lstlisting}

\noindent\textit{Remark.} The window skeleton does the two-pointer bookkeeping;
you supply state as closures. Writing \texttt{valid}/\texttt{add}/\texttt{rem}
separately is exactly the decomposition that survives a mutating interview
prompt.

\section{grid\_path: Shortest path through a grid}
\begin{lstlisting}[style=spec]
Steps from top-left to bottom-right through '.' cells, around '#'
walls, moving orthogonally. -1 if unreachable (or an endpoint is a
wall); a 1x1 open grid is 0.

Contract:
    grid_path(grid: list[str]) -> int

Hint:
    BFS: the prelude deque is the frontier, neighbors4 the edges.
    Mark seen when you ENQUEUE, not when you pop.
\end{lstlisting}
\noindent\textit{Doctests:}
\begin{lstlisting}[style=spec]
>>> grid_path(["..", ".."])
2
>>> grid_path([".#", "#."])
-1
>>> grid_path(["."])
0
>>> grid_path(["..#", "#.#", "#.."])
4
\end{lstlisting}
\noindent\textbf{Solution.}
\begin{lstlisting}[style=code]
def grid_path(grid):
    R, C = len(grid), len(grid[0])
    if grid[0][0] == "#" or grid[R - 1][C - 1] == "#":
        return -1
    q, seen = deque([(0, 0, 0)]), {(0, 0)}
    while len(q):
        r, c, d = q.popleft()
        if (r, c) == (R - 1, C - 1):
            return d
        for nr, nc in neighbors4(r, c, R, C):
            if grid[nr][nc] == "." and (nr, nc) not in seen:
                seen.add((nr, nc))
                q.append((nr, nc, d + 1))
    return -1
\end{lstlisting}

\section{rate\_limiter: Sliding-window rate limiter}
\begin{lstlisting}[style=spec]
A request at time ts is denied when its client already has `limit`
ALLOWED requests in (ts - window, ts]. Denied requests consume no
budget. Timestamps arrive non-decreasing.

Contract:
    rate_limiter(window: int, limit: int,
                 events: list[tuple[int, str]]) -> list[tuple[int, str]]
    The denied (ts, client) pairs, in stream order.

Hint:
    Tree(1, deque): each client's allowed timestamps in a deque. The
    prelude deque has no peek -- evict by popleft, and appendleft the
    survivor back. Each stamp enters and leaves once: O(1) amortized.
\end{lstlisting}
\noindent\textit{Doctests:}
\begin{lstlisting}[style=spec]
>>> rate_limiter(10, 2, [(1, 'alice'), (2, 'alice'), (3, 'alice'),
...                      (5, 'bob'), (12, 'alice'), (13, 'alice')])
[(3, 'alice')]
>>> rate_limiter(5, 3, [(1, 'a'), (2, 'a'), (3, 'a')])
[]
>>> rate_limiter(5, 1, [(1, 'a'), (2, 'b'), (3, 'a'), (7, 'a'), (8, 'b')])
[(3, 'a')]
\end{lstlisting}
\noindent\textbf{Solution.}
\begin{lstlisting}[style=code]
def rate_limiter(window, limit, events):
    allowed = Tree(1, deque)
    denied = []
    for ts, client in events:
        q = allowed[client]
        while len(q):
            t = q.popleft()
            if t > ts - window:
                q.appendleft(t)
                break
        if len(q) >= limit:
            denied.append((ts, client))
        else:
            q.append(ts)
    return denied
\end{lstlisting}

\section{running\_median: Median after each arrival}
\begin{lstlisting}[style=spec]
Even counts average the two middle values; one decimal is the
caller's problem -- return floats.

Contract:
    running_median(xs: list[int]) -> list[float]

Hint:
    Two heaps: lo holds the lower half NEGATED (the prelude heap is
    min-only), hi the upper half. Push through lo into hi, rebalance
    so lo keeps the extra; the median reads off the tops.
\end{lstlisting}
\noindent\textit{Doctests:}
\begin{lstlisting}[style=spec]
>>> running_median([12, 4, 5, 3, 8, 7])
[12.0, 8.0, 5.0, 4.5, 5.0, 6.0]
>>> running_median([10])
[10.0]
>>> running_median([1, 2, 3, 4])
[1.0, 1.5, 2.0, 2.5]
\end{lstlisting}
\noindent\textbf{Solution.}
\begin{lstlisting}[style=code]
def running_median(xs):
    lo, hi, out = [], [], []
    for x in xs:
        heappush(lo, -x)
        heappush(hi, -heappop(lo))
        if len(hi) > len(lo):
            heappush(lo, -heappop(hi))
        m = -lo[0] if len(lo) > len(hi) else (-lo[0] + hi[0]) / 2
        out.append(float(m))
    return out
\end{lstlisting}

\section{dijkstra: Challenge: shortest paths from weighted (from, to, dist) edges}
\begin{lstlisting}[style=spec]
Contract:
    dijkstra(edges: list[tuple], src) -> (dist: dict, parent: dict)
        Directed edges (u, v, w), w >= 0. dist[n] = shortest distance
        src -> n; unreachable nodes absent. parent[n] = predecessor on
        one shortest path (parent[src] = None).
    shortest_path(edges, src, dst) -> list | None
        The node sequence of a shortest src -> dst path, None if
        unreachable.

Hint:
    Adjacency map = Tree(1, list) fold over edges. Frontier = the heap.
    Lazy deletion: skip a popped entry unless it matches the current
    best -- the heap may hold stale distances. Path reconstruction =
    an unfold: iterate the parent chain, takeWhile not None, reverse.
    Graphs are not trees -- cycles make tree unfolds diverge, so fold
    over STATES to a fixed point instead.
\end{lstlisting}
\noindent\textit{Doctests:}
\begin{lstlisting}[style=spec]
>>> dist, _ = dijkstra([('a','b',1), ('b','c',2), ('a','c',4)], 'a')
>>> dist == {'a': 0, 'b': 1, 'c': 3}
True
>>> dist, _ = dijkstra([('a','b',5), ('a','c',1), ('c','b',1)], 'a')
>>> dist['b']
2
>>> dijkstra([('x','y',1)], 'y')[0]
{'y': 0}
>>> shortest_path([('a','b',5), ('a','c',1), ('c','b',1)], 'a', 'b')
['a', 'c', 'b']
>>> shortest_path([('a','b',1)], 'a', 'a')
['a']
>>> shortest_path([('a','b',1)], 'b', 'a') is None
True
\end{lstlisting}
\noindent\textbf{Solution.}
\begin{lstlisting}[style=code]
INF = float("inf")

def dijkstra(edges, src):
    adj = Tree(1, list)
    for u, v, w in edges:
        adj[u].append((w, v))
    dist, parent, h = {src: 0}, {src: None}, [(0, src)]
    while h:
        d, u = heappop(h)
        if d == dist.get(u, INF):
            for w, v in adj[u]:
                if d + w < dist.get(v, INF):
                    dist[v] = d + w
                    parent[v] = u
                    heappush(h, (d + w, v))
    return dist, parent

def shortest_path(edges, src, dst):
    dist, parent = dijkstra(edges, src)
    if dst not in dist:
        return None
    back = takeWhile(lambda n: n is not None, iterate(lambda n: parent[n], dst))
    return list(reversed(back))
\end{lstlisting}

\noindent\textit{Remark.} Lazy deletion is the grown-up move: the heap may hold stale
distances, so skip a popped entry unless it matches the current best. Path
reconstruction is an unfold along the parent chain. Graphs are not trees --
cycles make tree unfolds diverge, so we fold over \emph{states} to a fixed
point.

\section{dag\_longest\_path: Challenge: longest path in a weighted DAG}
\begin{lstlisting}[style=spec]
Contract:
    dag_longest(edges: list[tuple]) -> (total: int, path: list)
    Directed acyclic edges (u, v, w). The maximum-weight path anywhere
    in the DAG and its node sequence. Ties broken toward the
    lexically-first end node. Input guaranteed acyclic.

Hint:
    Longest path is NP-hard on general graphs; acyclicity is what
    makes it linear. Kahn's topo sort (indegree fold + deque), then a
    fold over topo order relaxing best[v] = max(best[v], best[u] + w).
    Same skeleton as Dijkstra's relax with max for min and topo order
    replacing the heap -- the LCS-vs-edit-distance move on graphs.
    maxOn best picks the end node; feed it the SORTED nodes so
    ties break lexically.
\end{lstlisting}
\noindent\textit{Doctests:}
\begin{lstlisting}[style=spec]
>>> dag_longest([('a','b',3), ('b','c',4), ('a','c',5)])
(7, ['a', 'b', 'c'])
>>> dag_longest([('s','a',1), ('s','b',5), ('a','t',6), ('b','t',1)])
(7, ['s', 'a', 't'])
>>> dag_longest([('u','v',2)])
(2, ['u', 'v'])
\end{lstlisting}
\noindent\textbf{Solution.}
\begin{lstlisting}[style=code]
def dag_longest(edges):
    adj, indeg, nodes = Tree(1, list), Tree(1, int), set()
    for u, v, w in edges:
        adj[u].append((w, v))
        indeg[v] += 1
        nodes |= {u, v}
    q = deque([n for n in sorted(nodes) if indeg[n] == 0])
    topo = []
    while len(q):
        u = q.popleft()
        topo.append(u)
        for w, v in adj[u]:
            indeg[v] -= 1
            if indeg[v] == 0:
                q.append(v)
    best = {n: 0 for n in nodes}
    parent = {n: None for n in nodes}
    for u in topo:
        for w, v in adj[u]:
            if best[u] + w > best[v]:
                best[v] = best[u] + w
                parent[v] = u
    end = maxOn(lambda n: best[n], sorted(nodes))
    back = takeWhile(lambda n: n is not None, iterate(lambda n: parent[n], end))
    return best[end], list(reversed(back))
\end{lstlisting}

\section{bellman\_ford: Challenge: shortest paths as a fixed point (negative weights ok)}
\begin{lstlisting}[style=spec]
Contract:
    bellman_ford(edges: list[tuple], src) -> dict | None
    Directed edges (u, v, w); w may be negative. dist[n] = shortest
    distance src -> n, unreachable nodes absent. Returns None if a
    negative cycle is reachable from src (distances then have no
    fixed point -- they diverge to -inf).

Hint:
    Shortest paths ARE a fixed point: dist is the least map satisfying
    dist[v] <= dist[u] + w for every edge. So: relax ALL edges as one
    step function, `until` the map stops changing. Convergence within
    |V| - 1 steps is guaranteed (a shortest path has at most |V| - 1
    edges); a change on step |V| proves a negative cycle. Dijkstra is
    the same relax with a clever ORDER; Bellman-Ford orders nothing
    and pays for it in passes.
\end{lstlisting}
\noindent\textit{Doctests:}
\begin{lstlisting}[style=spec]
>>> bellman_ford([('a','b',1), ('b','c',2), ('a','c',4)], 'a') == {'a': 0, 'b': 1, 'c': 3}
True
>>> bellman_ford([('a','b',5), ('b','c',-4), ('a','c',2)], 'a')['c']
1
>>> bellman_ford([('a','b',1), ('b','c',-3), ('c','b',1), ('b','d',1)], 'a') is None
True
>>> bellman_ford([('x','y',7)], 'y')
{'y': 0}
>>> bellman_ford([('a','b',-2), ('b','a',3)], 'a') == {'a': 0, 'b': -2}
True
\end{lstlisting}
\noindent\textbf{Solution.}
\begin{lstlisting}[style=code]
INF = float("inf")

def bellman_ford(edges, src):
    n_nodes = len({n for u, v, _ in edges for n in (u, v)} | {src})

    def relax(state):
        steps, dist = state
        new = dict(dist)
        for u, v, w in edges:
            if new.get(u, INF) + w < new.get(v, INF):
                new[v] = new[u] + w
        return (steps + 1, new)

    def settled(state):
        steps, dist = state
        return steps >= n_nodes or relax(state)[1] == dist

    steps, dist = until(settled, relax, (0, {src: 0}))
    return None if relax((steps, dist))[1] != dist else dist
\end{lstlisting}

\noindent\textit{Remark.} Shortest paths \emph{are} a fixed point: the least map
satisfying every edge's triangle inequality. \texttt{until} iterates the relax
step to stability; a change on pass $|V|$ proves a negative cycle. Dijkstra is
the same relax with a clever order; Bellman--Ford orders nothing and pays in
passes.

\section{course\_schedule: NeetCode: Course Schedule (cycle detection via Kahn)}
\begin{lstlisting}[style=spec]
Contract:
    can_finish(n: int, prereqs: list[tuple]) -> bool
        Courses 0..n-1; (a, b) means b must precede a. True iff every
        course can be taken, i.e. the prerequisite graph is acyclic.
    find_order(n: int, prereqs: list[tuple]) -> list | None
        A valid course order (smallest-numbered course first on ties),
        or None if impossible.

Hint:
    Kahn's insight inverted: topo sort doesn't FAIL on a cycle, it
    just leaves the cycle behind -- nodes on a cycle never reach
    indegree 0. So: run 06's Kahn, then len(topo) == n is the whole
    cycle test. Heap instead of deque only to make ties deterministic.
\end{lstlisting}
\noindent\textit{Doctests:}
\begin{lstlisting}[style=spec]
>>> find_order(2, [(1, 0)])
[0, 1]
>>> find_order(4, [(1, 0), (2, 0), (3, 1), (3, 2)])
[0, 1, 2, 3]
>>> find_order(2, [(1, 0), (0, 1)]) is None
True
>>> find_order(3, [])
[0, 1, 2]
>>> can_finish(2, [(1, 0)])
True
>>> can_finish(2, [(1, 0), (0, 1)])
False
>>> can_finish(5, [(1, 4), (2, 4), (3, 1), (3, 2)])
True
\end{lstlisting}
\noindent\textbf{Solution.}
\begin{lstlisting}[style=code]
def find_order(n, prereqs):
    adj, indeg = Tree(1, list), Tree(1, int)
    for a, b in prereqs:
        adj[b].append(a)
        indeg[a] += 1
    heap = []
    for c in range(n):
        if indeg[c] == 0:
            heappush(heap, c)
    topo = []
    while heap:
        u = heappop(heap)
        topo.append(u)
        for v in adj[u]:
            indeg[v] -= 1
            if indeg[v] == 0:
                heappush(heap, v)
    return topo if len(topo) == n else None

def can_finish(n, prereqs):
    return find_order(n, prereqs) is not None
\end{lstlisting}
\chapter{Class VIII --- Dynamic Programming and Control}

Choices now change what is possible later -- so memoise the future's needs.
The craft is naming the \emph{state}: index pairs for alignments
(\texttt{lcs}, \texttt{edit\_distance}), remaining capacity for knapsacks, a
rectangle for guillotine cuts, a bitmask plus room for the partition problem.
If you can say ``the future only needs $X$'' in one sentence, \texttt{memo}
over $X$ writes the rest. The parser closes the class: recursive descent
is control flow as grammar.
\medskip\noindent\textit{Practice files live in }\texttt{g8\_dp\_and\_control/}\textit{;
each carries the prelude snippets it needs inline. Attempt every problem there
before reading on -- the doctests are the referee.}

\section{coin\_change: NeetCode: Coin Change}
\begin{lstlisting}[style=spec]
Contract:
    coin_change(coins: list[int], amount: int) -> int
    Fewest coins summing to amount (unlimited supply), -1 if impossible.
    amount >= 0, coins positive.

Hint:
    Knapsack, dp[remaining]. Top-down with memo on the single
    index `amount` -- the future only needs what remains.
\end{lstlisting}
\noindent\textit{Doctests:}
\begin{lstlisting}[style=spec]
>>> coin_change([1, 2, 5], 11)
3
>>> coin_change([2], 3)
-1
>>> coin_change([1], 0)
0
>>> coin_change([186, 419, 83, 408], 6249)
20
\end{lstlisting}
\noindent\textbf{Solution.}
\begin{lstlisting}[style=code]
INF = float("inf")

def coin_change(coins, amount):
    @memo
    def go(rem):
        if rem == 0:
            return 0
        return min([1 + go(rem - c) for c in coins if c <= rem],
                       default=INF)
    best = go(amount)
    return -1 if best == INF else best
\end{lstlisting}

\section{lcs: NeetCode: Longest Common Subsequence}
\begin{lstlisting}[style=spec]
Contract:
    lcs(a: str, b: str) -> int
    Length of the longest common subsequence (not substring).

Hint:
    Alignment pattern, dp[i][j] -- same skeleton as 03's edit
    distance, but max instead of min. Index-based args for the cache.
\end{lstlisting}
\noindent\textit{Doctests:}
\begin{lstlisting}[style=spec]
>>> lcs("abcde", "ace")
3
>>> lcs("abc", "abc")
3
>>> lcs("abc", "def")
0
>>> lcs("", "abc")
0
\end{lstlisting}
\noindent\textbf{Solution.}
\begin{lstlisting}[style=code]
def lcs(a, b):
    @memo
    def go(i, j):
        if i == len(a) or j == len(b):
            return 0
        if a[i] == b[j]:
            return 1 + go(i + 1, j + 1)
        return max(go(i + 1, j), go(i, j + 1))
    return go(0, 0)
\end{lstlisting}

\section{edit\_distance: Levenshtein distance}
\begin{lstlisting}[style=spec]
Minimum single-character insertions, deletions, or substitutions
turning a into b.

Contract:
    edit_distance(a: str, b: str) -> int

Hint:
    The alignment DP again -- 02's skeleton with min-of-three plus
    one. Matching heads cost nothing; otherwise pay 1 and recurse on
    the three repairs. memo over the index pair.
\end{lstlisting}
\noindent\textit{Doctests:}
\begin{lstlisting}[style=spec]
>>> edit_distance("kitten", "sitting")
3
>>> edit_distance("", "abc")
3
>>> edit_distance("same", "same")
0
>>> edit_distance("flaw", "lawn")
2
>>> edit_distance("a", "b")
1
\end{lstlisting}
\noindent\textbf{Solution.}
\begin{lstlisting}[style=code]
def edit_distance(a, b):
    @memo
    def go(i, j):
        if i == len(a):
            return len(b) - j
        if j == len(b):
            return len(a) - i
        if a[i] == b[j]:
            return go(i + 1, j + 1)
        return 1 + min(go(i + 1, j), go(i, j + 1), go(i + 1, j + 1))
    return go(0, 0)
\end{lstlisting}

\section{knapsack: 0/1 knapsack}
\begin{lstlisting}[style=spec]
Each (weight, value) item usable at most once; maximize value within
capacity.

Contract:
    knapsack(capacity: int, items: list[tuple[int, int]]) -> int

Hint:
    memo over (i, room): take item i or skip it. Two integers of
    state -- the cache keys on the args tuple, so keep them hashable.
\end{lstlisting}
\noindent\textit{Doctests:}
\begin{lstlisting}[style=spec]
>>> knapsack(10, [(5, 10), (4, 40), (6, 30), (3, 50)])
90
>>> knapsack(0, [(1, 100)])
0
>>> knapsack(5, [])
0
>>> knapsack(3, [(4, 100), (2, 7)])
7
\end{lstlisting}
\noindent\textbf{Solution.}
\begin{lstlisting}[style=code]
def knapsack(capacity, items):
    items = list(items)
    @memo
    def go(i, room):
        if i == len(items) or room == 0:
            return 0
        w, v = items[i]
        skip = go(i + 1, room)
        if w > room:
            return skip
        return max(skip, v + go(i + 1, room - w))
    return go(0, capacity)
\end{lstlisting}

\section{calculator: Integer expression evaluator (capstone)}
\begin{lstlisting}[style=spec]
+, -, * and parentheses, optional spaces, normal precedence,
left-associative. No unary minus, no division.

Contract:
    evaluate(expr: str) -> int

Grammar (one function per rule, shared position cursor):
    expression := term (('+' | '-') term)*
    term       := factor ('*' factor)*
    factor     := NUMBER | '(' expression ')'

Hint:
    Recursive descent. Lex numbers with takeWhile str.isdigit on the
    remaining input -- span IS the lexer. And remember the trap you
    certified earlier: '' in "+-" is True; test membership against a
    tuple.
\end{lstlisting}
\noindent\textit{Doctests:}
\begin{lstlisting}[style=spec]
>>> evaluate("2*(3+4)")
14
>>> evaluate("2*(3+4)-5")
9
>>> evaluate("10-2-3")
5
>>> evaluate("2+3*4")
14
>>> evaluate("((1+2))*3")
9
>>> evaluate(" 12 * 2 ")
24
\end{lstlisting}
\noindent\textbf{Solution.}
\begin{lstlisting}[style=code]
def evaluate(expr):
    s = expr.replace(" ", "")
    pos = [0]

    def peek():
        return s[pos[0]] if pos[0] < len(s) else ""

    def number():
        digits = takeWhile(str.isdigit, s[pos[0]:])
        pos[0] += len(digits)
        return int("".join(digits))

    def factor():
        if peek() == "(":
            pos[0] += 1
            v = expression()
            pos[0] += 1
            return v
        return number()

    def term():
        v = factor()
        while peek() == "*":
            pos[0] += 1
            v *= factor()
        return v

    def expression():
        v = term()
        while peek() in ("+", "-"):
            op = peek()
            pos[0] += 1
            v = v + term() if op == "+" else v - term()
        return v

    return expression()
\end{lstlisting}

\section{max\_height\_cuboids: LC 1691: Maximum Height by Stacking Cuboids}
\begin{lstlisting}[style=spec]
Any rotation allowed; cuboid A sits on B only when ALL THREE of A's
dims are <= B's, compared side to side. Each cuboid used at most
once. Maximize stack height.

Contract:
    max_height(cuboids: list[tuple[int, int, int]]) -> int

Hint:
    The rotation freedom COLLAPSES: sort each cuboid's own dims
    ascending -- the biggest dim as height is always at least as good,
    and sorted-vs-sorted is the only comparison that can succeed (the
    photo_lineup exchange argument, twice). Then sort the cuboids and
    run weighted LIS: best ending at i = h_i + best over j fitting
    under i. Fold-with-history, O(n^2) -- the g4_07 shape with sum
    instead of length. (g8_06's mask version is the ground truth
    oracle for THIS file: same >= 17-style surprises, different fit
    rule -- all three dims here, base only there.)
\end{lstlisting}
\noindent\textit{Doctests:}
\begin{lstlisting}[style=spec]
>>> max_height([(50, 45, 20), (95, 37, 53), (45, 23, 12)])
190
>>> max_height([(38, 25, 45), (76, 35, 3)])
76
>>> max_height([(7, 11, 17), (7, 17, 11), (11, 7, 17), (11, 17, 7), (17, 7, 11), (17, 11, 7)])
102
>>> max_height([(1, 1, 1)])
1
\end{lstlisting}
\noindent\textbf{Solution.}
\begin{lstlisting}[style=code]
def max_height(cuboids):
    cs = sorted([sorted(c) for c in cuboids])
    fits = lambda a, b: all(zipWith(lambda x, y: x <= y, a, b))   # a under-or-equal b, all three dims
    best = []
    for i, c in enumerate(cs):
        best.append(c[2] + max([best[j] for j in range(i) if fits(cs[j], c)],
                               default=0))
    return max(best, default=0)
\end{lstlisting}

\section{guillotine\_cut: Max value cut from a sheet, guillotine cuts only}
\begin{lstlisting}[style=spec]
A W x H sheet; piece types (pw, ph, value), unlimited supply, fixed
orientation (add rotated copies yourself if rotation is allowed).
Every cut runs edge to edge, so each cut splits one rectangle into
two rectangles. Unused offcuts score 0.

Contract:
    guillotine(W: int, H: int, pieces: list[tuple[int, int, int]]) -> int

Hint:
    Guillotine is what re-admits DP: free-form 2D packing has no
    small state (the frontier goes irregular), but edge-to-edge cuts
    keep every subproblem a RECTANGLE -- (w, h), two ints. Interval
    pattern: best(w,h) = max of a piece that fits the sheet EXACTLY,
    every vertical split x, every horizontal split y. A smaller piece
    never needs a direct case: the cut that isolates it is one of the
    splits. memo over (w, h); O(W*H*(W+H)) -- say it.
\end{lstlisting}
\noindent\textit{Doctests:}
\begin{lstlisting}[style=spec]
>>> guillotine(4, 4, [(1, 1, 1)])
16
>>> guillotine(4, 4, [(3, 3, 7), (1, 1, 1)])
16
>>> guillotine(4, 4, [(3, 3, 20), (1, 1, 1)])
27
>>> guillotine(4, 4, [(4, 2, 10), (2, 2, 3)])
20
>>> guillotine(3, 3, [(2, 2, 5)])
5
>>> guillotine(2, 2, [(3, 1, 99)])
0
\end{lstlisting}
\noindent\textbf{Solution.}
\begin{lstlisting}[style=code]
def guillotine(W, H, pieces):
    pieces = list(pieces)

    @memo
    def best(w, h):
        exact  = [v for pw, ph, v in pieces if (pw, ph) == (w, h)]
        vsplit = [best(x, h) + best(w - x, h) for x in range(1, w // 2 + 1)]
        hsplit = [best(w, y) + best(w, h - y) for y in range(1, h // 2 + 1)]
        return max(exact + vsplit + hsplit, default=0)

    return best(W, H)
\end{lstlisting}

\noindent\textit{Remark.} Guillotine cuts are what re-admit dynamic programming:
free-form 2-D packing has no small state, but edge-to-edge cuts keep every
subproblem a rectangle -- two integers. A smaller piece never needs a direct
case; the cut that isolates it is one of the splits.

\section{partition\_k\_subsets: LC 698: Partition to K Equal Sum Subsets}
\begin{lstlisting}[style=spec]
Contract:
    can_partition(nums: list[int], k: int) -> bool
    True iff nums splits into k non-empty groups of equal sum.
    Positive ints, small n (say the bound: 2^n states).

Hint:
    The mask family, pure form: the future needs (which numbers
    remain, how much room is left in the group being filled). When a
    group closes (room hits 0) the room resets to target -- so the
    state is just (used, room): the number of CLOSED groups is
    implied by the used-sum. memo over that pair; fill greedily
    into one group at a time to avoid counting group orderings.
    (g8_06's go is this skeleton plus geometry.)
\end{lstlisting}
\noindent\textit{Doctests:}
\begin{lstlisting}[style=spec]
>>> can_partition([4, 3, 2, 3, 5, 2, 1], 4)
True
>>> can_partition([1, 2, 3, 4], 3)
False
>>> can_partition([2, 2, 2, 2], 2)
True
>>> can_partition([1], 1)
True
>>> can_partition([1, 1], 3)
False
\end{lstlisting}
\noindent\textbf{Solution.}
\begin{lstlisting}[style=code]
def can_partition(nums, k):
    nums = list(nums)
    total, n = sum(nums), len(nums)
    if k <= 0 or n < k or total % k:
        return False
    target = total // k

    @memo
    def go(used, room):
        if used == (1 << n) - 1:
            return room == target            # last group closed cleanly
        return any(go(used | (1 << i),
                      target if room == nums[i] else room - nums[i])
                   for i in range(n)
                   if not used & (1 << i) and nums[i] <= room)

    return go(0, target)
\end{lstlisting}

\noindent\textit{Remark.} The mask family in pure form: the future needs only
(which numbers remain, room left in the open group), because a closed group's
room resets to the target. Fill one group at a time so orderings are not counted
twice. $2^n$ states -- say the bound before you code.

\section{min\_path\_sum: Cheapest right/down path through a grid}
\begin{lstlisting}[style=spec]
Contract:
    min_path_sum(grid: list[list[int]]) -> int
    Top-left to bottom-right, moving only right or down.

Hint:
    Grid DP: state (r, c) -- the future only cares where you stand.
    memo over the index pair; base case is the goal cell.
    Compare g7_02: uniform steps there needed BFS; weighted cells
    with only right/down moves make a DAG, so plain DP suffices.
\end{lstlisting}
\noindent\textit{Doctests:}
\begin{lstlisting}[style=spec]
>>> min_path_sum([[1, 3, 1], [1, 5, 1], [4, 2, 1]])
7
>>> min_path_sum([[5]])
5
>>> min_path_sum([[1, 2], [3, 4]])
7
\end{lstlisting}
\noindent\textbf{Solution.}
\begin{lstlisting}[style=code]
def min_path_sum(grid):
    R, C = len(grid), len(grid[0])

    @memo
    def go(r, c):
        if (r, c) == (R - 1, C - 1):
            return grid[r][c]
        opts = ([go(r + 1, c)] if r + 1 < R else []) + \
               ([go(r, c + 1)] if c + 1 < C else [])
        return grid[r][c] + min(opts)

    return go(0, 0)
\end{lstlisting}
\chapter{Class IX --- Stack Folds}

The prelude has no stack section because a plain Python list already is one:
\texttt{append} pushes, \texttt{pop} pops, both $O(1)$. What this class adds is
the \emph{fold} discipline: the stack is the accumulator, and in
\texttt{balanced\_brackets} the failed state is \texttt{None} -- a poison value
that propagates untouched, Maybe as an accumulator. \texttt{next\_greater}
graduates to the monotonic stack, the pro tier, with its lovely amortised
$O(n)$ argument.
\medskip\noindent\textit{Practice files live in }\texttt{g9\_stack\_folds/}\textit{;
each carries the prelude snippets it needs inline. Attempt every problem there
before reading on -- the doctests are the referee.}

\section{balanced\_brackets: Matching pairs: valid bracket nesting}
\begin{lstlisting}[style=spec]
Only ()[]{} appear.

Contract:
    balanced(s: str) -> bool

Hint:
    The stack fold, purest form: a plain list IS a stack, and the
    fold's accumulator is one. The twist worth learning: use None as
    the FAILED state -- once a closer mismatches, the poison
    propagates through the rest of the fold untouched (Maybe as
    accumulator). Balanced == the fold ends on an empty stack.
\end{lstlisting}
\noindent\textit{Doctests:}
\begin{lstlisting}[style=spec]
>>> balanced("([{}])")
True
>>> balanced("([)]")
False
>>> balanced("(((")
False
>>> balanced("")
True
>>> balanced("]")
False
\end{lstlisting}
\noindent\textbf{Solution.}
\begin{lstlisting}[style=code]
MATCH = {")": "(", "]": "[", "}": "{"}

def balanced(s):
    def step(stack, c):
        if stack is None:
            return None                     # poison: already failed
        if c in "([{":
            return stack + [c]
        if not stack or stack[-1] != MATCH[c]:
            return None
        return stack[:-1]
    return foldl(step, s, []) == []
\end{lstlisting}

\noindent\textit{Remark.} A fold whose accumulator is a stack -- and whose
failure state is \texttt{None}, propagating untouched through the rest of the
fold. That is Maybe as an accumulator, and it removes every flag and early
return from the classic solution.

\section{eval\_rpn: Evaluate Reverse Polish notation}
\begin{lstlisting}[style=spec]
Integer operands, operators + - *.

Contract:
    eval_rpn(tokens: list[str]) -> int

Hint:
    One fold over the tokens: operands push; an operator pops two and
    pushes the result. The classic slip: stack[-2] is the LEFT
    operand. The answer is whatever the fold leaves on the stack.
\end{lstlisting}
\noindent\textit{Doctests:}
\begin{lstlisting}[style=spec]
>>> eval_rpn(["2", "3", "+", "4", "*"])
20
>>> eval_rpn(["5", "1", "2", "+", "4", "*", "-"])
-7
>>> eval_rpn(["7"])
7
\end{lstlisting}
\noindent\textbf{Solution.}
\begin{lstlisting}[style=code]
OPS = {"+": lambda a, b: a + b, "-": lambda a, b: a - b,
       "*": lambda a, b: a * b}

def eval_rpn(tokens):
    def step(stack, tok):
        if tok in OPS:
            return stack[:-2] + [OPS[tok](stack[-2], stack[-1])]
        return stack + [int(tok)]
    return foldl(step, tokens, [])[0]
\end{lstlisting}

\section{next\_greater: First later element that's larger}
\begin{lstlisting}[style=spec]
For each element, the next strictly greater element to its right;
-1 if none.

Contract:
    next_greater(xs: list[int]) -> list[int]

Hint:
    The MONOTONIC stack: fold over enum(xs); the newcomer pops every
    stacked index whose value it beats (that newcomer is their
    answer), then pushes itself. Each index pushes and pops at most
    once -- O(n) despite the inner while. The pro tier of stacks.
\end{lstlisting}
\noindent\textit{Doctests:}
\begin{lstlisting}[style=spec]
>>> next_greater([2, 1, 2, 4, 3])
[4, 2, 4, -1, -1]
>>> next_greater([5, 4, 3])
[-1, -1, -1]
>>> next_greater([])
[]
\end{lstlisting}
\noindent\textbf{Solution.}
\begin{lstlisting}[style=code]
def next_greater(xs):
    out = [-1] * len(xs)
    def step(stack, ix):
        i, x = ix
        while stack and xs[stack[-1]] < x:
            out[stack.pop()] = x
        return stack + [i]
    foldl(step, enum(xs), [])
    return out
\end{lstlisting}

\noindent\textit{Remark.} The monotonic stack: the newcomer pops every index it
beats (it is their answer) and then waits for its own. Each index pushes and
pops at most once, so the nested \texttt{while} is still $O(n)$ -- amortised
analysis in one picture.
\appendix
\part{Appendices}
\chapter{prelude.py, Complete}
The whole file, exactly as it ships -- 15 sections, define-before-use,
no line over 105 columns.
\begin{lstlisting}[style=file]
"""prelude.py -- Haskell-style Prelude in pure Python, in pedagogical order.

Pascal discipline: strict define-before-use. Read top to bottom and every
name is already defined when you meet it. Sections build on each other:

   1 combinators          2 pairs                3 arithmetic
   4 list basics          5 higher-order lists   6 folds
   7 scans                8 streams (lazy)       9 slicing & spans
  10 sorting & searching 11 strings             12 containers
  13 nodes               14 grids & windows     15 control

Import qualified (import prelude as P): Haskell names shadow builtins by design.
Beyond the Prelude proper: the working parts of Data.List, Data.Map, Data.Maybe,
Data.Function and friends are folded into the same sections -- grouped by what
they do, not by the module they came from. None is Nothing throughout.
foldr = foldl(flip(f), reversed(xs), init): iterative, stack-safe, valid only
because Python is strict.
"""

# ============ 1. combinators ============ (functions about functions)

id        = lambda x: x                             # shadows builtin id() by design
const     = lambda x: lambda _: x
flip      = lambda f: (lambda x, y: f(y, x))        # flip f x y = f y x
curry     = lambda f: lambda x: lambda y: f(x, y)
uncurry   = lambda f: lambda p: f(*p)
partial   = lambda f, *bound: lambda *args: f(*bound, *args)
on        = lambda f, g: lambda x, y: f(g(x), g(y))  # Data.Function: combine via a key -- cmp `on` fst
NOTHING   = object()                                 # Maybe's Nothing: "no arg given"; test with `is`
fromMaybe = lambda d, x: d if x is None else x       # Data.Maybe: default for every None-returning tool
isJust    = lambda x: x is not None                  # Data.Maybe: None test as a handable predicate
isNothing = lambda x: x is None

# ============ 2. pairs ============

fst  = lambda p: p[0]
snd  = lambda p: p[1]
swap = lambda p: (p[1], p[0])               # Data.Tuple

# ============ 3. arithmetic & logic ============

succ      = lambda x: chr(ord(x) + 1) if isinstance(x, str) else x + 1   # Enum: numbers and Chars
pred      = lambda x: chr(ord(x) - 1) if isinstance(x, str) else x - 1
even      = lambda n: n % 2 == 0
odd       = lambda n: n % 2 == 1
not_      = lambda b: not b                                              # `not` is a Python keyword
otherwise = True
signum    = lambda x: (x > 0) - (x < 0)
div       = lambda a, b: a // b                                          # floors, like Haskell div
mod       = lambda a, b: a % b
quot      = lambda a, b: -(-a // b) if (a < 0) != (b < 0) else a // b    # truncates, like Haskell quot
rem       = lambda a, b: a - b * quot(a, b)
quotRem   = lambda a, b: (quot(a, b), rem(a, b))
gcd       = lambda a, b: abs(a) if b == 0 else gcd(b, a % b)             # >= 0 like Haskell; depth <= 90
lcm       = lambda a, b: abs(a // gcd(a, b) * b) if a and b else 0       # >= 0 like Haskell
hypot     = lambda x, y: (x*x + y*y) ** 0.5

def isqrt(n):  # floor sqrt without floats (Newton)
    if n < 0:
        raise ValueError("isqrt of negative")
    x = n
    y = (x + 1) // 2
    while y < x:
        x, y = y, (y + n // y) // 2
    return x if n else 0

# ============ 4. list basics ============

head        = lambda xs: xs[0]
tail        = lambda xs: xs[1:]
init        = lambda xs: xs[:-1]
last        = lambda xs: xs[-1]
null        = lambda xs: len(xs) == 0
elem        = lambda x, xs: x in xs
notElem     = lambda x, xs: x not in xs
replicate   = lambda n, x: [x] * n
drop        = lambda n, xs: list(xs)[max(n, 0):]                # finite lists only; n<0 drops nothing
splitAt     = lambda n, xs: (xs[:n], xs[n:]) if n >= 0 else (xs[:0], xs)   # n<0 splits at the front
nub         = lambda xs: list(dict.fromkeys(xs))    # unique, first occurrence wins (dicts keep order)
lookup      = lambda k, pairs: next((v for kk, v in pairs if kk == k), None)   # Nothing -> None
zip_        = lambda a, b: list(zip(a, b))          # shortest wins, any iterable
zip3        = lambda a, b, c: list(zip(a, b, c))
unzip       = lambda ps: tuple(map(list, zip(*ps))) if ps else ([], [])
enum        = lambda xs, start=0: zip_(list(range(start, start + len(xs))), xs)
pairwise    = lambda xs: list(zip(xs, xs[1:]))              # zip xs (tail xs)
isPrefixOf  = lambda p, xs: list(xs[:len(p)]) == list(p)             # Data.List; strings and lists alike
isSuffixOf  = lambda p, xs: list(xs[len(xs) - len(p):]) == list(p)   # NB not [-len(p):] -- [-0:] is ALL
stripPrefix = lambda p, xs: xs[len(p):] if isPrefixOf(p, xs) else None  # Just the rest, or Nothing
def transpose(xss):                    # Data.List: rows <-> cols, RAGGED-SAFE like Haskell
    xss = [list(xs) for xs in xss]     # any iterables in (rows may be one-shot); materialise once
    n = max(map(len, xss), default=0)  # short rows just drop out of later columns (no truncation)
    return [[xs[i] for xs in xss if i < len(xs)] for i in range(n)]

# ============ 5. higher-order lists ============

# Data.Maybe mapMaybe: map and drop the Nothings in ONE pass -- the parse-and-filter shape
mapMaybe  = lambda f, xs: [y for y in (f(x) for x in xs) if y is not None]
map_      = lambda f, xs: [f(x) for x in xs]
filter_   = lambda pred, xs: [x for x in xs if pred(x)]
# Data.List find: lazy first-match -- folds can't stop early, find can
find      = lambda pred, xs: next((x for x in xs if pred(x)), None)   # first match, else None
catMaybes = lambda xs: [x for x in xs if x is not None]                      # mapMaybe id
concat    = lambda xss: [x for xs in xss for x in xs]
concatMap = lambda f, xs: [y for x in xs for y in f(x)]
starmap   = lambda f, pairs: [f(*p) for p in pairs]   # map . uncurry
zipWith   = lambda f, a, b: [f(x, y) for x, y in zip(a, b)]
zipWith3  = lambda f, a, b, c: [f(x, y, z) for x, y, z in zip(a, b, c)]
cross     = lambda a, b: [(x, y) for x in a for y in b]     # cartesian (`product` names the fold)
def partition(pred, xs):                   # (keepers, rest) in ONE pass; pred called once per element
    yes, no = [], []
    for x in xs:
        (yes if pred(x) else no).append(x)
    return (yes, no)

# ============ 6. folds ============ (collapse a list to one value)

def compose(*fns): # foldr (.) id -- RIGHTMOST runs first; constant stack depth
    def composed(x):
        for f in reversed(fns):
            x = f(x)
        return x
    return composed

def foldl(f, xs, init=NOTHING):  # THE left fold; Python buried its own in functools as reduce
    it = iter(xs)
    if init is NOTHING:
        try:
            acc = next(it)
        except StopIteration:
            raise TypeError("fold of empty sequence with no initial value")
    else:
        acc = init
    for x in it:
        acc = f(acc, x)
    return acc

foldl1  = lambda f, xs: foldl(f, xs)
foldr   = lambda f, xs, init: foldl(flip(f), reversed(list(xs)), init)
foldr1  = lambda f, xs: foldl(flip(f), reversed(list(xs)))
product = lambda xs: foldl(lambda a, x: a * x, xs, 1)   # the numeric fold

# enumeration folds -- brute-force licenses for small n (say the bound out loud):
# Data.List subsequences: the powerset; each element DOUBLES acc -- 2^n, fine to n ~ 20
subsequences = lambda xs: foldl(lambda acc, x: acc + [s + [x] for s in acc], list(xs), [[]])
# Control.Monad replicateM: every length-n word over xs -- cross, n times; |xs|^n
replicateM   = lambda n, xs: foldl(lambda acc, _: [s + [x] for s in acc for x in xs], range(n), [[]])

# ============ 7. scans ============ (a fold that keeps its history)

def accumulate(xs, f=None, initial=NOTHING):     # scanl / scanl1
    if f is None:
        f = lambda a, b: a + b
    it = iter(xs)
    if initial is NOTHING:
        try:
            acc = next(it)
        except StopIteration:
            return
    else:
        acc = initial
    yield acc
    for x in it:
        acc = f(acc, x)
        yield acc

scanl  = lambda f, z, xs: list(accumulate(xs, f, initial=z))
scanl1 = lambda f, xs: list(accumulate(xs, f))
scanr  = lambda f, z, xs: list(reversed(scanl(flip(f), z, list(reversed(list(xs))))))
scanr1 = lambda f, xs: list(reversed(scanl1(flip(f), list(reversed(list(xs))))))

# ============ 8. streams ============ (lazy, possibly infinite)

def count(start=0, step=1):                 # [start, start+step ..]
    n = start
    while True:
        yield n
        n += step

def repeat(x, n=None):
    if n is None:
        while True:
            yield x
    else:
        for _ in range(n):
            yield x

def cycle(xs):
    xs = list(xs)
    if not xs:
        raise ValueError("cycle: empty list")   # Haskell errors too; the alternative is a silent hang
    while True:
        yield from xs

def iterate(f, x):                          # iterate f x = [x, f x, f (f x), ..]
    while True:
        yield x
        x = f(x)

def unfoldr(f, seed):                       # Data.List unfoldr: f(seed) -> None | (value, seed')
    out = []                                # iterate's finite twin: grows a list until f says stop --
    step = f(seed)                          # no more threading (state, acc) tuples through until
    while step is not None:
        val, seed = step
        out.append(val)
        step = f(seed)
    return out

def chain(*iterables):                      # concat
    for it in iterables:
        yield from it

# ============ 9. slicing & spans ============ (finite views of streams)

def islice(iterable, stop):                 # take
    it = iter(iterable)
    for _ in range(stop):
        try:
            yield next(it)
        except StopIteration:
            return

take      = lambda n, xs: list(islice(iter(xs), n))           # works on infinite streams

def takewhile(pred, xs):
    for x in xs:
        if not pred(x):
            return
        yield x

def takeuntil(pred, xs):                    # like takewhile(not . pred), but INCLUDES the first hit
    for x in xs:
        yield x
        if pred(x):
            return

def dropwhile(pred, xs):
    it = iter(xs)
    for x in it:
        if not pred(x):
            yield x
            break
    yield from it

takeWhile = lambda p, xs: list(takewhile(p, xs))
dropWhile = lambda p, xs: list(dropwhile(p, xs))
def span(p, xs):                            # split at first failure, ONE pass; safe on one-shot iters
    xs = list(xs)
    i = next((i for i, x in enumerate(xs) if not p(x)), len(xs))
    return (xs[:i], xs[i:])
break_    = lambda p, xs: span(lambda x: not p(x), xs)
inits     = lambda xs: [xs[:i] for i in range(len(xs) + 1)]   # Data.List: every prefix, [] first
tails     = lambda xs: [xs[i:] for i in range(len(xs) + 1)]   # every suffix; substrings start here
chunksOf  = lambda n, xs: [xs[i:i + n] for i in range(0, len(xs), n)]   # Data.List.Split; short last ok

# ============ 10. sorting & searching ============

sortOn = lambda f, xs: sorted(xs, key=f)    # sortOn (schwartzian, f called once per element)
minOn  = lambda f, xs: min(xs, key=f)    # best by key in O(n) -- sortOn + head without the sort
maxOn  = lambda f, xs: max(xs, key=f)    # (minimumBy/maximumBy + comparing; first wins ties)

def cmp_to_key(cmp):                        # resurrect Python 2's cmp for sortBy
    class K:
        def __init__(self, x): self.x = x
        def __lt__(self, other): return cmp(self.x, other.x) < 0
    return K

sortBy = lambda cmp, xs: sorted(xs, key=cmp_to_key(cmp))   # comparator sort, Python 2 style

def merge(a, b):                            # merge two sorted lists, stable; the heart of merge sort
    out, i, j = [], 0, 0                    # (section 13's merge_lists is this same loop on ListNodes)
    while i < len(a) and j < len(b):
        if a[i] <= b[j]:
            out.append(a[i]); i += 1
        else:
            out.append(b[j]); j += 1
    return out + a[i:] + b[j:]

def bisect_left(a, x):                      # first index where a[i] >= x
    lo, hi = 0, len(a)
    while lo < hi:
        mid = (lo + hi) // 2
        if a[mid] < x:
            lo = mid + 1
        else:
            hi = mid
    return lo

def bisect_right(a, x):                     # first index where a[i] > x
    lo, hi = 0, len(a)
    while lo < hi:
        mid = (lo + hi) // 2
        if a[mid] <= x:
            lo = mid + 1
        else:
            hi = mid
    return lo

# ============ 11. strings ============

words   = str.split
unwords = " ".join
lines   = str.splitlines
unlines = lambda ls: "".join(l + "\n" for l in ls)

# ============ 12. containers ============

class defaultdict(dict):
    def __init__(self, factory=None, *args, **kw):
        super().__init__(*args, **kw)
        self.factory = factory
    def __missing__(self, key):
        if self.factory is None:
            raise KeyError(key)
        self[key] = self.factory()
        return self[key]

def Counter(xs):                            # count-by-value; a defaultdict(int) fold
    d = defaultdict(int)
    for x in xs:
        d[x] += 1
    return d

def insertWith(f, k, v, d):                 # Data.Map insertWith f k v m: PERSISTENT -- returns a new
    out = dict(d)                           # dict untouched; collision -> f(new, old), Haskell order
    out[k] = f(v, out[k]) if k in out else v   # order (new value first)
    return out

def fromListWith(f, pairs):                 # Data.Map fromListWith: THE dict-building fold; f(new, old)
    d = {}                                  # like insertWith -- so grouping with concat REVERSES each
    for k, v in pairs:                      # group (the classic Haskell gotcha); use add for counts
        d[k] = f(v, d[k]) if k in d else v
    return d                                # Counter == fromListWith(add) over (x, 1);
                                            # grouping == fromListWith(++) over (k, [v])

def unionWith(f, a, b):                     # Data.Map unionWith: merge two dicts, f on shared keys
    out = dict(a)                           # f(a's value, b's value) -- Haskell's left/right order
    for k, v in b.items():
        out[k] = f(out[k], v) if k in out else v
    return out                              # one combinator, many merges:
                                            # bag_union == unionWith(max), merge-sum == unionWith(add)

def groupBy(eq, xs):                        # Data.List groupBy: runs of CONSECUTIVE elements; each new
    out = []                                # element is compared via eq against the run's FIRST element,
    for x in xs:                            # exactly like Haskell (span-based) -- not its neighbor
        if out and eq(out[-1][0], x):
            out[-1].append(x)
        else:
            out.append([x])
    return out

group = lambda xs: groupBy(lambda a, b: a == b, xs)   # Data.List group: runs of equals, [[a]] -- no keys

# multiset ops on Counters (use .get, not [] -- indexing a Counter autovivifies zeros):
bag_union = lambda a, b: {k: max(a.get(k, 0), b.get(k, 0)) for k in a.keys() | b.keys()}
bag_inter = lambda a, b: {k: v for k in a.keys() & b.keys() if (v := min(a.get(k, 0), b.get(k, 0)))}
bag_diff  = lambda a, b: {k: a[k] - b[k] for k in a if a[k] > b.get(k, 0)}    # clipped at 0
bag_sub   = lambda a, b: all(b.get(k, 0) >= n for k, n in a.items())          # a <= b (inclusion)

Tree  = lambda depth, leaf: (defaultdict(leaf) if depth == 1
                             else defaultdict(lambda: Tree(depth - 1, leaf)))
ITree = lambda: defaultdict(ITree)

def paths(t):                               # cata for Tree/ITree: [(keypath, leaf)]
    # recursion depth = tree height (word length / len(nums)) -- well under the ~1000 limit
    if not isinstance(t, dict) or not t:    # leaf = non-dict value, or empty node
        return [([], t)]
    return [([k] + p, leaf) for k, sub in t.items() for p, leaf in paths(sub)]

leaves = lambda t: [l for _, l in paths(t)]

def setpath(t, ks, v):                      # write leaf v at key path ks (autovivifies interior)
    for k in ks[:-1]:
        t = t[k]
    t[ks[-1]] = v                           # NB: clobbers any subtree already at ks

def getpath(t, ks, default=None):           # read along ks WITHOUT autovivifying
    for k in ks:
        if not isinstance(t, dict) or k not in t:
            return default
        t = t[k]
    return t

class deque:                                # two-stack; all four ends amortized O(1)
    def __init__(self, xs=()):
        self._in, self._out = [], list(reversed(list(xs)))
    def append(self, x): self._in.append(x)
    def appendleft(self, x): self._out.append(x)
    def pop(self):
        if not self._in:
            self._in, self._out = list(reversed(self._out)), []
        return self._in.pop()
    def popleft(self):
        if not self._out:
            self._out, self._in = list(reversed(self._in)), []
        return self._out.pop()
    def __len__(self): return len(self._in) + len(self._out)
    def __iter__(self): return iter(self._out[::-1] + self._in)

def dsu(n):                                 # union-find, path halving; union -> False if already joined
    parent = list(range(n))
    def find(x):
        while parent[x] != x:
            parent[x] = parent[parent[x]]
            x = parent[x]
        return x
    def union(a, b):
        ra, rb = find(a), find(b)
        if ra == rb:
            return False
        parent[ra] = rb
        return True
    return find, union

def heappush(h, x):                         # min-heap on a plain list: sift-up
    h.append(x)
    i = len(h) - 1
    while i and h[(i - 1) // 2] > h[i]:
        h[(i - 1) // 2], h[i] = h[i], h[(i - 1) // 2]
        i = (i - 1) // 2

def heappop(h):                             # min-heap on a plain list: sift-down
    h[0], h[-1] = h[-1], h[0]
    top = h.pop()
    i, n = 0, len(h)
    while True:
        s, l, r = i, 2*i + 1, 2*i + 2
        if l < n and h[l] < h[s]: s = l
        if r < n and h[r] < h[s]: s = r
        if s == i:
            return top
        h[i], h[s] = h[s], h[i]
        i = s

# ============ 13. nodes ============ (the LeetCode givens: binary trees & linked lists)

class TreeNode:                             # the LeetCode binary tree
    def __init__(self, val=0, left=None, right=None):
        self.val, self.left, self.right = val, left, right

def tfold(f, z, t):                         # cata: f(val, l_acc, r_acc); depth = tree height
    return z if t is None else f(t.val, tfold(f, z, t.left), tfold(f, z, t.right))

inorder   = lambda t: tfold(lambda v, l, r: l + [v] + r, [], t)
preorder  = lambda t: tfold(lambda v, l, r: [v] + l + r, [], t)
postorder = lambda t: tfold(lambda v, l, r: l + r + [v], [], t)
tdepth    = lambda t: tfold(lambda _, l, r: 1 + max(l, r), 0, t)
tsize     = lambda t: tfold(lambda _, l, r: 1 + l + r, 0, t)

def levelorder(t):                          # BFS: the deque earning its keep
    out, q = [], deque([t] if t else [])
    while len(q):
        n = q.popleft()
        out.append(n.val)
        if n.left:  q.append(n.left)
        if n.right: q.append(n.right)
    return out

class ListNode:                             # the LeetCode linked list
    def __init__(self, val=0, next=None):
        self.val, self.next = val, next

def from_list(xs):                          # build; foldr of ListNode
    node = None
    for x in reversed(list(xs)):
        node = ListNode(x, node)
    return node

def to_list(node):                          # unfold back to a Python list
    out = []
    while node:
        out.append(node.val)
        node = node.next
    return out

def reverse_list(node):                     # three-pointer; prev is a fold accumulator
    prev = None
    while node:
        node.next, prev, node = prev, node, node.next
    return prev

def middle(node):                           # fast/slow: second middle for even lengths
    slow = fast = node
    while fast and fast.next:
        slow, fast = slow.next, fast.next.next
    return slow

def has_cycle(node):                        # Floyd: fast laps slow iff a cycle exists
    slow = fast = node
    while fast and fast.next:
        slow, fast = slow.next, fast.next.next
        if slow is fast:
            return True
    return False

def merge_lists(a, b):                      # merge two sorted lists; dummy-head idiom
    dummy = tail_ = ListNode()
    while a and b:
        if a.val <= b.val:
            tail_.next, a = a, a.next
        else:
            tail_.next, b = b, b.next
        tail_ = tail_.next
    tail_.next = a or b
    return dummy.next

# ============ 14. grids & windows ============

def neighbors4(r, c, R, C):                 # in-bounds orthogonal neighbors
    for dr, dc in ((1, 0), (-1, 0), (0, 1), (0, -1)):
        if 0 <= r + dr < R and 0 <= c + dc < C:
            yield r + dr, c + dc

def neighbors8(r, c, R, C):                 # ... plus diagonals
    for dr in (-1, 0, 1):
        for dc in (-1, 0, 1):
            if (dr or dc) and 0 <= r + dr < R and 0 <= c + dc < C:
                yield r + dr, c + dc

def longest_window(xs, valid, add, rem):    # sliding-window skeleton; state lives in the closures
    lo = best = 0
    for hi in range(len(xs)):
        add(xs[hi])
        while not valid():                  # shrink until legal again
            rem(xs[lo])
            lo += 1
        best = max(best, hi - lo + 1)
    return best

# ============ 15. control ============

def memo(f):                                # unbounded memoizer; enough for DP (no eviction by design)
    cache = {}
    def wrapped(*args, **kw):
        key = (args, frozenset(kw.items())) if kw else args
        if key not in cache:
            cache[key] = f(*args, **kw)
        return cache[key]
    wrapped.cache = cache                   # peek at the DP table if curious
    return wrapped

def until(p, f, x):                         # until p f x -- loop, not recursion (unbounded depth)
    while not p(x):
        x = f(x)
    return x
\end{lstlisting}
\chapter{Haskell Provenance}
Where each borrowed name comes from, for the day you meet the originals.

\begin{center}\small
\begin{tabular}{@{}p{0.25\textwidth}p{0.68\textwidth}@{}}
\toprule
\textbf{Module} & \textbf{Names} \\
\midrule
Prelude &
id, const, flip, curry, uncurry, fst, snd, succ, pred, even, odd, signum,
div, mod, quot, rem, gcd, lcm, head, tail, init, last, null, elem,
replicate, drop, splitAt, lookup, zip, unzip, map, filter,
concat, concatMap, zipWith, foldl, foldr, product, iterate, repeat,
cycle, take, takeWhile, dropWhile, span, break, words, unwords, lines,
unlines, until \\
Data.List &
find, partition, nub, sortOn, sortBy, transpose, unfoldr, subsequences,
inits, tails, isPrefixOf, isSuffixOf, stripPrefix, minimumBy/maximumBy
(as minOn/maxOn), groupBy, group \\
Data.Map &
insertWith, fromListWith, unionWith \\
Data.Maybe &
fromMaybe, isJust, isNothing, mapMaybe, catMaybes (None is Nothing) \\
Data.Function & on \\
Data.Tuple & swap \\
Control.Monad & replicateM \\
Data.List.Split & chunksOf \\
Python itertools (names) &
count, islice, takewhile, dropwhile, chain, accumulate, starmap, pairwise --
kept under Python's own names so both vocabularies stay warm \\
\bottomrule
\end{tabular}
\end{center}

Deliberately not ported: laziness-dependent knot-tying (\texttt{fix}),
typeclass machinery, and monadic control flow -- Python is strict and
untyped, and the prelude stays honest about it.
\chapter{How to Practice}
The course directory is the textbook's other half:

\begin{itemize}
\item \texttt{fpython/prelude.py} -- the toolkit. Study loop: read a section,
close the file, retype its definitions from memory, diff.
\item \texttt{fpython/g1\_\ldots{} g9\_\ldots} -- 74 practice files. Each is
self-contained: statement, contract, hint, the needed snippets copied inline,
a \texttt{\# solution goes here} slot, and a doctest runner. Run with
\texttt{python3 file.py} until it reports $N/N$.
\item \texttt{fpython/solutions/} -- the same files solved, flat, named
\texttt{g2\_08\_chunks.py}. For \emph{after} the fight: compare shapes, not
before.
\item \texttt{fpython/prelude\_flowchart.svg} -- the recognition ladder of
Chapter~2 as a printable picture.
\end{itemize}

Work the classes in order; within a class, any order. Narrate out loud as you
solve -- shape, tools, complexity -- because the interview grades the
narration as much as the code. When a solution of yours beats the book's,
trust the doctests and keep yours.
\backmatter
\end{document}
